\chapter{Biochemistry}

\medterm{Arginase Deficiency}-- An autosomal-recessive urea-cycle disorder caused by deficiency of arginase 1, resulting in elevated arginine and impaired nitrogen disposal. Unlike many urea-cycle disorders, it typically causes progressive spasticity, developmental delay, and growth impairment rather than recurrent severe neonatal hyperammonemia. Management includes protein and arginine reduction, specialized nutrition, ammonia-scavenging therapy, and metabolic specialist care.

\medterm{Argininemia}-- The biochemical and clinical disorder caused by arginase deficiency, characterized by persistent hyperargininemia and variable hyperammonemia. It is usually recognized in childhood through progressive motor dysfunction, spasticity, seizures, developmental delay, or growth failure; diagnosis relies on plasma amino acids and molecular or enzyme testing.

\medterm{Acetyl-CoA Carboxylase}-- The rate-limiting and committed enzyme of fatty acid synthesis, catalyzing the ATP-dependent carboxylation of acetyl-CoA to malonyl-CoA in the cytoplasm. It functions as an active polymeric filament requiring biotin (vitamin B7), ATP, and bicarbonate ($HCO_3^-$). Allosterically activated by citrate (which signals an excess of mitochondrial energy and carbon substrates, promoting polymerization) and allosterically inhibited by long-chain fatty acyl-CoAs (such as palmitoyl-CoA, mediating feedback inhibition). Hormonally, insulin stimulates protein phosphatase 2A to dephosphorylate and activate the enzyme, favoring lipogenesis; in contrast, glucagon, epinephrine, and AMP-activated protein kinase (AMPK) phosphorylate and inactivate the enzyme during energy deficit. The reaction product, malonyl-CoA, allosterically inhibits carnitine palmitoyltransferase-1 (CPT-1), effectively preventing newly synthesized fatty acids from entering the mitochondria for beta-oxidation.

\medterm{Acetyl-Coenzyme A}-- A central two-carbon metabolic crossroad intermediate consisting of an acetyl group linked to the reactive terminal sulfhydryl (-SH) group of coenzyme A (derived from pantothenic acid, vitamin B5) via a high-energy thioester bond ($\Delta G^{\circ\prime}$ of hydrolysis $\approx -31.5$ kJ/mol). Acetyl-CoA is generated within the mitochondrial matrix through three major convergent catabolic pathways: (1) The oxidative decarboxylation of pyruvate by the pyruvate dehydrogenase complex (linking aerobic glycolysis to the citric acid cycle); (2) The cyclic $\beta$-oxidation of fatty acids; and (3) The oxidative catabolism of ketogenic amino acids (leucine, isoleucine, lysine, phenylalanine, tyrosine, tryptophan). Mitochondrial acetyl-CoA condenses with oxaloacetate via citrate synthase to fuel the citric acid cycle for ATP generation, or, during states of fasting or untreated diabetes, is funneled into ketogenesis (forming acetoacetate and $\beta$-hydroxybutyrate). Translocated to the cytoplasm via the citrate shuttle, acetyl-CoA serves as the essential building block for de novo fatty acid synthesis (via acetyl-CoA carboxylase) and cholesterol biosynthesis (via HMG-CoA reductase). In the brain, acetyl-CoA condenses with choline via choline acetyltransferase to synthesize the neurotransmitter acetylcholine, while in the nucleus, it acts as the obligatory acetyl donor for histone acetyltransferases (HATs) to regulate epigenetic chromatin remodeling.

\medterm{Acid}-- A substance that donates protons (H+) or accepts electron pairs (Lewis acid). In aqueous solution, acids increase H+ concentration (pH <7). Strong acids (HCl, H2SO4, HNO3) dissociate completely; weak acids (acetic acid, carbonic acid H2CO3, phosphate, amino acid side chains) partially dissociate (Ka, pKa). In biochemistry, key acids include: carbonic acid (H2CO3, from CO2 + H2O, major blood buffer with bicarbonate), lactic acid (anaerobic glycolysis), ketone bodies (acetoacetate, beta-hydroxybutyrate, diabetic ketoacidosis), uric acid (purine catabolism, gout), fatty acids, amino acids (carboxyl group, pKa ~2), and nucleotide phosphates. Acid-base balance is maintained by respiratory (CO2) and renal (H+, HCO3-) compensation. Anion gap = [Na+] - ([Cl-] + [HCO3-]) identifies unmeasured anions.

\medterm{Activated Partial Thromboplastin Time}-- A clinical screening coagulation assay that evaluates the functional integrity of the intrinsic and common pathways of the coagulation cascade (specifically factors XII, XI, IX, VIII, X, V, II [prothrombin], and I [fibrinogen]). The test is performed by incubating citrated, platelet-poor plasma with a contact activator (such as micronized silica, celite, or kaolin) and partial thromboplastin (phospholipid substitute for platelets) at 37$^\circ$C, followed by recalcification with $CaCl_2$, measuring the time required for fibrin clot formation (reference range: 25--35 seconds). The aPTT is the standard laboratory modality used to monitor unfractionated heparin therapy, with a therapeutic goal of 1.5 to 2.5 times the patient's baseline value (typically 60--80 seconds). An isolated prolonged aPTT occurs in inherited factor deficiencies (hemophilia A [factor VIII], hemophilia B [factor IX], factor XI deficiency), severe von Willebrand disease (due to loss of factor VIII carrier stabilization), factor XII deficiency (asymptomatic in vivo), and in the presence of lupus anticoagulants (antiphospholipid antibodies that cause an in vitro artifact of prolonged clotting time despite promoting in vivo thrombosis).

\medterm{Active Site}-- The specific three-dimensional region of an enzyme where substrate binding and catalysis occur. Formed by amino acid residues that may be distant in primary sequence but brought together by protein folding. The active site provides a unique microenvironment (specific pH, hydrophobicity, electrostatics) that stabilizes the transition state, lowering activation energy. Substrate binding follows induced fit (conformational change) or lock-and-key models. Active site residues directly participate in catalysis (acid-base, covalent, metal ion catalysis).

\medterm{Adiponectin}-- A 30-kDa peptide hormone secreted almost exclusively by mature adipocytes that plays an essential anti-atherogenic, anti-inflammatory, and insulin-sensitizing role. Unlike most adipokines (such as leptin and TNF-$\alpha$), circulating adiponectin levels are inversely correlated with body mass index and visceral adiposity. Adiponectin binds to AdipoR1 (predominant in skeletal muscle) and AdipoR2 (predominant in liver) receptors, stimulating AMP-activated protein kinase (AMPK) and peroxisome proliferator-activated receptor alpha (PPAR-$\alpha$). This downstream signaling enhances glucose uptake, promotes fatty acid beta-oxidation, decreases hepatic gluconeogenesis, and inhibits endothelial vascular cell adhesion molecule-1 (VCAM-1) expression. Hypoadiponectinemia is strongly linked to insulin resistance, metabolic syndrome, type 2 diabetes mellitus, and coronary artery disease.

\medterm{ADP (Adenosine Diphosphate)}-- A nucleotide composed of adenine, ribose, and two phosphate groups. Formed by hydrolysis of ATP (ATP $\to$ ADP + P$_i$, $\Delta G^\circ$ = -30.5 kJ/mol). Serves as the immediate precursor for ATP synthesis via oxidative phosphorylation (mitochondrial ATP synthase), substrate-level phosphorylation (glycolysis, TCA cycle), and photophosphorylation. The ADP/ATP ratio regulates cellular energy state: high ADP stimulates catabolic pathways (glycolysis, beta-oxidation, TCA cycle) and inhibits anabolic pathways.

\medterm{Albumin}-- The most abundant plasma protein (35-50 g/L, ~60% of total protein), synthesized in the liver. A single polypeptide (66.5 kDa, 585 amino acids) with high negative charge at physiological pH. Functions: (1) Maintains colloid osmotic pressure (oncotic pressure, ~75% of total), retaining fluid in vasculature; (2) Transport: binds and transports fatty acids, bilirubin, hormones (thyroxine, cortisol), drugs (warfarin, penicillins), calcium, and metal ions; (3) Antioxidant: free thiol (Cys34) scavenges reactive oxygen species; (4) Buffering: contributes to plasma buffering capacity. Low albumin (hypoalbuminemia) causes edema, reduced drug binding, and is a marker of malnutrition, liver disease, nephrotic syndrome, inflammation, or protein-losing enteropathy. Albumin infusion used for volume expansion in specific settings.

\medterm{Alkaptonuria}-- An autosomal recessive inborn error of tyrosine catabolism caused by deficiency of homogentisate 1,2-dioxygenase (homogentisic acid oxidase), encoded by the HGD gene on chromosome 3. In the normal tyrosine degradation pathway, 4-hydroxyphenylpyruvate is converted to homogentisic acid (HGA); homogentisate oxidase oxidizes HGA to 4-maleylacetoacetate. Enzyme deficiency prevents HGA breakdown, causing massive accumulation of homogentisic acid in tissues and its excretion in large quantities in urine. Clinical hallmarks include: (1) Urine that turns dark brownish-black upon standing, alkalinization, or exposure to air, caused by spontaneous oxidation of HGA into a melanin-like polyphenol pigment; (2) Ochronosis: progressive deposition of dark blue-black pigment in connective tissues, sclerae, and cartilage (notably ears and nose), typically becoming visible in the third decade of life; and (3) Ochronotic arthropathy: severe, debilitating degenerative arthritis affecting the spine (with characteristic radiographic intervertebral disc calcification and narrowing) and large weight-bearing joints (hips, knees, shoulders). Pigment accumulation also causes cardiac valvular calcification (aortic stenosis) and renal calculi. Treatment includes dietary restriction of phenylalanine and tyrosine, and nitisinone (an inhibitor of 4-hydroxyphenylpyruvate dioxygenase that blocks the synthesis of homogentisic acid).

\medterm{Allosteric Regulation}-- The regulation of an enzyme or protein activity by the binding of an effector molecule (activator or inhibitor) at a site other than the active site (the allosteric site). This binding induces a conformational change that alters the protein's affinity for its substrate. Allosteric enzymes typically exhibit sigmoidal kinetics rather than hyperbolic Michaelis-Menten kinetics. A classic example is the feedback inhibition of phosphofructokinase-1 by ATP and citrate, or its activation by AMP and fructose-2,6-bisphosphate.

\medterm{Alpha-Ketoglutarate}-- A foundational five-carbon dicarboxylic $\alpha$-ketoacid intermediate of the citric acid cycle (Krebs cycle) that serves as a central metabolic junction connecting carbohydrate energy metabolism with amino acid transamination and nitrogen disposal. In the mitochondrial matrix, $\alpha$-ketoglutarate is generated from the oxidative decarboxylation of D-isocitrate by isocitrate dehydrogenase (generating NADH and $CO_2$); it is subsequently converted to succinyl-CoA by the multienzyme $\alpha$-ketoglutarate dehydrogenase complex (generating NADH, $CO_2$, and requiring TPP, lipoate, CoA, FAD, and $NAD^+$). In amino acid metabolism, $\alpha$-ketoglutarate acts as the universal amino group acceptor for virtually all transaminases (ALT, AST), accepting an $\alpha$-amino group from donor amino acids to form L-glutamate while liberating the corresponding $\alpha$-ketoacid. Mitochondrial glutamate dehydrogenase then reversibly deaminates glutamate back to $\alpha$-ketoglutarate, releasing toxic free ammonia ($NH_4^+$) for detoxification in the urea cycle. Beyond metabolism, $\alpha$-ketoglutarate functions as an essential co-substrate for Fe(II)/$\alpha$-ketoglutarate-dependent dioxygenases, including prolyl hydroxylases (which target HIF-1$\alpha$ for proteasomal destruction under normoxia), ten-eleven translocation (TET) DNA demethylases, and Jumonji histone demethylases; accumulation of structural oncometabolites (such as 2-hydroxyglutarate in IDH-mutant tumors or succinate in SDH-deficient paragangliomas) competitively inhibits these enzymes, driving pseudo-hypoxia and tumorigenesis.

\medterm{Amino Acid}-- The building blocks of proteins, 20 standard L-$\alpha$-amino acids sharing a central carbon ($\alpha$-carbon) bonded to an amino group (-NH2), carboxyl group (-COOH), hydrogen, and a unique side chain (R group). Classification by side chain properties: nonpolar (glycine, alanine, valine, leucine, isoleucine, proline, methionine, phenylalanine, tryptophan), polar uncharged (serine, threonine, cysteine, asparagine, glutamine), acidic (aspartate, glutamate), basic (lysine, arginine, histidine). Nine essential amino acids (cannot synthesize): histidine, isoleucine, leucine, lysine, methionine, phenylalanine, threonine, tryptophan, valine. Amino acids undergo transamination (PLP-dependent), deamination, and decarboxylation. Carbon skeletons feed into glycolysis/TCA (glucogenic: pyruvate, oxaloacetate, fumarate, succinyl-CoA, alpha-ketoglutarate; ketogenic: acetyl-CoA, acetoacetate; both: isoleucine, phenylalanine, tyrosine, tryptophan). Branched-chain amino acids (BCAA: leucine, isoleucine, valine) are primarily metabolized in muscle. Branched-chain ketoacid dehydrogenase deficiency causes maple syrup urine disease.

\medterm{AMP (Adenosine Monophosphate)}-- A nucleotide composed of adenine, ribose, and a single phosphate group. Formed by hydrolysis of ADP (ADP -> AMP + Pi) or directly from ATP via adenylate kinase (2 ADP <-> ATP + AMP). AMP is a key allosteric signal of low energy status: activates AMPK (AMP-activated protein kinase), phosphorylase kinase (glycogenolysis), and PFK-1 (glycolysis). Degraded to IMP by AMP deaminase in purine catabolism. The ATP:ADP:AMP ratio (energy charge) regulates metabolic flux.

\medterm{AMPK}-- AMP-activated protein kinase, a phylogenetically conserved heterotrimeric serine/threonine kinase that operates as the master energy sensor and metabolic rheostat of the eukaryotic cell. Composed of a catalytic $\alpha$-subunit and regulatory $\beta$- and $\gamma$-subunits, AMPK is allosterically activated when AMP or ADP binds to cystathionine $\beta$-synthase (CBS) domains on the $\gamma$-subunit, which also promotes phosphorylation at Thr172 by upstream kinases (notably LKB1 and CaMKK$\beta$) and inhibits dephosphorylation. Activated AMPK coordinates cellular survival during energy stress (hypoxia, glucose deprivation, exercise) by shutting down ATP-consuming anabolic pathways and switching on ATP-generating catabolic pathways. It phosphorylates and inactivates acetyl-CoA carboxylase (halting lipogenesis and disinhibiting CPT-1 for beta-oxidation), inhibits HMG-CoA reductase (suppressing cholesterol synthesis), inhibits glycogen synthase, and downregulates mTORC1 to halt translation, while stimulating GLUT4-mediated glucose uptake, glycolysis, mitochondrial biogenesis (via PGC-1$\alpha$), and autophagy. Pharmacologically, AMPK is activated by metformin and exercise.

\medterm{Amylase}-- A metalloenzyme requiring both calcium and chloride that catalyzes the endohydrolysis of $\alpha$-(1$\to$4)-glucosidic linkages in starch, glycogen, and related polysaccharides into maltose, maltotriose, and $\alpha$-limit dextrins. Synthesized and secreted predominantly by pancreatic acinar cells (P-type isoamylase, 40--45\% of serum activity) into pancreatic juice and by salivary glands (S-type isoamylase, 55--60\%) into saliva. Normal serum amylase ranges from 30 to 110 U/L. In acute pancreatitis, serum amylase rises rapidly within 3 to 6 hours of onset, peaks at 12 to 24 hours, and normalizes within 3 to 5 days. Because amylase is also elevated in salivary gland inflammation (parotitis, mumps), bowel obstruction or infarction, ruptured ectopic pregnancy, renal failure (decreased clearance), and macroamylasemia (complexing with immunoglobulins, preventing renal excretion), serum lipase is preferred for its greater diagnostic specificity for acute pancreatitis.

\medterm{Anabolism}-- The set of metabolic pathways that construct complex molecules from simpler precursors, requiring energy input (ATP, reducing equivalents NADPH). Examples: protein synthesis from amino acids, glycogen synthesis (glycogenesis), fatty acid synthesis (lipogenesis), cholesterol synthesis, nucleotide synthesis, gluconeogenesis. Anabolic pathways are typically stimulated by insulin and inhibited by glucagon, epinephrine, and AMPK. They share intermediates with catabolic pathways but use distinct regulated enzymes (reciprocal regulation).

\medterm{Anemia of Chronic Disease}-- Also known as anemia of inflammation, the second most prevalent form of anemia worldwide, developing in patients with chronic infections, autoimmune diseases (rheumatoid arthritis, SLE), chronic kidney disease, or malignancies. The central pathogenic driver is the pro-inflammatory cytokine interleukin-6 (IL-6), which transcriptionally upregulates hepatic production of the master iron-regulatory hormone hepcidin. Elevated hepcidin binds, internalizes, and degrades ferroportin-1, the sole cellular iron exporter on basolateral membranes of duodenal enterocytes and reticuloendothelial macrophages. Consequently, iron is locked within intracellular ferritin stores, starving developing erythroid precursors and causing functional iron deficiency despite abundant total body iron. Cytokines (TNF-$\alpha$, IFN-$\gamma$) also directly suppress erythropoietin production in the kidney and blunt bone marrow erythroid progenitor proliferation. Typical laboratory findings reveal a normocytic normochromic or mildly microcytic anemia, low serum iron, low total iron-binding capacity (TIBC), normal or elevated serum ferritin ($>100$ ng/mL), and elevated inflammatory markers (ESR, CRP).

\medterm{Anion Gap}-- A clinical biochemical calculation derived from standard serum electrolyte measurements that estimates the concentration of unmeasured anions in plasma, primarily calculated as: $\text{Anion Gap} = [Na^+] - ([Cl^-] + [HCO_3^-])$, with a modern normal reference range of 8 to 12 mEq/L (or 10 to 14 mEq/L if potassium is included: $[Na^+ + K^+] - [Cl^- + HCO_3^-]$). Because plasma maintains strict electrical neutrality (total cations equal total anions), the gap represents the surplus of unmeasured physiological plasma anions (predominantly negatively charged albumin [contributing ~11 mEq/L], along with phosphate, sulfate, and organic acids) over unmeasured cations ($Ca^{2+}, Mg^{2+}$, and $\gamma$-globulins). Because serum albumin accounts for the vast majority of normal unmeasured anions, the expected normal anion gap MUST be adjusted downward in hypoalbuminemia: corrected anion gap = calculated anion gap + $2.5 \times (4.0 - [\text{serum albumin in g/dL}])$. Calculation of the anion gap is the essential first step in evaluating metabolic acidosis, categorizing it into: (1) High Anion Gap Metabolic Acidosis (HAGMA, $>12$ mEq/L), caused by accumulation of nonvolatile unmeasured organic acids, summarized by the mnemonic MUDPILES (Methanol, Uremia, Diabetic ketoacidosis, Propylene glycol, Iron/Isoniazid, Lactic acidosis, Ethylene glycol, Salicylates); and (2) Normal Anion Gap (Hyperchloremic) Metabolic Acidosis (NAGMA), where bicarbonate loss is compensated by an equimolar retention of chloride (severe diarrhea, renal tubular acidosis, acetazolamide use).

\medterm{Anticodon}-- A three-nucleotide sequence at one end of transfer RNA that base-pairs with a complementary messenger RNA codon during translation, ensuring the correct amino acid is incorporated into the growing polypeptide. Wobble pairing at the first anticodon position permits one tRNA to recognize multiple synonymous codons, explaining codon degeneracy.

\medterm{Antioxidant}-- Any substance, nutrient, or endogenous enzymatic system that, when present at low concentrations compared to an oxidizable substrate, significantly delays, intercepts, or prevents the oxidation of that substrate by neutralizing reactive oxygen species (ROS) and reactive nitrogen species (RNS). The human body maintains a multilayered antioxidant defense network categorized into enzymatic and non-enzymatic components: (1) Primary Enzymatic Antioxidants: Superoxide dismutase (SOD) catalyzes the dismutation of superoxide radicals into oxygen and hydrogen peroxide ($2 O_2^{\bullet-} + 2 H^+ \to O_2 + H_2O_2$), operating via cytosolic Cu/Zn-SOD (SOD1) and mitochondrial Mn-SOD (SOD2); Catalase, a heme-containing tetrameric enzyme localized to peroxisomes, decomposes hydrogen peroxide into water and molecular oxygen ($2 H_2O_2 \to 2 H_2O + O_2$); and Glutathione peroxidase (GPx), a selenium-dependent cytosolic and mitochondrial enzyme that reduces hydrogen peroxide and organic hydroperoxides to water using reduced glutathione (GSH) as the electron donor ($H_2O_2 + 2 GSH \to GSSG + 2 H_2O$), with oxidized glutathione (GSSG) subsequently regenerated by NADPH-dependent glutathione reductase; (2) Non-Enzymatic Endogenous Scavengers: Reduced glutathione (GSH, a tripeptide $\gamma$-Glu-Cys-Gly providing a critical sulfhydryl buffer), ubiquinol (reduced CoQ10), bilirubin, and uric acid; and (3) Exogenous Dietary Antioxidants: $\alpha$-Tocopherol (vitamin E, the premier chain-breaking lipid-soluble antioxidant protecting cell membranes), ascorbic acid (vitamin C, water-soluble radical scavenger that recycles oxidized vitamin E), and carotenoids.

\medterm{Antiphospholipid Antibody}-- A family of pathogenic autoantibodies directed against plasma proteins bound to anionic phospholipids, primarily $\beta_2$-glycoprotein I ($\beta_2$GPI), cardiolipin, and prothrombin. These antibodies define antiphospholipid syndrome (APS), which can occur as a primary autoimmune entity or secondary to systemic lupus erythematosus (SLE). Clinically, laboratory detection requires at least one positive test on two occasions separated by at least 12 weeks: (1) Lupus anticoagulant (demonstrating phospholipid-dependent clot prolongation on dilute Russell viper venom time [dRVVT] or sensitive aPTT that fails to correct on a 1:1 mixing study but corrects upon the addition of excess phospholipids); (2) Anti-cardiolipin IgG/IgM antibodies; or (3) Anti-$\beta_2$-glycoprotein I IgG/IgM antibodies. Paradoxically, while these antibodies prolong phospholipid-dependent in vitro clotting assays by competing with coagulation factors for reagent lipids, they drive a severe hypercoagulable state in vivo: $\beta_2$GPI-antibody complexes cross-link endothelial cells, platelets, and monocytes, triggering tissue factor expression, platelet aggregation, and complement activation. Patients present with recurrent deep vein thrombosis, pulmonary embolism, arterial strokes, and obstetric complications (unexplained recurrent early pregnancy losses or severe preeclampsia). Management requires long-term therapeutic anticoagulation with vitamin K antagonists (warfarin), as direct oral anticoagulants (DOACs) are inferior in triple-positive APS.

\medterm{Antithrombin}-- A 58-kDa plasma glycoprotein belonging to the serpin (serine protease inhibitor) superfamily, synthesized by hepatocytes, that operates as the body's primary physiological brake on the coagulation cascade. Antithrombin inactivates thrombin (factor IIa) and activated factor X (FXa), as well as factors IXa, XIa, and XIIa, by presenting a reactive center loop (Arg393-Ser394) that forms an irreversible, equimolar suicide-substrate acyl-enzyme complex with the protease's catalytic serine residue. In the absence of catalysts, this inhibition proceeds at a slow baseline rate; however, binding of unfractionated heparin, low-molecular-weight heparin (LMWH), or vascular endothelial heparan sulfate proteoglycans to a specific pentasaccharide domain on antithrombin induces an allosteric conformational change that accelerates the rate of factor Xa and thrombin inhibition by $>1,000$-fold. Hereditary antithrombin deficiency (autosomal dominant) confers a 20- to 50-fold increased risk of venous thromboembolism, typically presenting before age 40 with deep vein thrombosis, pulmonary embolism, or mesenteric vein thrombosis. Patients exhibit marked resistance to unfractionated heparin, requiring antithrombin concentrate infusion or direct thrombin inhibitors (argatroban, bivalirudin).

\medterm{Apoenzyme}-- The protein component of an enzyme without its required cofactor or coenzyme. Inactive or minimally active until the cofactor binds. The complete, catalytically active complex of apoenzyme + cofactor is called the holoenzyme. Many enzymes require metal ions (Zn2+, Mg2+, Fe2+/Fe3+, Cu2+, Mn2+) or organic cofactors (NAD+, FAD, coenzyme A, TPP, PLP, biotin) for activity. Apoenzyme folding may depend on cofactor binding.

\medterm{ATP (Adenosine Triphosphate)}-- The primary energy currency of the cell, composed of adenine, ribose, and three phosphate groups. Hydrolysis of the terminal phosphoanhydride bond releases $\Delta G^\circ$ = -30.5 kJ/mol (-7.3 kcal/mol), driving endergonic cellular processes such as muscle contraction, active transport, and biosynthesis. Generated via substrate-level phosphorylation (glycolysis, TCA cycle) and oxidative phosphorylation in mitochondria.

\medterm{ATP Synthase}-- Also designated Complex V of the mitochondrial respiratory chain, a multisubunit rotary nanomotor embedded in the inner mitochondrial membrane that couples proton translocation down an electrochemical proton gradient ($\Delta p$) to the synthesis of ATP from ADP and inorganic phosphate ($P_i$). The enzyme comprises two functional assemblies: the membrane-spanning $F_o$ domain (subunits $a$, $b$, and a rotating oligomeric $c$-ring) that forms the proton channel, and the soluble matrix $F_1$ domain ($\alpha_3\beta_3\gamma\delta\epsilon$) that houses the catalytic sites. As protons pass from the intermembrane space through $F_o$ into the matrix, the $c$-ring rotates, driving the asymmetrical central $\gamma$-subunit rotor within the stationary $\alpha_3\beta_3$ hexamer. According to the Boyer binding-change mechanism, rotation drives sequential conformational changes in the three catalytic $\beta$-subunits through three distinct states: Open (ATP released, ADP + $P_i$ bind), Loose (substrates bound loosely), and Tight (substrates condensed into ATP). The complex is specifically inhibited by the antibiotic oligomycin, which binds the $F_o$ subunit and arrests both ATP synthesis and electron transport, while uncoupling proteins (UCP-1/thermogenin) dissipate the gradient as heat.

\medterm{Base}-- A substance that accepts protons (H+) or donates electron pairs (Lewis base). In aqueous solution, bases increase OH- concentration (pH >7). Strong bases (NaOH, KOH) dissociate completely; weak bases (NH3, HCO3-, protein side chains) partially. In biochemistry, key bases include: bicarbonate (HCO3-, major blood buffer), phosphate (HPO42-), amino acid side chains (histidine imidazole, lysine, arginine), and nucleotide bases (adenine, guanine, cytosine, thymine, uracil). Base excess/deficit quantifies metabolic acid-base disturbance.

\medterm{Base Excision Repair}-- The repair pathway for small, non-helix-distorting base lesions including oxidized bases (8-oxoguanine), deaminated bases (uracil from cytosine deamination), alkylated bases, and abasic (AP) sites. Steps: (1) a specific DNA glycosylase recognizes and removes the damaged base by hydrolyzing the N-glycosidic bond, creating an AP (apurinic/apyrimidinic) site. Different glycosylases recognize different lesions: uracil-DNA glycosylase (UNG) removes uracil from DNA; OGG1 removes 8-oxoguanine; MPG (methylpurine DNA glycosylase) removes alkylated bases. After glycosylase action, AP endonuclease 1 (APE1) nicks the DNA backbone 5' to the AP site, DNA polymerase beta replaces the missing nucleotide, and DNA ligase seals the nick.

\medterm{Beriberi}-- A nutritional deficiency disease caused by inadequate thiamine (vitamin B1), occurring in two forms. Wet beriberi presents with high-output heart failure, peripheral vasodilation, edema, and tachycardia due to impaired cardiac energy metabolism. Dry beriberi presents with peripheral neuropathy (symmetric sensory and motor deficits, paresthesias, muscle wasting). Risk factors: polished rice as dietary staple, alcoholism, malabsorption, and pregnancy. Treatment: thiamine supplementation (parenteral initially for severe cases). Wernicke-Korsakoff syndrome is the neuropsychiatric manifestation of thiamine deficiency in alcoholism.

\medterm{Beta-Oxidation}-- The catabolic process by which fatty acid molecules are broken down in the mitochondria (and peroxisomes) to generate acetyl-CoA, NADH, and FADH2. Fatty acids are activated to acyl-CoA by acyl-CoA synthetase, transported into the mitochondrial matrix via the carnitine shuttle, and undergo a four-step repeating cycle: oxidation, hydration, oxidation, and thiolysis. Yields 106 ATP net per molecule of palmitate (16-carbon fatty acid).

\medterm{Bilirubin}-- A linear tetrapyrrole yellow catabolic pigment derived primarily from the physiological breakdown of heme (80--85\% originating from senescent red blood cell hemoglobin degradation by reticuloendothelial macrophages in the spleen and liver, with 15--20\% derived from ineffective erythropoiesis and non-erythroid hemeproteins such as myoglobin and cytochrome P450). Macrophage heme oxygenase cleaves heme into biliverdin, iron, and carbon monoxide; biliverdin reductase subsequently reduces biliverdin to unconjugated (indirect) bilirubin. Highly lipophilic and insoluble in aqueous plasma, unconjugated bilirubin circulates tightly bound to albumin. In the liver, hepatocytes take up bilirubin via OATP1B1/1B3 transporters and conjugate it with two molecules of glucuronic acid via the microsomal enzyme UDP-glucuronosyltransferase (UGT1A1), forming water-soluble conjugated (direct) bilirubin diglucuronide. Conjugated bilirubin is actively secreted across the canalicular membrane into bile by the ATP-dependent transporter MRP2. Normal total serum bilirubin is 0.2--1.2 mg/dL ($<0.3$ mg/dL direct). Accumulation in blood ($>2.5--3.0$ mg/dL) deposits in elastin-rich tissues, causing jaundice and scleral icterus. Unconjugated hyperbilirubinemia results from hemolysis, Gilbert syndrome (mild UGT1A1 promoter defect), or Crigler-Najjar syndrome (severe/absent UGT1A1, risking kernicterus in infants). Conjugated hyperbilirubinemia reflects biliary obstruction (gallstones, pancreatic cancer) or canalicular transporter defects (Dubin-Johnson syndrome).

\medterm{Biomolecule}-- Any organic molecule produced by living organisms, essential for structure, function, and regulation of biological processes. Four major classes: (1) Proteins (amino acid polymers--enzymes, structural, transport, signaling); (2) Nucleic acids (DNA, RNA--genetic information storage and expression); (3) Carbohydrates (sugars, polysaccharides--energy, structure); (4) Lipids (fats, phospholipids, steroids--membranes, energy storage, signaling). Biomolecules are built from a small set of elements (C, H, O, N, P, S) and share universal metabolic pathways across life.

\medterm{Biotin (Vitamin B7)}-- A water-soluble B-complex vitamin containing an imidazole and thiophene ring fused together with a valeric acid side chain. Biotin serves as an indispensable, covalently bound prosthetic cofactor for four major carboxylase enzymes in human metabolism, catalyzing ATP-dependent carbon dioxide ($CO_2$) transfer: (1) Pyruvate Carboxylase: converts pyruvate to oxaloacetate in mitochondria, replenishing citric acid cycle intermediates (anaplerosis) and initiating gluconeogenesis; (2) Acetyl-CoA Carboxylase: converts acetyl-CoA to malonyl-CoA in the cytoplasm, constituting the rate-limiting committed step of de novo fatty acid synthesis; (3) Propionyl-CoA Carboxylase: converts propionyl-CoA to methylmalonyl-CoA during the catabolism of odd-chain fatty acids and branched-chain amino acids (valine, isoleucine, methionine, threonine); and (4) $\beta$-Methylcrotonyl-CoA Carboxylase: required for leucine degradation. The enzyme holocarboxylase synthetase covalently attaches biotin to a specific lysine $\epsilon$-amino group on the apoenzyme, forming biocytin, while biotinidase cleaves and recycles biotin during protein turnover. Nutritional deficiency is rare because biotin is produced by colonic microflora, but occurs with prolonged parenteral nutrition without vitamins or excessive consumption of raw egg whites (which contain avidin, a glycoprotein that binds biotin with femtomolar affinity, preventing absorption). Deficiency causes perioral and periorbital erythematous scaly dermatitis, diffuse alopecia, glossitis, paresthesias, depression, and lactic acidosis.

\medterm{Blood Urea Nitrogen}-- A clinical biochemical measurement quantifying the nitrogen concentration present as urea circulating in the blood, reflecting the balance between hepatic amino acid deamination and renal clearance. Urea is synthesized exclusively in hepatocytes via the mitochondrial and cytosolic reactions of the urea cycle to neutralize neurotoxic ammonia. Normal serum BUN concentration ranges from 7 to 20 mg/dL (2.5--7.1 mmol/L). Because urea is freely filtered at the glomerulus and reabsorbed (40--50\%) by passive diffusion in the renal tubules (which increases markedly during states of reduced tubular flow, volume depletion, and renal hypoperfusion), BUN is a sensitive but non-specific indicator of renal clearance. The serum BUN-to-creatinine ratio is of profound diagnostic value: a ratio $>20:1$ indicates prerenal azotemia (dehydration, congestive heart failure, hemorrhage, or excessive protein catabolism), where enhanced tubular reabsorption outpaces creatinine clearance; a ratio $<15:1$ indicates intrinsic renal failure (acute tubular necrosis, glomerulonephritis), where tubular damage prevents urea reabsorption; and postrenal obstruction frequently causes an initial BUN/Cr ratio $>20:1$ that normalizes as parenchymal injury develops.

\medterm{BNP}-- B-type (or brain) natriuretic peptide, a 32-amino-acid peptide hormone synthesized, stored, and secreted predominantly by ventricular cardiomyocytes in response to increased myocardial wall stretch, volume overload, and intracardiac filling pressures. Synthesized as a 134-amino-acid precursor (pre-proBNP), it is cleaved into the prohormone proBNP (108 amino acids) and subsequently processed by the endoprotease corin or furin into biologically active BNP (32 amino acids, C-terminal) and an equimolar, biologically inert N-terminal fragment (NT-proBNP, 76 amino acids). Biologically active BNP binds to natriuretic peptide receptor-A (NPR-A), activating transmembrane guanylyl cyclase to elevate intracellular cyclic GMP (cGMP), which induces systemic venous and arterial vasodilation, promotes renal glomerular filtration, stimulates natriuresis and diuresis, and suppresses the renin-angiotensin-aldosterone system (RAAS) and sympathetic tone. Clinically, serum BNP and NT-proBNP are essential biomarkers for evaluating acute dyspnea: a BNP $<100$ pg/mL (or NT-proBNP $<300$ pg/mL) demonstrates a $>95\%$ negative predictive value, effectively ruling out acute decompensated heart failure. Values are falsely elevated in chronic kidney disease (reduced clearance), female sex, and elderly patients, and paradoxically reduced in obese individuals.

\medterm{Buffer}-- A solution that resists pH change upon addition of acid or base, composed of a weak acid and its conjugate base (or weak base and conjugate acid). Buffer capacity is maximal when pH = pKa ± 1. Major physiological buffers: (1) Bicarbonate/carbonic acid (HCO3-/H2CO3, pKa 6.1, primary blood buffer, linked to respiratory CO2 elimination); (2) Phosphate (H2PO4-/HPO42-, pKa 7.2, intracellular and urinary buffer); (3) Protein buffers (hemoglobin, albumin, pKa ~6.8-7.8, major intracellular and plasma buffer); (4) Ammonia/ammonium (NH3/NH4+, renal buffer). Henderson-Hasselbalch equation: pH = pKa + log([base]/[acid]).

\medterm{Calcium}-- The most abundant mineral cation ($Ca^{2+}$) in the human body (totaling ~1,000 g in an adult, 99\% sequestered as hydroxyapatite in bones and teeth). Total serum calcium is maintained within an exceptionally narrow reference range of 8.5 to 10.5 mg/dL (2.15--2.55 mmol/L), distributed across three distinct fractions: (1) Free Ionized Calcium ($Ca^{2+}$, 50\%, 1.15--1.33 mmol/L): the sole biologically active fraction, directly regulated by parathyroid hormone (PTH) and calcitriol; (2) Protein-Bound Calcium (40\%): predominantly bound to negatively charged albumin, fluctuating directly with serum albumin concentration (corrected calcium = measured total calcium + $0.8 \times [4.0 - \text{serum albumin in g/dL}]$); and (3) Complexed Calcium (10\%): chelated to citrate, phosphate, and sulfate. Alterations in systemic pH shift calcium-albumin binding: in metabolic alkalemia, excess negative charges on deprotonated albumin bind ionized $Ca^{2+}$, triggering acute symptomatic hypocalcemia (paresthesias, tetany, Chvostek and Trousseau signs) without changing total calcium. Intracellularly, cytosolic calcium is maintained at nanomolar concentrations ($~100$ nM), functioning as a ubiquitous second messenger: influx through voltage-gated or receptor-operated channels binds calmodulin, activates protein kinase C (PKC), triggers neurotransmitter exocytosis, drives excitation-contraction coupling in striated and smooth muscle, and complexes with vitamin K-dependent Gla residues to anchor coagulation factors on platelet membranes.

\medterm{Carbohydrate}-- A class of organic compounds composed of carbon, hydrogen, and oxygen (general formula Cm(H2O)n), serving as the primary energy source for the body. Classification: monosaccharides (glucose, fructose, galactose--single sugar units), disaccharides (sucrose, lactose, maltose--two units), oligosaccharides (3-10 units), and polysaccharides (starch, glycogen, cellulose, fiber--many units). Glucose is the essential fuel for the brain and red blood cells. The body stores limited glycogen (liver ~100 g, muscle ~400 g). Dietary carbohydrates are digested to monosaccharides, absorbed in the small intestine, and metabolized via glycolysis, the TCA cycle, and oxidative phosphorylation to produce ATP. Excess carbohydrate is converted to fat (lipogenesis). The glycemic index ranks foods by their effect on blood glucose.

\medterm{Cardiac Biomarkers}-- A panel of endogenous intracellular myocardial enzymes, contractile proteins, and neurohormones released into the peripheral circulation following myocardial injury, necrosis, or hemodynamic ventricular wall stress, utilized in emergency and cardiovascular medicine to diagnose acute coronary syndromes (ACS) and heart failure. The modern gold standard is high-sensitivity cardiac Troponin I and T (hs-cTnI and hs-cTnT), contractile regulatory proteins possessing near-absolute myocardial tissue specificity; following sarcolemmal rupture from ischemic necrosis, hs-cTn levels rise within 2 to 3 hours, peak at 12 to 24 hours, and persist for 7 to 14 days. Creatine Kinase-MB (CK-MB), a dimeric enzyme present in myocardium (15--20\% of total cardiac CK) and skeletal muscle, rises within 4 to 6 hours, peaks at 18 to 24 hours, and normalizes within 48 to 72 hours; its rapid clearance makes CK-MB specifically useful for detecting early reinfarction following an initial acute event. Myoglobin is a small hemeprotein that rises exceptionally early (within 1--2 hours) but lacks cardiac specificity (released during any skeletal muscle injury). For hemodynamic wall stress and volume overload, B-type natriuretic peptide (BNP) and N-terminal proBNP (NT-proBNP) are synthesized by stressed ventricular myocytes, serving as definitive biomarkers to rule in or rule out acute decompensated heart failure.

\medterm{Carnitine}-- An essential quaternary ammonium compound ($\beta$-hydroxy-$\gamma$-trimethylaminobutyrate) that functions as an indispensable carrier molecule for the transport of activated long-chain fatty acids (LCFAs) across the impermeable inner mitochondrial membrane for beta-oxidation. Carnitine is obtained from dietary red meat and dairy products (75\%) or synthesized endogenously (25\%) in the liver and kidneys from lysine and methionine residues, requiring ascorbic acid (vitamin C), iron, and pyridoxine as cofactors. Because the inner mitochondrial membrane lacks a transporter for long-chain acyl-CoA esters, the carnitine shuttle system operates sequentially: carnitine palmitoyltransferase-1 (CPT-1) on the outer mitochondrial membrane conjugates the acyl group to carnitine, forming acylcarnitine; carnitine-acylcarnitine translocase (CACT) shuttles acylcarnitine across the inner membrane in exchange for free carnitine; and CPT-2 on the inner membrane reconverts acylcarnitine back into matrix acyl-CoA and free carnitine. Primary carnitine deficiency (caused by autosomal recessive mutations in the OCTN2 carnitine transporter gene, SLC22A5) impairs renal tubular reabsorption and cellular uptake, leading to massive urinary carnitine wasting. Patients present in early childhood with hypoketotic hypoglycemia during fasting, encephalopathy, hepatomegaly, skeletal muscle weakness, and dilated cardiomyopathy. Treatment requires lifelong high-dose oral L-carnitine supplementation.

\medterm{Carnitine Palmitoyltransferase (CPT)}-- A pair of mitochondrial membrane transferase enzymes (CPT-1 and CPT-2) that catalyze the reversible transesterification between acyl-CoA and carnitine, constituting the rate-limiting machinery for long-chain fatty acid entry into the mitochondrial matrix for beta-oxidation. CPT-1 resides on the outer mitochondrial membrane and transfers the fatty acyl moiety from cytosolic acyl-CoA to carnitine, generating acylcarnitine and releasing free CoA-SH. CPT-1 is the primary committed, rate-limiting control point for fatty acid oxidation; it is potently allosterically inhibited by malonyl-CoA (the product of acetyl-CoA carboxylase), ensuring that fatty acid oxidation is suppressed whenever active lipogenesis is occurring in the fed state. Three tissue-specific isoforms exist: CPT-1A (liver), CPT-1B (skeletal muscle and heart), and CPT-1C (brain). CPT-2 resides on the inner mitochondrial membrane facing the matrix, where it reverses the reaction, converting acylcarnitine and matrix CoA-SH into acyl-CoA (which enters beta-oxidation) and free carnitine. CPT-1 deficiency causes severe fasting hypoketotic hypoglycemia and hepatomegaly, while CPT-2 deficiency (the most common inherited disorder of lipid metabolism in adults) presents as exercise-induced myalgia, muscle stiffness, rhabdomyolysis, and myoglobinuria triggered by fasting, cold exposure, or prolonged exertion.

\medterm{Catabolism}-- The set of metabolic pathways that break down complex molecules into simpler units, releasing energy captured as ATP, NADH, FADH2, NADPH. Examples: glycolysis, beta-oxidation, proteolysis, glycogenolysis, nucleotide catabolism. Catabolic pathways converge on central intermediates (acetyl-CoA, pyruvate, TCA cycle intermediates) and are typically stimulated by glucagon, epinephrine, cortisol, and AMPK; inhibited by insulin. The energy released drives anabolic processes, ion gradients, and mechanical work.

\medterm{Catalyst}-- A substance that increases the rate of a chemical reaction without being consumed, by lowering the activation energy (providing an alternative reaction pathway). Enzymes are biological catalysts (proteins or ribozymes). Catalysts do not alter the reaction equilibrium ($\Delta G$), only the kinetics. They accelerate both forward and reverse reactions equally. Enzyme catalysis mechanisms: proximity/orientation, acid-base, covalent, metal ion, electrostatic stabilization, and strain/distortion.

\medterm{Celiac Disease}-- A chronic, immune-mediated enteropathy of the small intestine triggered by the ingestion of dietary gluten (a storage protein composite found in wheat, rye, and barley) in genetically predisposed individuals carrying human leukocyte antigen HLA-DQ2 or HLA-DQ8. Gluten peptides rich in proline and glutamine resist gastric and pancreatic endopeptidases; undigested 33-mer gliadin peptides cross the intestinal epithelial barrier into the lamina propria, where tissue transglutaminase (tTG-2) deamidates neutral glutamine residues into negatively charged glutamic acid. Deamidated gliadin fits with high affinity into the peptide-binding grooves of HLA-DQ2/DQ8 on antigen-presenting cells, initiating a vigorous CD4+ Th1 response (secreting IFN-$\gamma$) and activating cytotoxic CD8+ intraepithelial lymphocytes (IELs) that destroy enterocytes. This leads to mucosal inflammation, crypt hyperplasia, and progressive blunting and atrophy of duodenal and jejunal villi. Classic presentation includes malabsorption, chronic osmotic diarrhea, steatorrhea, weight loss, iron-deficiency anemia, and fat-soluble vitamin (A, D, E, K) deficiencies, associated with dermatitis herpetiformis. Serologic screening is performed using IgA anti-tissue transglutaminase (anti-tTG) and anti-endomysial antibodies; confirmation requires duodenal biopsy (Marsh classification). A strict lifelong gluten-free diet achieves complete clinical and histological recovery.

\medterm{Cholesterol}-- A 27-carbon sterol essential for cell membrane structure (modulates fluidity), precursor of steroid hormones (cortisol, aldosterone, sex hormones), bile acids, and vitamin D. Synthesized de novo (~1 g/day, primarily liver) from acetyl-CoA via the mevalonate pathway (HMG-CoA reductase is rate-limiting, targeted by statins). Obtained from diet (animal products). Transported in blood as lipoproteins (LDL, HDL). Cellular uptake via LDL receptor (endocytosis). Excess cholesterol is esterified (ACAT) for storage or effluxed to HDL (ABCA1) for reverse cholesterol transport.

\medterm{Cholesterol Synthesis}-- Occurs in all nucleated cells (primarily liver and intestine, approximately 1 g/day production), using acetyl-CoA as the precursor through approximately 30 enzymatic steps. Key stages: (1) two acetyl-CoA condense to acetoacetyl-CoA, then to HMG-CoA (3-hydroxy-3-methylglutaryl-CoA); (2) HMG-CoA reductase (ER-membrane enzyme, the rate-limiting step) reduces HMG-CoA to mevalonate using 2 NADPH. This step is inhibited by statins (competitive inhibitors, the most widely prescribed cholesterol-lowering drugs). After mevalonate, isopentenyl pyrophosphate (IPP) is formed, and six IPP molecules condense to form squalene, which is cyclized to lanosterol and then converted through nineteen additional steps to cholesterol.

\medterm{Chromatin}-- The macromolecular complex of genomic DNA, basic histone proteins, and non-histone chromosomal proteins localized within the eukaryotic nucleus that packages the linear genome, prevents DNA entanglement, protects against damage, and dynamically regulates transcription, replication, and repair. The fundamental repeating structural unit of chromatin is the nucleosome, consisting of 146--147 base pairs of negatively charged DNA wrapped approximately 1.65 times around an octameric core of positively charged basic histone proteins (two copies each of H2A, H2B, H3, and H4). Linker histone H1 binds the entry/exit site of DNA, facilitating higher-order folding into a 30-nm chromatin fiber and looped chromosomal domains. Chromatin exists in two distinct structural and functional states: (1) Euchromatin: a dispersed, lightly packed, transcriptionally active conformation characterized by hyperacetylated histone tails (relaxed DNA accessible to transcription factors and RNA polymerase II); and (2) Heterochromatin: a condensed, densely staining, transcriptionally silent conformation, subdivided into constitutive heterochromatin (permanently silent repetitive regions, such as centromeres and telomeres, enriched in H3K9me3) and facultative heterochromatin (developmentally regulated silenced regions, such as the inactive female X chromosome [Barr body], enriched in H3K27me3 and dense DNA CpG methylation).

\medterm{Chylomicron}-- The largest (75--1,200 nm diameter) and least dense ($<0.95$ g/mL) class of plasma lipoproteins, synthesized exclusively by enterocytes in the small intestine to package and transport dietary (exogenous) lipids into systemic circulation. Composed of an apolar core rich in triacylglycerols ($>85\%$ by mass) and cholesteryl esters, surrounded by an amphipathic phospholipid monolayer containing free cholesterol and the structural signature protein apolipoprotein B-48 (ApoB-48), produced by post-transcriptional cytidine-to-uridine mRNA editing of the APOB transcript by APOBEC-1. Nascent chylomicrons are secreted into intestinal lacteals, travel through the thoracic duct, and enter the subclavian vein, thereby bypassing the liver on initial transit. In the systemic circulation, chylomicrons acquire apolipoprotein C-II (ApoC-II, an obligatory activator of lipoprotein lipase) and apolipoprotein E (ApoE, for receptor uptake) from circulating HDL particles. At capillary endothelial surfaces in adipose tissue, skeletal muscle, and myocardium, lipoprotein lipase (LPL) hydrolyzes core triglycerides into free fatty acids and glycerol. The resulting chylomicron remnants, depleted of core triglycerides but retaining ApoB-48, ApoE, and cholesteryl esters, are cleared rapidly and completely by hepatic LDL receptor-related protein-1 (LRP-1) and LDL receptors via ApoE-mediated receptor endocytosis.

\medterm{Citrate Synthase}-- The initial pacemaker enzyme of the citric acid cycle (Krebs cycle), localized entirely within the mitochondrial matrix. It catalyzes the irreversible condensation of oxaloacetate (4 carbons) with the acetyl moiety of acetyl-CoA (2 carbons) and water to yield citrate (6 carbons) and free coenzyme A (CoA-SH). The reaction is driven strongly forward by the exergonic cleavage of the high-energy thioester bond ($\Delta G^{\circ\prime} = -31.5$ kJ/mol). Citrate synthase follows an ordered bi-bi kinetic mechanism: oxaloacetate binds first, inducing a major conformational shift that closes the active site and creates the binding cleft for acetyl-CoA, preventing wasteful non-productive hydrolysis of acetyl-CoA. Activity is regulated primarily by substrate availability (oxaloacetate and acetyl-CoA) and allosterically inhibited by high ATP, NADH, succinyl-CoA, and long-chain acyl-CoA esters, reflecting high mitochondrial energy charge.

\medterm{Citric Acid Cycle (Krebs Cycle)}-- A central metabolic pathway occurring in the mitochondrial matrix that oxidizes acetyl-CoA (derived from carbohydrates, lipids, and proteins) into CO2, generating 3 NADH, 1 FADH2, and 1 GTP (ATP) per turn. Rate-limiting enzyme: isocitrate dehydrogenase (inhibited by ATP and NADH; activated by ADP). Also serves an amphibolic role by providing precursors for amino acids, heme, and fatty acid synthesis.

\medterm{Codon}-- A specific sequence of three contiguous ribonucleotide bases situated within a messenger RNA (mRNA) transcript that directs the incorporation of a specific amino acid into a growing nascent polypeptide chain or dictates the termination of translation. With four possible RNA bases (adenine [A], uracil [U], guanine [G], cytosine [C]), there are $4^3 = 64$ possible triplet codons. Sixty-one codons specify the 20 standard amino acids, while three codons (UAA, UAG, and UGA) function as non-coding Stop (nonsense) codons that bind protein release factors (RFs) to terminate ribosomal translation. AUG functions as the universal Start (initiation) codon, establishing the translational reading frame and encoding L-methionine (or N-formylmethionine in prokaryotes and mitochondria). The genetic code exhibits four foundational characteristics: (1) Degenerate / Redundant: multiple synonymous codons specify the same amino acid (e.g., leucine and arginine are each encoded by 6 codons), typically differing at the third 'wobble' position, which buffers against harmful missense mutations; (2) Unambiguous: each individual codon specifies only one amino acid; (3) Non-overlapping and Commaless: codons are read consecutively in triplets without gaps or shared nucleotides from the start codon to the stop codon; and (4) Nearly Universal: conserved across virtually all living organisms from bacteria to humans, with minor exceptions in mitochondrial DNA (e.g., UGA encodes tryptophan instead of stop).

\medterm{Coenzyme}-- A non-protein organic molecule that binds to an enzyme (apoenzyme) and is required for its catalytic activity. Coenzymes often act as transient carriers of specific chemical groups or electrons. Many are vitamin-derived: NAD+/NADP+ (niacin, B3), FAD/FMN (riboflavin, B2), coenzyme A (pantothenic acid, B5), thiamine pyrophosphate/TPP (B1), pyridoxal phosphate/PLP (B6), biotin (B7), tetrahydrofolate (B9), cobalamin (B12). Coenzymes may be tightly bound (prosthetic groups, e.g., FAD in succinate dehydrogenase) or loosely bound (cosubstrates, e.g., NAD+ in dehydrogenases).

\medterm{Cofactor}-- A non-protein chemical compound or metallic ion required for an enzyme's catalytic activity. Cofactors include: (1) Metal ions (Mg2+, Zn2+, Fe2+/Fe3+, Cu2+, Mn2+, Ca2+, Mo, Se) that stabilize structure, participate in redox, or act as Lewis acids; (2) Coenzymes (organic cofactors, often vitamin-derived). The complete active enzyme (apoenzyme + cofactor) is the holoenzyme. Cofactor deficiency causes metabolic dysfunction (e.g., Zn2+ for carbonic anhydrase, Mg2+ for kinases, Fe2+ for cytochromes).

\medterm{Collagen}-- The most abundant structural protein in the human body (comprising $>25--30\%$ of total body protein mass), providing tensile strength, mechanical durability, and architectural scaffolding to connective tissues, including bone, cartilage, tendons, skin, blood vessels, and the extracellular matrix. Collagen possesses a characteristic triple-helical quaternary structure composed of three polypeptide $\alpha$-chains featuring a repeating tripeptide sequence: $(\text{Gly}-X-Y)_n$, where $X$ is frequently proline (inducing a tight polyproline II twist) and $Y$ is frequently 4-hydroxyproline or hydroxylysine. Small glycine residues, occurring at every third position, fit into the restricted central axis of the triple helix. Collagen biosynthesis involves complex intracellular and extracellular steps: (1) Inside the rough ER, procollagen $\alpha$-chains are synthesized, and specific proline and lysine residues undergo enzymatic hydroxylation catalyzed by prolyl and lysyl hydroxylases (requiring ascorbic acid / vitamin C, $Fe^{2+}$, and $\alpha$-ketoglutarate as cofactors); (2) Selected hydroxylysines undergo glycosylation with glucose and galactose, and three $\alpha$-chains assemble via C-terminal disulfide bonding into a triple helix; (3) Following Golgi exocytosis, extracellular procollagen peptidases cleave the terminal non-helical registration propeptides, generating insoluble tropocollagen; and (4) Lysyl oxidase (a copper-dependent enzyme) deaminates lysine residues into reactive allysine aldehydes, which condense covalently to form mature, cross-linked collagen fibrils. Major types include: Type I (bone, skin, tendon, dentin, defective in Osteogenesis Imperfecta); Type II (hyaline and elastic cartilage, vitreous humor); Type III (reticulin fibers, blood vessels, skin, defective in vascular Ehlers-Danlos syndrome); and Type IV (basement membranes, defective in Alport syndrome). Vitamin C deficiency impairs hydroxylation, causing Scurvy.

\medterm{Cori Cycle}-- Also designated the lactic acid cycle, a crucial inter-organ metabolic pathway involving the coordinated cooperation of peripheral tissues undergoing anaerobic glycolysis (particularly vigorously contracting skeletal muscle and mature red blood cells lacking mitochondria) with the liver. Under conditions of tissue hypoxia or intense muscular exertion where rate of ATP demand outstrips oxidative phosphorylation, glycolysis rapidly converts glucose to pyruvate, generating a net of 2 ATP molecules per glucose. To sustain glycolytic flux through glyceraldehyde-3-phosphate dehydrogenase, lactate dehydrogenase (LDH) reduces pyruvate to L-lactate, oxidizing NADH back to $NAD^+$. Because accumulation of lactic acid in contracting muscle would cause severe intracellular acidosis and inhibit phosphofructokinase-1, lactate is exported into the bloodstream via monocarboxylate transporters (MCTs) and carried to the liver. In hepatocytes, the high $NAD^+/NADH$ ratio drives hepatic LDH to oxidize lactate back to pyruvate with the formation of NADH. Hepatocytes then convert pyruvate back into glucose via gluconeogenesis, consuming 6 high-energy phosphate bonds (4 ATP and 2 GTP). The newly synthesized glucose is released into the systemic circulation to replenish muscle glycogen stores or sustain red blood cell energy needs. The Cori cycle effectively shifts the energetic burden and metabolic acidosis of anaerobic work from muscle to the liver at a net systemic cost of 4 ATP equivalents.

\medterm{Cortisol}-- The primary endogenous glucocorticoid steroid hormone, synthesized from cholesterol in the zona fasciculata of the adrenal cortex under the regulatory control of the hypothalamic-pituitary-adrenal (HPA) axis via hypothalamic CRH and pituitary ACTH. Lipophilic cortisol diffuses across plasma membranes and binds cytoplasmic glucocorticoid receptors (GR), inducing chaperone dissociation, receptor homodimerization, and nuclear translocation where it binds glucocorticoid response elements (GREs) to modulate gene transcription. Cortisol is an indispensable stress and counter-regulatory catabolic hormone: in hepatocytes, it transcriptionally induces key gluconeogenic enzymes (PEPCK and glucose-6-phosphatase) to elevate blood glucose; in peripheral tissues, it stimulates proteolysis in skeletal muscle (supplying glucogenic amino acid precursors) and lipolysis in adipocytes. Cortisol exerts profound anti-inflammatory and immunosuppressive actions by inducing annexin-1 (lipocortin), which inhibits phospholipase $A_2$ (blocking arachidonic acid liberation and subsequent prostaglandin and leukotriene synthesis), and by repressing NF-$\kappa$B, halting cytokine transcription (IL-1, IL-2, IL-6, TNF-$\alpha$). Pathologically, chronic excess causes Cushing syndrome (central obesity, moon facies, buffalo hump, purple striae, osteoporosis, muscle wasting, and hyperglycemia); primary adrenal insufficiency causes Addison disease (hypotension, hyperkalemia, hypoglycemia, and skin hyperpigmentation due to elevated ACTH/POMC).

\medterm{Creatinine}-- A cyclic nitrogenous waste product formed non-enzymatically at a constant daily rate (approximately 1--2\% per day) from the spontaneous, irreversible dehydration and dephosphorylation of creatine and creatine phosphate in skeletal muscle. Normal serum creatinine concentrations range from 0.7 to 1.3 mg/dL in adult men and 0.5 to 1.1 mg/dL in adult women, directly proportional to total functional muscle mass. Creatinine is freely filtered across the glomerular filtration barrier, is not reabsorbed by renal tubular cells, and undergoes only minor active secretion (10--15\%) by organic cation transporters (OCT2) in the proximal convoluted tubule. Because of these physiological properties, serum creatinine and 24-hour urinary creatinine clearance serve as the universal foundation for estimating glomerular filtration rate (eGFR) via the CKD-EPI and MDRD equations. Due to the non-linear hyperbolic relationship between serum creatinine and GFR, an acute doubling of serum creatinine reflects an approximately 50\% loss of functioning nephron mass. Levels are elevated in acute kidney injury, chronic kidney disease, and rhabdomyolysis, and artificially elevated by drugs that inhibit tubular secretion (cimetidine, trimethoprim).

\medterm{CRP}-- C-reactive protein, a 118-kDa cyclic pentameric acute-phase reactant belonging to the pentraxin protein superfamily, synthesized and secreted rapidly by hepatocytes predominantly under the transcriptional stimulation of interleukin-6 (IL-6), with synergistic amplification by IL-1$\beta$ and TNF-$\alpha$. CRP serves as an innate pattern recognition receptor (PRR), binding with high, $Ca^{2+}$-dependent affinity to phosphocholine residues exposed on the surface of apoptotic host cells and microbial polysaccharides (such as the C-polysaccharide of Streptococcus pneumoniae). Surface-bound CRP binds C1q to activate the classical complement cascade and engages Fc$\gamma$ receptors on macrophages and neutrophils, acting as an opsonin to promote targeted phagocytosis. Serum CRP levels rise precipitously (up to 1,000-fold, exceeding 100 mg/L) within 6 to 8 hours of acute infection, tissue infarction (e.g., myocardial infarction), surgical trauma, or autoimmune flare-ups, and decay rapidly (half-life $\approx 19$ hours) upon clinical resolution. High-sensitivity CRP (hs-CRP) assays measure basal subclinical vascular inflammation (normal $<1.0$ mg/L, average risk 1.0--3.0 mg/L, high risk $>3.0$ mg/L), functioning as an independent biomarker for atherothrombotic cardiovascular risk.

\medterm{Cytochrome c}-- A small (12 kDa), highly conserved, water-soluble peripheral membrane hemeprotein loosely associated with the outer surface of the inner mitochondrial membrane within the intermembrane space. In the electron transport chain, cytochrome c operates as a mobile single-electron shuttle, accepting an electron from ubiquinol-cytochrome c oxidoreductase (Complex III) via its central heme iron ($Fe^{3+} \to Fe^{2+}$) and transferring it to cytochrome c oxidase (Complex IV), where oxygen is reduced to water. Beyond bioenergetics, cytochrome c is the pivotal executioner messenger in the intrinsic (mitochondrial) pathway of apoptosis. Under cellular stress, DNA damage, or deprivation of survival factors, pro-apoptotic Bcl-2 proteins (Bax and Bak) oligomerize to permeabilize the outer mitochondrial membrane (MOMP). Released into the cytoplasm, cytochrome c binds apoptotic protease activating factor-1 (Apaf-1) in the presence of dATP, triggering the assembly of the heptameric apoptosome, which recruits and cleaves procaspase-9, launching the downstream proteolytic executioner caspase cascade (caspases 3 and 7).

\medterm{D-dimer}-- A specific soluble fibrin degradation product composed of two covalently cross-linked D domains of adjacent fibrin monomers joined by an isopeptide bond introduced by activated factor XIII (FXIIIa). During physiological coagulation, thrombin cleaves fibrinogen to form fibrin monomers that assemble into polymers, which FXIIIa covalently stabilizes via intermolecular covalent cross-links; subsequent fibrinolysis by the serine protease plasmin cleaves this fibrin mesh, releasing various soluble fragments, among which the D-dimer moiety is unique to cross-linked fibrin. Circulating D-dimer is measured quantitatively via monoclonal antibody immunoassay (normal threshold $<500$ ng/mL fibrinogen equivalent units [FEU]). Because D-dimer generation demands both active thrombin generation (clot formation) and plasmin activity (clot breakdown), it represents an exceptionally sensitive marker of intravascular thrombosis. It is clinically utilized for its outstanding negative predictive value ($>95\%$): a normal D-dimer reliably excludes deep vein thrombosis (DVT) and pulmonary embolism (PE) in patients with low-to-intermediate pretest probability. Elevated levels occur in venous thromboembolism, disseminated intravascular coagulation (DIC), sepsis, malignancy, trauma, and pregnancy.

\medterm{DNA (Deoxyribonucleic Acid)}-- Double-stranded helical molecule carrying genetic information. Composed of nucleotides (sugar-phosphate backbone with bases adenine, thymine, guanine, cytosine). Replication is semi-conservative. Genes are specific sequences encoding proteins or functional RNAs. Chromosomes are highly condensed DNA-protein complexes. Human genome: ~3 billion base pairs, ~20,000-25,000 genes.

\medterm{DNA Repair}-- A collection of conserved cellular processes that detect, excise, and correct damaged nucleotides and backbone lesions within the DNA double helix to preserve genomic integrity and prevent oncogenesis. Five major repair pathways exist: (1) Base Excision Repair (BER): DNA glycosylases recognize and cleave abnormal bases (uracil, 8-oxoguanine), generating an AP site repaired by AP endonuclease, DNA polymerase $\beta$, and ligase; (2) Nucleotide Excision Repair (NER): excinuclease endonuclease complexes excise bulky, helix-distorting lesions (such as UV-induced thymine dimers and chemical adducts), defective in Xeroderma Pigmentosum; (3) Mismatch Repair (MMR): recognizes unpaired bases and insertion/deletion loops post-replication during S-phase, defective in Lynch Syndrome (hereditary non-polyposis colorectal cancer) due to mutations in MLH1, MSH2, MSH6, or PMS2, causing microsatellite instability (MSI); (4) Homologous Recombination (HR): an error-free mechanism repairing double-strand breaks during late S/G2 phase using the homologous sister chromatid as a template, requiring BRCA1, BRCA2, ATM, and RAD51 proteins; and (5) Non-Homologous End Joining (NHEJ): directly ligates double-strand breaks without a homologous template throughout the cell cycle, error-prone and frequently mutagenic, defective in severe combined immunodeficiency (SCID) and ataxia-telangiectasia.

\medterm{Electrolytes}-- Minerals in human bodily fluids that dissociate into electrically charged ions, indispensable for maintaining osmotic equilibrium, fluid distribution between intracellular and extracellular compartments, acid-base buffering, neuromuscular excitability, and enzymatic catalysis. The major extracellular electrolytes are sodium ($Na^+$, 135--145 mEq/L, the primary determinant of plasma osmolality and extracellular fluid volume) and chloride ($Cl^-$, 96--106 mEq/L, maintaining electrical neutrality and acid-base balance), alongside bicarbonate ($HCO_3^-$, 22--28 mEq/L). The dominant intracellular electrolytes are potassium ($K^+$, intracellular $\approx 150$ mEq/L vs. extracellular 3.5--5.0 mEq/L, maintaining resting membrane potentials and cardiac conduction via the $Na^+/K^+$-ATPase), magnesium ($Mg^{2+}$, essential kinase and ATP chelation cofactor), and organic phosphates. Divalent calcium ($Ca^{2+}$, total 8.5--10.5 mg/dL, ionized 1.15--1.33 mmol/L) regulates excitation-contraction coupling, neurotransmission, and hemostasis. Electrolyte concentrations are tightly regulated by the kidneys, lungs, and endocrine hormones (aldosterone, antidiuretic hormone [ADH], parathyroid hormone [PTH], and natriuretic peptides); pathological derangements provoke life-threatening cardiac arrhythmias, seizures, and neuromuscular failure.

\medterm{Electron Transport Chain}-- Series of protein complexes (I-IV) and mobile carriers (ubiquinone, cytochrome c) in the inner mitochondrial membrane that transfer electrons from NADH and FADH2 to oxygen, pumping protons to create electrochemical gradient used by ATP synthase. Complex I (NADH dehydrogenase): NADH -> ubiquinone. Complex II (succinate dehydrogenase): FADH2 -> ubiquinone. Complex III (cytochrome bc1): ubiquinone -> cytochrome c. Complex IV (cytochrome c oxidase): cytochrome c -> O2 -> H2O. Uncouplers (2,4-dinitrophenol, thermogenin) dissipate proton gradient as heat.

\medterm{Enzyme}-- Biological catalyst (protein or RNA) that accelerates chemical reactions by lowering activation energy. Specificity: lock-and-key or induced fit models. Kinetics: Michaelis-Menten equation (Vmax, Km). Regulation: allosteric control, covalent modification (phosphorylation), proteolytic activation, feedback inhibition. Inhibition: competitive, non-competitive, uncompetitive. Clinical: enzyme levels as diagnostic markers (ALT, AST, CK, amylase, lipase, troponin).

\medterm{Enzyme Inhibition}-- The reduction or abolition of an enzyme's catalytic rate by specific chemical agents through reversible or irreversible molecular interactions, fundamentally characterized by alterations in Michaelis-Menten kinetic parameters ($K_m$ and $V_{\max}$) and Lineweaver-Burk double-reciprocal plots: (1) Competitive Inhibition: The inhibitor structurally resembles the substrate and binds reversibly to the free enzyme's active site, competing directly with substrate. Apparent $K_m$ increases (reduced apparent affinity), while $V_{\max}$ remains unchanged because excess substrate can outcompete the inhibitor (e.g., statins inhibiting HMG-CoA reductase; methotrexate inhibiting dihydrofolate reductase); (2) Non-Competitive Inhibition: The inhibitor binds to an allosteric regulatory site distinct from the active site with equal affinity for free enzyme and enzyme-substrate complex. $V_{\max}$ decreases (catalytic turnover is impaired), while $K_m$ remains unchanged because substrate binding affinity is unaltered and cannot overcome the inhibition (e.g., lead inhibiting ferrochelatase); (3) Uncompetitive Inhibition: The inhibitor binds exclusively to the enzyme-substrate ($ES$) complex, locking the substrate in the catalytic pocket; both $V_{\max}$ and $K_m$ decrease proportionally, maintaining a constant $K_m/V_{\max}$ slope; and (4) Irreversible Inhibition: The inhibitor forms covalent bonds with catalytic residues or destroys essential active-site functional groups, permanently destroying enzyme activity (e.g., aspirin acetylating Ser530 of COX-1/2; organophosphates phosphorylating the active serine of acetylcholinesterase).

\medterm{Enzyme Kinetics}-- Study of reaction rates and factors affecting them. Michaelis-Menten kinetics: V = Vmax[S]/(Km + [S]). Km (Michaelis constant): substrate concentration at half Vmax, inversely related to enzyme affinity. Vmax: maximum reaction rate at saturating substrate. Lineweaver-Burk plot: double-reciprocal plot (1/V vs. 1/[S]) for determining Km and Vmax. Inhibition patterns: competitive (increases Km, same Vmax), non-competitive (same Km, decreases Vmax), uncompetitive (decreases both).

\medterm{Epigenetics}-- The study of heritable, stable alterations in gene expression and cellular phenotype that occur without alterations to the underlying primary DNA nucleotide sequence. Major epigenetic mechanisms include: (1) DNA Methylation: DNA methyltransferases (DNMT1, DNMT3a, DNMT3b) transfer methyl groups from S-adenosylmethionine (SAM) to the 5-position of cytosine within CpG dinucleotides. Dense CpG island promoter methylation recruits methyl-CpG-binding domain proteins (MBDs) and histone deacetylases, condensing chromatin into transcriptionally silent heterochromatin; (2) Post-Translational Histone Modifications: Histone acetyltransferases (HATs) acetylate lysine residues on amino-terminal histone tails, neutralizing positive charges, relaxing chromatin into transcriptionally accessible euchromatin; histone deacetylases (HDACs) remove acetyl groups, driving transcriptional repression. Histone methylation can activate or repress depending on the residue (e.g., H3K4me3 activates, whereas H3K9me3 and H3K27me3 repress); (3) Genomic Imprinting: Parent-of-origin-specific gene silencing via differential DNA methylation during gametogenesis, exemplified by the 15q11-q13 locus (maternal imprint/paternal deletion causes Prader-Willi syndrome, whereas paternal imprint/maternal deletion causes Angelman syndrome); and (4) Non-Coding RNA Regulation: MicroRNAs and long non-coding RNAs (e.g., Xist in X-chromosome inactivation) that guide chromatin remodeling and post-transcriptional gene silencing.

\medterm{Epinephrine}-- Adrenaline, a primary catecholamine hormone synthesized and secreted by chromaffin cells of the adrenal medulla (comprising ~80\% of adrenal catecholamine output) under sympathetic nervous system stimulation via preganglionic acetylcholine acting on nicotinic receptors. Synthesized from L-tyrosine via L-DOPA, dopamine, and norepinephrine; the final step involves the methylation of norepinephrine by the cytoplasmic enzyme phenylethanolamine N-methyltransferase (PNMT), an enzyme heavily induced by cortisol reaching the medulla via intra-adrenal portal blood. Epinephrine acts as the master hormonal mediator of the acute 'fight-or-flight' stress response, exerting systemic effects by binding to multiple adrenergic receptor subtypes: (1) $\beta_1$-adrenergic receptors (myocardium): couples to $G_{\alpha s}$, increasing cAMP and PKA activity, producing positive inotropy (contractility), chronotropy (heart rate), and dromotropy (conduction velocity); (2) $\beta_2$-adrenergic receptors: couples to $G_{\alpha s}$, relaxing smooth muscle to cause bronchodilation, skeletal muscle vasodilation, and stimulating glycogenolysis and gluconeogenesis in liver and muscle, as well as lipolysis in adipose tissue (activating hormone-sensitive lipase); and (3) $\alpha_1$-adrenergic receptors: couples to $G_{\alpha q}$, increasing $IP_3$ and intracellular $Ca^{2+}$, causing potent peripheral cutaneous and splanchnic vasoconstriction and pupillary dilation (mydriasis). Clinically, intramuscular epinephrine is the first-line, life-saving emergency treatment for systemic anaphylaxis and a core component of cardiopulmonary resuscitation (ACLS).

\medterm{ESR}-- Erythrocyte sedimentation rate, a non-specific laboratory hematology test measuring the distance (in millimeters) that red blood cells settle in an anticoagulated column of whole blood (standard Westergren method) over a one-hour period (reference range: $<15$ mm/hr in men, $<20$ mm/hr in women, increasing with age: roughly age/2 for men, (age+10)/2 for women). Under basal conditions, erythrocytes carry a net negative surface electrical charge (zeta potential) mediated by sialic acid residues on glycophorin A, which causes them to repel one another and remain suspended. During systemic inflammation, acute-phase proteins (predominantly asymmetric, positively charged fibrinogen, as well as immunoglobulins and $\alpha_2$-macroglobulin) coat the erythrocyte membrane, neutralizing the zeta potential. This permits erythrocytes to stack into columnar rouleaux formations; because rouleaux have a smaller surface-area-to-volume ratio than individual cells, they settle through plasma significantly faster. Marked elevations ($>100$ mm/hr) are strongly suggestive of serious underlying pathology, such as temporal (giant cell) arteritis, polymyalgia rheumatica, multiple myeloma, systemic osteomyelitis, or metastatic malignancy. ESR responds and resolves much more slowly than CRP.

\medterm{Factor V Leiden}-- The most common inherited thrombophilia in Caucasian populations (carrier frequency 3--8\%), caused by a single point mutation in the F5 gene ($1691G \to A$) that results in an arginine-to-glutamine amino acid substitution at position 506 (Arg506Gln). Physiologically, activated protein C (APC), in complex with protein S, dampens coagulation by cleaving activated factor V (FVa) at three specific peptide sites, the primary and most rapid of which is Arg506. The Arg506Gln substitution abolishes this initial cleavage site, rendering mutant FVa approximately tenfold slower to inactivate by APC (activated protein C resistance). Consequently, mutant FVa remains active within the prothrombinase complex ($FXa \cdot FVa \cdot Ca^{2+} \cdot phospholipid$), generating excessive thrombin and fostering a hypercoagulable state. Heterozygotes exhibit a 4- to 8-fold increased lifetime risk of venous thromboembolism (VTE, such as deep vein thrombosis and pulmonary embolism), whereas homozygotes face an 80-fold increased risk. The diagnosis is screened via an APC resistance clot-based assay and confirmed definitively by DNA PCR testing. Routine anticoagulation is typically reserved for patients who develop an unprovoked thrombotic event.

\medterm{FAD (Flavin Adenine Dinucleotide)}-- A critical redox-active prosthetic coenzyme composed of riboflavin (vitamin B2) linked to adenosine diphosphate. FAD exists tightly or covalently bound to specific flavoprotein enzymes (such as succinate dehydrogenase in Complex II of the electron transport chain, acyl-CoA dehydrogenase in fatty acid beta-oxidation, and dihydrolipoyl dehydrogenase in the pyruvate dehydrogenase complex). The isoalloxazine ring of FAD acts as a versatile electron acceptor capable of undergoing reversible one- or two-electron reductions, accepting two electrons and two protons ($FAD + 2e^- + 2H^+ \rightleftharpoons FADH_2$) via a stable intermediate semiquinone radical ($FADH^\bullet$). This unique thermodynamic property enables flavoproteins to bridge obligatory two-electron donor reactions (such as carbon-carbon alkane dehydrogenation to alkenes) with one-electron carrier systems (such as iron-sulfur clusters and cytochromes in the mitochondrial respiratory chain). Electrons transferred from $FADH_2$ into the electron transport chain enter downstream at ubiquinone (coenzyme Q), bypassing Complex I and yielding approximately 1.5 ATP molecules per $FADH_2$ oxidized.

\medterm{Fatty Acid}-- A carboxylic acid comprising a terminal hydrophilic carboxyl group (-COOH) linked to an unbranched, hydrophobic aliphatic hydrocarbon tail of varying chain length and saturation, functioning as the premier compact energy reserve of the human body and the structural backbone of membrane lipids (phospholipids and sphingolipids). Fatty acids are classified by: (1) Carbon Chain Length: Short-chain (SCFA, 2--5 carbons, such as acetate, propionate, and butyrate, produced by colonic bacterial fermentation of dietary fiber); Medium-chain (MCFA, 6--12 carbons, absorbed directly into the portal vein without chylomicrons); Long-chain (LCFA, 14--20 carbons, predominantly palmitate [16:0] and oleate [18:1], requiring the carnitine shuttle for mitochondrial entry); and Very long-chain (VLCFA, $\ge 22$ carbons, catabolized initially in peroxisomes); (2) Degree of Saturation: Saturated fatty acids contain no carbon-carbon double bonds, packing tightly into rigid membranes; Monounsaturated fatty acids (MUFAs, e.g., oleic acid 18:1 cis-$\Delta^9$) contain one cis double bond; and Polyunsaturated fatty acids (PUFAs) contain two or more double bonds; and (3) Nutritional Essentiality: Humans cannot introduce double bonds beyond carbon-9 (counting from the carboxyl end); thus, $\alpha$-linolenic acid (an $\omega$-3 fatty acid, 18:3) and linoleic acid (an $\omega$-6 fatty acid, 18:2) are nutritionally essential fatty acids, serving as indispensable precursors for eicosanoids (prostaglandins, thromboxanes, leukotrienes), resolvins, and cell membranes.

\medterm{Fatty Acid Synthase}-- A multifunctional homodimeric enzyme complex ($~540$ kDa) in the cytoplasm that carries out the de novo synthesis of palmitate (a 16-carbon saturated fatty acid) from acetyl-CoA, malonyl-CoA, and NADPH. Each identical monomer contains seven discrete catalytic domains and an acyl carrier protein (ACP) holding a 4'-phosphopantetheine prosthetic group (derived from pantothenic acid, vitamin B5). The reaction begins with an acetyl primer transferred from acetyl-CoA to ACP by acetyl transferase, followed by sequential two-carbon additions donated by malonyl-CoA. Each cycle of elongation involves four coordinated steps: condensation (forming $\beta$-ketoacyl-ACP, releasing $CO_2$), reduction (requiring NADPH, forming $\beta$-hydroxyacyl-ACP), dehydration (forming trans-$\Delta^2$-enoyl-ACP), and a second reduction (requiring NADPH, forming saturated acyl-ACP). After seven successive cycles, the thioesterase domain cleaves palmitate from ACP. NADPH required for fatty acid synthesis is generated by the pentose phosphate pathway (G6PD) and malic enzyme.

\medterm{Feedback Inhibition}-- A fundamental cellular regulatory mechanism (also designated end-product inhibition) in which the final downstream product of a multienzyme metabolic pathway acts as an allosteric inhibitor of an earlier, usually the first committed and rate-limiting, enzyme of that pathway. This negative feedback loop ensures homeostatic metabolic economy, automatically preventing the wasteful overproduction and intracellular accumulation of intermediate metabolites and end-products when physiological requirements are met, while immediately disinhibiting flux when end-product levels decline. The inhibitor binds to a regulatory allosteric site distinct from the catalytic active site, inducing a conformational transition that reduces enzyme affinity for substrate ($K_m$) or decreases catalytic turnover rate ($V_{\max}$). Classic medical biochemistry examples include: (1) Purine and pyrimidine biosynthesis: CTP allosterically inhibits aspartate transcarbamoylase (ATCase) in pyrimidine synthesis; IMP, AMP, and GMP feedback-inhibit PRPP amidotransferase in de novo purine synthesis; (2) Cholesterol biosynthesis: intracellular free cholesterol and oxysterols accelerate the ubiquitination and proteasomal degradation of HMG-CoA reductase while blocking SREBP-2 processing; (3) Heme biosynthesis: free mitochondrial heme directly feedback-inhibits the rate-limiting enzyme $\delta$-aminolevulinic acid synthase-1 (ALAS1) in hepatocytes; (4) Fatty acid synthesis: palmitoyl-CoA allosterically inhibits acetyl-CoA carboxylase; and (5) Glycolysis: ATP and citrate feedback-inhibit phosphofructokinase-1 (PFK-1).

\medterm{Fibrinogen}-- Factor I of the coagulation cascade, a large 340-kDa hexameric plasma glycoprotein ($(\alpha\beta\gamma)_2$) synthesized exclusively by hepatocytes, circulating at high concentrations (200--400 mg/dL). Fibrinogen is the precursor of the definitive fibrin clot and serves as an essential acute-phase reactant whose transcription is strongly upregulated by interleukin-6. In the final common pathway of blood coagulation, thrombin proteolytically cleaves four small, highly negatively charged peptides (fibrinopeptides A and B) from the central E domain of fibrinogen. The unmasked central domain then spontaneously polymerizes with outer D domains of neighboring molecules via non-covalent end-to-middle bonds to form an insoluble protofibril mesh. Factor XIIIa (fibrin-stabilizing factor) subsequently introduces covalent isopeptide bonds between lysine and glutamine residues, converting soft fibrin into a robust, cross-linked hemostatic plug. Fibrinogen also binds platelet integrin $\alpha_{IIb}\beta_3$ (GPIIb/IIIa), mediating platelet aggregation. Acquired hypofibrinogenemia ($<100$ mg/dL) occurs in severe liver failure, massive transfusion dilution, and consumption coagulopathies like disseminated intravascular coagulation (DIC), requiring cryoprecipitate or fibrinogen concentrate replacement.

\medterm{Fluid Imbalance}-- A clinical state of pathological excess, deficit, or abnormal compartmental distribution of total body water and its dissolved electrolyte constituents. Total body water comprises intracellular fluid (ICF, $\approx 2/3$) and extracellular fluid (ECF, $\approx 1/3$, subdivided into interstitial and intravascular plasma spaces), maintained in osmotic equilibrium primarily by sodium ($Na^+$, the dominant extracellular osmole) and potassium ($K^+$, the dominant intracellular osmole). Fluid imbalances are classified by volume and tonicity: (1) Volume Depletion (hypovolemia): loss of sodium and water from gastrointestinal losses (vomiting, diarrhea), renal wasting (diuretics, Addison disease), or third-spacing (burns, pancreatitis), triggering baroreceptor-mediated sympathetic activation, renin-angiotensin-aldosterone system (RAAS) stimulation, and non-osmotic ADH secretion; (2) Volume Overload (hypervolemia): pathological sodium and water retention seen in congestive heart failure, hepatic cirrhosis, and nephrotic syndrome, where reduced effective circulating arterial blood volume (EABV) drives persistent renal sodium retention despite massive interstitial edema; and (3) Osmolar Imbalances: Hypernatremia ($>145$ mEq/L, indicating free water deficit and cellular dehydration) and Hyponatremia ($<135$ mEq/L, indicating relative water excess and risking cerebral edema or central pontine myelinolysis upon rapid correction).

\medterm{Folate (Vitamin B9)}-- A water-soluble B-complex vitamin composed of a pteridine ring, para-aminobenzoic acid (PABA), and one or more L-glutamate residues. Dietary polyglutamate folates are hydrolyzed to monoglutamates by brush-border conjugases and reduced intracellularly by dihydrofolate reductase (DHFR) into the biologically active coenzyme tetrahydrofolate (THF). THF functions as an essential carrier of one-carbon units (methyl, methylene, formyl) in vital biosynthetic pathways: (1) De novo pyrimidine synthesis: conversion of dUMP to dTMP by thymidylate synthase (using $N^5,N^{10}$-methylene-THF); (2) De novo purine synthesis: supplying carbons 2 and 8 of the purine ring; and (3) Amino acid metabolism: remethylation of homocysteine to methionine via methionine synthase (using $N^5$-methyl-THF and vitamin B12). Because the human body stores relatively small reserves of folate in the liver (sufficient for only 3 to 4 months), nutritional deficiency develops rapidly in alcohol use disorder, malabsorption (celiac disease), pregnancy, chronic hemolysis, or during therapy with DHFR inhibitors (methotrexate, trimethoprim) or anticonvulsants (phenytoin). Folate deficiency manifests as megaloblastic anemia (characterized by macrocytic erythrocytes $[MCV > 100\text{ fL}]$ and hypersegmented neutrophils $[>5\text{ lobes}]$ on blood smear) with elevated plasma homocysteine and NORMAL methylmalonic acid (distinguishing it from vitamin B12 deficiency). Crucially, folate deficiency causes NO neurological symptoms. Maternal peri-conceptional folate deficiency causes fetal neural tube defects (spina bifida, anencephaly), prompting mandatory folic acid food fortification.

\medterm{Free Radical}-- Any chemical molecular species, atom, or ion capable of independent existence that possesses one or more unpaired valence electrons in its outer atomic orbitals, a configuration that imparts extreme thermodynamic instability and chemical reactivity. To regain electron pairing, free radicals rapidly abstract electrons or hydrogen atoms from neighboring biological macromolecules (lipids, proteins, nucleic acids), transforming the targeted molecule into a secondary free radical and launching self-propagating, destructive free-radical chain reactions. The most clinically important biological free radicals are reactive oxygen species (ROS) and reactive nitrogen species (RNS), including: the superoxide anion radical ($O_2^{\bullet-}$), generated by electron leakage from the mitochondrial electron transport chain or by NADPH oxidase during phagocytic respiratory burst; the hydroxyl radical ($^{\bullet}OH$), the most highly reactive and damaging radical in human biology, produced from hydrogen peroxide and ferrous iron ($Fe^{2+}$) via the transition-metal-catalyzed Fenton reaction ($Fe^{2+} + H_2O_2 \to Fe^{3+} + ^{\bullet}OH + OH^-$) or the Haber-Weiss reaction; and the nitric oxide radical ($NO^\bullet$), synthesized by nitric oxide synthase, which reacts rapidly with superoxide to generate the powerful, cytotoxic non-radical oxidant peroxynitrite ($ONOO^-$). Free-radical-mediated damage plays a pivotal role in reperfusion injury, radiation-induced cellular necrosis, paracetamol hepatotoxicity (via NAPQI depletion of glutathione), and oncogenic mutagenesis.

\medterm{Fructose Metabolism}-- The specialized hepatic pathway that converts dietary D-fructose (derived from sucrose hydrolysis or high-fructose corn syrup) into glycolytic intermediates. In hepatocytes, fructose metabolism bypasses the primary rate-limiting step of glycolysis (PFK-1), allowing unrestricted, rapid flux into downstream triose phosphates. The pathway involves three sequential enzymes: (1) Fructokinase (ketohexokinase) phosphorylates fructose at carbon-1 using ATP to generate fructose-1-phosphate (F-1-P); (2) Aldolase B (fructose-1,6-bisphosphate aldolase) cleaves F-1-P into dihydroxyacetone phosphate (DHAP) and D-glyceraldehyde; and (3) Triokinase phosphorylates glyceraldehyde using ATP to yield glyceraldehyde-3-phosphate (G3P). Both DHAP and G3P directly enter glycolysis, lipogenesis, or gluconeogenesis. Two major inborn errors exist: Essential Fructosuria, an entirely benign, asymptomatic autosomal recessive condition caused by fructokinase deficiency, resulting in elevated fructose in blood and urine without clinical toxicity; and Hereditary Fructose Intolerance (HFI), a life-threatening autosomal recessive disorder caused by aldolase B deficiency. In HFI, ingestion of fructose, sucrose, or sorbitol traps intracellular phosphate as accumulated F-1-P, causing severe ATP depletion, inhibition of glycogenolysis and gluconeogenesis, profound hypoglycemia, vomiting, jaundice, hepatomegaly, and renal proximal tubular dysfunction.

\medterm{Fructose-1,6-Bisphosphatase}-- The committed, rate-limiting, and key regulatory enzyme of gluconeogenesis, localized to the cytoplasm of hepatocytes and renal tubular cells. It catalyzes the irreversible hydrolytic dephosphorylation of D-fructose-1,6-bisphosphate to D-fructose-6-phosphate and inorganic phosphate ($P_i$), bypassing the irreversible glycolytic step catalyzed by phosphofructokinase-1 (PFK-1). Fructose-1,6-bisphosphatase is subjected to reciprocal allosteric regulation relative to PFK-1 to prevent futile cycling: it is potently allosterically inhibited by AMP (indicating low cellular energy charge) and by fructose-2,6-bisphosphate (F-2,6-BP), while being stimulated by ATP and citrate. Under fasting conditions, elevated glucagon acts through protein kinase A (PKA) to phosphorylate the bifunctional PFK-2/FBPase-2 enzyme, activating FBPase-2 and lowering F-2,6-BP levels; this relieves inhibition of fructose-1,6-bisphosphatase, driving hepatic glucose output. Hereditary fructose-1,6-bisphosphatase deficiency is an autosomal recessive inborn error of metabolism characterized by recurrent episodes of severe fasting hypoglycemia, lactic acidosis, ketosis, and hepatomegaly, triggered by fasting or febrile illness.

\medterm{G6PD Deficiency}-- The most common human enzymopathy worldwide (affecting $>400$ million individuals), an X-linked recessive disorder caused by mutations in the G6PD gene on the X chromosome, conferring partial evolutionary resistance against falciparum malaria. Glucose-6-phosphate dehydrogenase (G6PD) catalyzes the first, rate-limiting step of the pentose phosphate pathway (hexose monophosphate shunt): the oxidation of glucose-6-phosphate to 6-phosphoglucono-$\delta$-lactone with concomitant reduction of $NADP^+$ to NADPH. Because mature erythrocytes lack mitochondria and nuclei, the HMP shunt is their sole source of NADPH. Erythrocytic NADPH is indispensable for glutathione reductase to reduce oxidized glutathione (GSSG) back to reduced glutathione (GSH). GSH, in turn, is used by glutathione peroxidase to detoxify hydrogen peroxide and organic peroxides, protecting hemoglobin and erythrocyte membranes from oxidative stress. Under baseline conditions, patients are typically asymptomatic; however, exposure to oxidative precipitants---including infections, fava beans (favism), and medications (dapsone, primaquine, sulfonamides, rasburicase, nitrofurantoin)---leads to overwhelming oxidative stress. Hemoglobin sulfhydryl groups are oxidized, causing hemoglobin to precipitate into insoluble aggregates called Heinz bodies. As erythrocytes transit splenic sinusoids, splenic macrophages excise these rigid inclusions, generating characteristic 'bite cells' and 'blister cells' seen on peripheral blood smear, triggering acute intravascular and extravascular hemolytic anemia with jaundice, dark hemoglobinuric urine, and elevated indirect bilirubin. Diagnosis is confirmed by G6PD enzyme activity assays, performed several weeks after an acute hemolytic crisis (to avoid false-normal results from surviving young reticulocytes with higher enzyme levels).

\medterm{Galactose Metabolism}-- The biochemical pathway (the Leloir pathway) responsible for the conversion of dietary D-galactose into glucose-1-phosphate for entry into glycolysis or glycogenesis. Galactose is derived primarily from the intestinal brush-border hydrolysis of milk lactose by lactase. Once absorbed into hepatocytes, galactose metabolism proceeds through four coordinated enzymatic steps: (1) Phosphorylation: Galactokinase (GALK) uses ATP to phosphorylate galactose at carbon-1, generating galactose-1-phosphate; (2) Uridylyl Transfer: Galactose-1-phosphate uridyltransferase (GALT) transfers a UMP moiety from UDP-glucose to galactose-1-phosphate, yielding UDP-galactose and glucose-1-phosphate (which enters glycolysis via phosphoglucomutase); (3) Epimerization: UDP-galactose 4-epimerase (GALE) reversibly inverts the stereochemistry at carbon-4 using $NAD^+$, converting UDP-galactose back to UDP-glucose, thus maintaining the cellular pool of UDP-galactose for glycoprotein, glycolipid, and lactose synthesis; and (4) Secondary polyol pathway: When circulating galactose is elevated, non-specific aldose reductase reduces galactose to galactitol (a polyol sugar alcohol that does not cross cellular membranes), causing osmotic swelling, fiber denaturation, and cataract formation in the ocular lens.

\medterm{Galactosemia}-- An inherited autosomal recessive disorder of galactose metabolism, most commonly and severely presenting as Classic Galactosemia (Type 1), caused by homozygous mutations in the GALT gene encoding galactose-1-phosphate uridyltransferase. Dietary lactose (milk sugar) is hydrolyzed in the intestinal brush border to glucose and D-galactose. In the Leloir pathway, galactose is phosphorylated by galactokinase (GALK) to galactose-1-phosphate; GALT then converts galactose-1-phosphate and UDP-glucose into UDP-galactose and glucose-1-phosphate. GALT deficiency causes catastrophic intracellular accumulation of toxic galactose-1-phosphate in liver, brain, and renal tubular cells, alongside accumulation of free galactose, which aldose reductase converts into the toxic sugar alcohol galactitol. Symptoms manifest rapidly in the first days of life upon ingestion of breast milk or dairy formula: infants develop lethargy, severe vomiting, diarrhea, jaundice, hepatosplenomegaly, bleeding diathesis from hepatic failure, renal Fanconi syndrome, and life-threatening neonatal Escherichia coli sepsis. Galactitol accumulation in the lens causes bilateral oil-drop cataracts. Type 2 galactosemia (GALK deficiency) causes a milder presentation characterized solely by early cataracts due to galactitol deposition without hepatic or renal disease. Diagnosis is confirmed by demonstrating non-glucose reducing substances in urine (Clinitest positive, glucose dipstick negative) and erythrocyte GALT enzyme assay. Treatment mandates immediate and lifelong exclusion of lactose and galactose from the diet.

\medterm{Ghrelin}-- A 28-amino-acid orexigenic peptide hormone synthesized and secreted predominantly by neuroendocrine P/D1 cells in the fundus of the stomach (and to a lesser extent by the duodenum and hypothalamus). Ghrelin is unique among human gastrointestinal hormones as the sole circulating peptide that directly stimulates hunger and appetite. To achieve biological activity, ghrelin undergoes post-translational O-octanoylation at its serine-3 residue, catalyzed by the membrane-bound enzyme ghrelin O-acyltransferase (GOAT); this eight-carbon fatty acid modification is obligatory for ghrelin to cross the blood-brain barrier and bind the growth hormone secretagogue receptor-1a (GHSR-1a). In the arcuate nucleus of the hypothalamus, acyl-ghrelin binds GHSR-1a, stimulating neuropeptide Y (NPY) and agouti-related peptide (AgRP) neurons while suppressing pro-opiomelanocortin (POMC) neurons, thereby increasing food intake and promoting gastric motility and acid secretion. Ghrelin secretion rises sharply during fasting and preprandially, signaling hunger, and falls rapidly within one hour of nutrient ingestion. Ghrelin also acts directly on the anterior pituitary gland to stimulate growth hormone secretion. Pathologically, circulating ghrelin levels are massively elevated in Prader-Willi syndrome, contributing to severe hyperphagia and early morbid obesity, whereas levels are permanently suppressed following gastric sleeve or Roux-en-Y bariatric resection of the gastric fundus.

\medterm{Glucagon Signaling}-- A classic catabolic endocrine signaling cascade triggered when the 29-amino-acid peptide hormone glucagon (secreted by $\alpha$-cells of the pancreatic islets of Langerhans during fasting or hypoglycemia) binds to cell-surface glucagon receptors, predominantly in hepatocytes. The glucagon receptor is a class B G-protein coupled receptor (GPCR) that couples to heterotrimeric $G_{\alpha s}$ proteins. Ligand binding stimulates the exchange of GDP for GTP on $G_{\alpha s}$, dissociating the $\alpha$-subunit to activate adenylate cyclase in the plasma membrane. Adenylate cyclase catalyzes the conversion of cytosolic ATP into the second messenger cyclic adenosine monophosphate (cAMP). Elevated cAMP binds cooperatively to the regulatory subunits of protein kinase A (PKA), releasing active catalytic subunits. PKA mediates coordinated phosphorylation events that rapidly stimulate glucose mobilization while halting glucose storage: (1) Glycogenolysis: PKA phosphorylates and activates phosphorylase kinase, which phosphorylates and activates glycogen phosphorylase, while simultaneously phosphorylating and inactivating glycogen synthase; (2) Gluconeogenesis: PKA phosphorylates the bifunctional PFK-2/FBPase-2 enzyme, activating FBPase-2 and depleting fructose-2,6-bisphosphate (F-2,6-BP), which shuts down glycolysis (PFK-1) and stimulates gluconeogenesis (fructose-1,6-bisphosphatase); PKA also enters the nucleus to phosphorylate CREB, inducing transcription of PEPCK and glucose-6-phosphatase; and (3) Fatty Acid Oxidation: PKA phosphorylates and inactivates acetyl-CoA carboxylase, reducing malonyl-CoA and disinhibiting CPT-1 to drive mitochondrial beta-oxidation and ketogenesis.

\medterm{Glucokinase}-- Hexokinase IV, a monomeric enzyme expressed selectively in hepatocytes and pancreatic $\beta$-cells (as well as hypothalamic neurons) that phosphorylates glucose to glucose-6-phosphate. Glucokinase differs fundamentally from ubiquitously expressed hexokinases I--III in several key kinetic and regulatory features: (1) It exhibits a high Michaelis constant ($K_m \approx 5$--10 mM) corresponding to physiological postprandial portal blood glucose levels, displaying positive cooperativity (sigmoidal kinetics); (2) It has a high maximal velocity ($V_{\max}$), enabling the liver to efficiently clear massive glucose loads during the absorptive state; (3) It is NOT inhibited by its reaction product, glucose-6-phosphate; (4) In hepatocytes, its activity is regulated by glucokinase regulatory protein (GKRP), which sequesters glucokinase in the nucleus during fasting (favored by fructose-6-phosphate) and releases it into the cytoplasm when portal glucose rises (favored by fructose-1-phosphate); (5) In pancreatic $\beta$-cells, it acts as the primary glucose sensor controlling ATP generation and insulin secretion. Heterozygous inactivating mutations cause maturity-onset diabetes of the young type 2 (MODY-2), whereas homozygous mutations cause permanent neonatal diabetes mellitus.

\medterm{Gluconeogenesis}-- Metabolic pathway generating glucose from non-carbohydrate precursors (lactate, glycerol, glucogenic amino acids). Occurs primarily in liver (90%) and kidney (10%). Key steps bypassing irreversible glycolysis reactions: pyruvate -> oxaloacetate (pyruvate carboxylase) -> phosphoenolpyruvate (PEP carboxykinase); fructose-1,6-bisphosphate -> fructose-6-phosphate (FRU-1,6-bisphosphatase); glucose-6-phosphate -> glucose (glucose-6-phosphatase). Occurs during fasting, starvation, intense exercise. Hormonal regulation: stimulated by glucagon, cortisol; inhibited by insulin.

\medterm{Glucose-6-Phosphatase}-- A membrane-bound multicomponent enzyme complex situated on the luminal surface of the endoplasmic reticulum (ER), predominantly expressed in the liver and renal cortex (and minimally in intestinal enterocytes). It catalyzes the final terminal step of both gluconeogenesis and glycogenolysis: the hydrolytic dephosphorylation of glucose-6-phosphate to free D-glucose and inorganic phosphate ($P_i$). The functional complex requires the catalytic subunit (G6PC) paired with specific ER translocases: T1 (transporting G6P into the ER lumen), T2 (transporting $P_i$ out), and T3 (transporting free glucose into the cytosol for release into blood via GLUT2). Skeletal muscle lacks this enzyme, which prevents muscle glycogen from directly contributing to systemic blood glucose. Deficiencies in the catalytic subunit cause Glycogen Storage Disease Type Ia (Von Gierke disease), characterized by severe fasting hypoglycemia, massive hepatomegaly, lactic acidosis, hyperuricemia (gout), and hyperlipidemia; deficiency in the T1 translocase causes Type Ib, which uniquely presents with neutropenia and recurrent pyogenic infections.

\medterm{Glycogen}-- Branched polymer of glucose residues stored in liver and muscle. Alpha-1,4 linkages in chains, alpha-1,6 linkages at branch points (every 8-12 residues). Glycogen synthesis (glycogenesis) stimulated by insulin; breakdown (glycogenolysis) stimulated by glucagon and epinephrine. Glycogen phosphorylase cleaves alpha-1,4 bonds; debranching enzyme handles alpha-1,6 bonds. Glycogenolysis provides glucose during fasting (liver) or rapid energy (muscle).

\medterm{Glycogen Storage Disease}-- A group of inherited autosomal recessive metabolic disorders caused by enzymatic deficiencies in the synthesis or catabolism of glycogen, resulting in abnormal glycogen accumulation and cellular damage in the liver, skeletal muscle, or cardiac tissue. Key clinical types include: (1) Type I (Von Gierke disease): deficiency of glucose-6-phosphatase (Type Ia) or G6P translocase (Type Ib), impairing both glycogenolysis and gluconeogenesis, causing profound fasting hypoglycemia, severe lactic acidosis, hyperuricemia (gout), hyperlipidemia, and massive hepatomegaly (normal heart); (2) Type II (Pompe disease): deficiency of lysosomal $\alpha$-(1$\to$4)-glucosidase (acid maltase), leading to massive glycogen buildup within lysosomes, causing hypertrophic cardiomyopathy, generalized hypotonia ('floppy infant'), macroglossia, and early death from heart failure; (3) Type III (Cori or Forbes disease): deficiency of glycogen debranching enzyme ($\alpha$-1,6-glucosidase and 4-$\alpha$-glucanotransferase), accumulating limit dextrin with short outer branches, causing mild fasting hypoglycemia and hepatomegaly with normal blood lactate levels (gluconeogenesis is intact); (4) Type IV (Andersen disease): deficiency of branching enzyme, accumulating poorly branched, insoluble polyglucosan that provokes fatal hepatic cirrhosis; and (5) Type V (McArdle disease): deficiency of skeletal muscle glycogen phosphorylase (myophosphorylase), causing painful muscle cramps, exercise intolerance, and myoglobinuria/rhabdomyolysis following strenuous exertion, with a characteristic flat venous lactate curve during ischemic exercise.

\medterm{Glycogenesis}-- The anabolic biochemical pathway of glycogen synthesis from $\alpha$-D-glucose, occurring primarily in the cytoplasm of hepatocytes (supplying systemic blood glucose during fasting) and skeletal myocytes (providing an immediate fuel reservoir for muscle contraction). The pathway is powerfully induced in the absorptive state by insulin and high intracellular glucose-6-phosphate (G6P). The process begins with glucose phosphorylation by hexokinase or glucokinase to G6P, followed by phosphoglucomutase-mediated isomerization to glucose-1-phosphate (G1P). G1P then reacts with UTP, catalyzed by UDP-glucose pyrophosphorylase, to generate high-energy uridine diphosphate glucose (UDP-glucose). De novo glycogen chain initiation requires the self-glucosylating primer protein glycogenin, which attaches eight glucose residues to a specific tyrosine residue ($Tyr194$). Subsequently, the committed, rate-limiting enzyme glycogen synthase transfers glucosyl units from UDP-glucose to the non-reducing ends of the growing glycogen polymer via $\alpha$-(1$\to$4)-glycosidic linkages. Glycogen synthase is allosterically activated by G6P and covalently regulated: insulin stimulates protein phosphatase 1 (PP1), which dephosphorylates and activates glycogen synthase; in contrast, glucagon, epinephrine, and PKA phosphorylate and inactivate it. To increase solubility and multiply terminal reaction sites, the branching enzyme (amylo-$\alpha$(1$\to$4)$\to$$\alpha$(1$\to$6)-transglucosidase) removes an outer 7-residue oligosaccharide segment and transfers it to an interior glucose via an $\alpha$-(1$\to$6)-glycosidic bond every 8 to 12 residues. Glycogen branching enzyme deficiency causes Type IV Glycogen Storage Disease (Andersen disease).

\medterm{Glycogenolysis}-- Breakdown of glycogen to glucose-1-phosphate (converted to glucose-6-phosphate). Liver glycogenolysis maintains blood glucose during fasting; muscle glycogenolysis provides energy for contraction. Regulated by hormonal signals (glucagon, epinephrine) and allosteric effectors (AMP activates phosphorylase). Deficiencies cause glycogen storage diseases (Von Gierke, Pompe, McArdle).

\medterm{Glycolysis}-- Metabolic pathway breaking down glucose (6C) to pyruvate (3C), generating 2 ATP (net) and 2 NADH. Occurs in cytoplasm. Ten enzymatic steps: investment phase (uses 2 ATP) and payoff phase (produces 4 ATP, 2 NADH). Key regulatory enzymes: hexokinase, phosphofructokinase-1 (PFK-1, rate-limiting), pyruvate kinase. Fate of pyruvate: aerobic (TCA cycle), anaerobic (lactate), or converted to acetyl-CoA (link reaction).

\medterm{Growth Hormone}-- Somatotropin, a 191-amino-acid, 22-kDa single-chain polypeptide hormone synthesized, stored, and secreted by somatotroph cells in the anterior pituitary gland in a pulsatile pattern (with peak surges occurring during slow-wave sleep). Secretion is stimulated by hypothalamic growth hormone-releasing hormone (GHRH) and ghrelin, and inhibited by somatostatin and IGF-1 negative feedback. GH binds to cell-surface dimeric cytokine receptors, recruiting and phosphorylating JAK2, which phosphorylates STAT5 to regulate target gene transcription. GH exerts dual metabolic and somatic actions: (1) Direct Metabolic Effects: GH acts as a powerful counter-regulatory, anti-insulin, and catabolic hormone that raises blood glucose and preserves protein during fasting. In adipose tissue, it stimulates hormone-sensitive lipase to promote lipolysis and elevate free fatty acids; in skeletal muscle, it inhibits glucose uptake by impairing insulin receptor substrate-1 (IRS-1) signaling, producing peripheral insulin resistance; in the liver, it stimulates gluconeogenesis; and (2) Indirect Anabolic Effects: In hepatocytes, GH stimulates the synthesis and secretion of insulin-like growth factor-1 (IGF-1 / somatomedin C). IGF-1 acts on epiphyseal growth plates in children to stimulate chondrocyte proliferation, linear bone growth, and lean muscle mass accretion. GH deficiency in childhood causes pituitary dwarfism / growth failure; in adults, it causes reduced muscle mass and osteopenia. GH-secreting pituitary adenomas cause Gigantism in children (prior to epiphyseal fusion) and Acromegaly in adults (frontal bossing, enlarged hands/feet, prognathism, macroglossia, cardiomegaly, and secondary diabetes mellitus).

\medterm{HbA1c (Glycated Hemoglobin)}-- A stable minor fraction of adult hemoglobin $A$ formed through the slow, non-enzymatic, irreversible covalent condensation of D-glucose with the N-terminal valine amino group of the hemoglobin $\beta$-polypeptide chains (glycation / Maillard reaction). The reaction proceeds through a reversible intermediate aldimine (Schiff base) that undergoes an Amadori rearrangement over several weeks to form a permanent, stable ketoamine adduct ($HbA_{1c}$). Because red blood cells are freely permeable to glucose via insulin-independent GLUT1 transporters and possess a mean circulatory lifespan of approximately 120 days, the proportion of circulating hemoglobin that is glycated directly reflects average weighted plasma glucose exposure over the preceding 8 to 12 weeks (with the preceding 30 days contributing ~50\% of the value). Normal HbA1c is $<5.7\%$; prediabetes is defined as $5.7--6.4\%$; and diabetes mellitus is diagnosed at $\ge 6.5\%$ (confirmed on repeat testing). HbA1c is the gold-standard clinical metric for monitoring long-term glycemic control and predicting risk of microvascular complications (diabetic retinopathy, nephropathy, and neuropathy). Results can be falsely lowered by conditions with shortened red cell survival (hemolytic anemias, acute blood loss, hemodialysis, pregnancy, high-dose erythropoietin) and falsely elevated in states of prolonged erythrocyte lifespan (iron deficiency anemia, vitamin B12/folate deficiency, asplenia).

\medterm{HDL}-- High-density lipoprotein, the smallest (5--12 nm) and densest (1.063--1.210 g/mL) class of plasma lipoproteins, synthesized and secreted as lipid-poor discoidal particles by hepatocytes and intestinal enterocytes. HDL is composed of a protein-rich shell (approximately 50\% protein by mass) dominated by apolipoprotein A-I (ApoA-I) and ApoA-II, and serves as the molecular centerpiece of Reverse Cholesterol Transport (RCT)---the physiological process of extracting excess cholesterol from peripheral tissues and macrophages and returning it to the liver for biliary excretion. Nascent pre-$\beta$-HDL acquires free unesterified cholesterol and phospholipids from peripheral cell membranes via the ATP-binding cassette transporter ABCA1. The surface enzyme lecithin-cholesterol acyltransferase (LCAT), activated by ApoA-I, transfers a fatty acyl chain from phosphatidylcholine to free cholesterol, forming hydrophobic cholesteryl esters that migrate into the particle's core, transforming discoidal HDL into mature spherical $\alpha$-HDL ($HDL_3$ and $HDL_2$). HDL delivers cholesteryl esters to the liver via scavenger receptor class B type I (SR-BI) without particle degradation, or transfers them to VLDL and LDL in exchange for triglycerides via cholesteryl ester transfer protein (CETP). HDL also acts as a circulating reservoir of ApoC-II and ApoE for chylomicrons and VLDL, and exhibits anti-inflammatory, antioxidant (via paraoxonase), and anti-thrombotic properties. Epidemiologically, serum HDL-C ($>40$ mg/dL in men, $>50$ mg/dL in women) is strongly and inversely correlated with atherosclerotic cardiovascular disease.

\medterm{Hemochromatosis}-- An autosomal recessive disorder of systemic iron overload caused predominantly by homozygous missense mutations (C282Y or H63D) in the HFE gene on chromosome 6, which modulates the body's iron sensor mechanism. Under physiological conditions, the HFE protein interacts with transferrin receptor-1 and transferrin receptor-2 on hepatocytes to upregulate the synthesis of hepcidin, the master iron regulatory hormone. Mutant HFE fails to signal, resulting in inappropriately low or absent circulating hepcidin. Unchecked, the iron export channel ferroportin remains constitutively active on basolateral membranes of duodenal enterocytes and reticuloendothelial macrophages, driving uncontrolled intestinal iron absorption (2--4 mg/day above normal). Because the human body lacks an active excretory pathway for iron, excess iron saturates transferrin and deposits in parenchymal cells of the liver, pancreas, heart, joints, and pituitary gland as toxic hemosiderin. Free iron catalyzes the generation of reactive oxygen species via the Fenton reaction, inducing lipid peroxidation and progressive tissue fibrosis. Symptoms typically manifest after age 40 (earlier in men due to lack of menstrual blood loss), presenting with the classic clinical triad of micronodular cirrhosis, 'bronze' hyperpigmentation of the skin (melanin and iron deposition), and diabetes mellitus ('bronze diabetes' from pancreatic islet destruction). Other features include dilated or restrictive cardiomyopathy, hypogonadism (pituitary gonadotropin failure), and pseudogout (calcium pyrophosphate deposition in joints). Diagnostic screening shows elevated transferrin saturation ($>45--50\%$) and markedly elevated serum ferritin ($>1,000$ ng/mL), confirmed by HFE genetic testing. Therapeutic phlebotomy (weekly until ferritin $<50$ ng/mL) prevents end-stage organ damage and restores normal life expectancy.

\medterm{Hemoglobin}-- Tetrameric protein in red blood cells (4 subunits: 2 alpha, 2 beta in adults) containing heme groups that bind oxygen. Each heme contains iron (Fe2+) that reversibly binds O2. Cooperative binding produces sigmoidal oxygen dissociation curve. P50 (pO2 at 50% saturation) ~27 mmHg. Factors shifting curve: Bohr effect (H+, CO2 decrease affinity), 2,3-BPG (decreases affinity), temperature (increases with heat). Fetal hemoglobin (HbF: 2 alpha, 2 gamma) has higher O2 affinity. Mutations cause hemoglobinopathies (sickle cell, thalassemia).

\medterm{Hemolytic Anemia}-- A diverse class of anemias characterized by the premature destruction and shortened lifespan of circulating red blood cells ($<120$ days), occurring whenever erythrocytic breakdown outpaces the compensatory 6- to 8-fold proliferative capacity of the bone marrow. Mechanistically classified as: (1) Intrinsic (intracorpuscular) defects, which are predominantly inherited and include hemoglobinopathies (sickle cell anemia, thalassemia), membrane cytoskeletal defects (hereditary spherocytosis, hereditary elliptocytosis), and enzyme deficiencies (G6PD deficiency, pyruvate kinase deficiency); and (2) Extrinsic (extracorpuscular) defects, which are predominantly acquired and include autoimmune hemolytic anemia (warm IgG or cold IgM autoantibodies, diagnosed by positive direct Coombs / DAT testing), microangiopathic hemolytic anemia (MAHA: TTP, HUS, DIC, showing fragmented schistocytes on blood smear), infections (malaria, babesiosis), and mechanical trauma (prosthetic heart valves). Hemolysis is further divided anatomically into extravascular hemolysis (phagocytosis of abnormal RBCs by macrophages in the spleen and liver) and intravascular hemolysis (direct complement-mediated or mechanical rupture within blood vessels). Universal laboratory hallmarks include elevated indirect (unconjugated) bilirubin, elevated serum lactate dehydrogenase (LDH, released from ruptured RBCs), reticulocytosis ($>2--3\%$, with polychromasia on smear), and severely depressed or absent serum haptoglobin (which binds free hemoglobin; intravascular hemolysis also produces hemoglobinemia, hemoglobinuria, and urine hemosiderin).

\medterm{Hemophilia A}-- An X-linked recessive bleeding disorder caused by pathogenic mutations or inversions (most commonly the intron 22 inversion, accounting for 45--50\% of severe cases) in the F8 gene, leading to quantitative deficiency or dysfunction of coagulation factor VIII. Factor VIII functions as an essential regulatory cofactor that complexes with activated factor IX (FIXa), calcium, and phospholipid membranes to form the intrinsic tenase complex, which accelerates the activation of factor X by $>100,000$-fold. Disease severity correlates directly with residual plasma factor VIII activity: severe ($<1\%$ activity, presenting with spontaneous, crippling hemarthroses, deep muscle hematomas, and retroperitoneal bleeding), moderate (1--5\% activity, bleeding after minor trauma), and mild (5--40\% activity, hemorrhage only after major surgery or dental extraction). Laboratory testing demonstrates an isolated, prolonged aPTT that corrects upon a 1:1 mixing study with normal pooled plasma (unless an inhibitory factor VIII alloantibody has developed), while PT, bleeding time, platelet count, and vWF are normal. Management consists of prophylactic or episodic infusion of recombinant or plasma-derived factor VIII concentrates, desmopressin (DDAVP) for mild cases, and the bispecific monoclonal antibody emicizumab (which bridges factor IXa and factor X, mimicking factor VIII function).

\medterm{Hemophilia B}-- An X-linked recessive coagulopathy (also historically termed Christmas disease) caused by mutations in the F9 gene resulting in deficiency of coagulation factor IX, a vitamin K-dependent serine protease zymogen. Activated factor IX (FIXa) serves as the enzymatic engine of the intrinsic tenase complex, partnering with cofactor VIIIa to proteolytically cleave and activate factor X on phospholipid membranes. Although clinically indistinguishable from hemophilia A (manifesting with spontaneous hemarthrosis, muscle hematomas, and life-threatening post-traumatic hemorrhage depending on factor levels: severe $<1\%$, moderate 1--5\%, mild 5--40\%), hemophilia B is five times less common. Laboratory findings mirror hemophilia A, showing an isolated prolonged aPTT with normal PT, platelet count, and fibrinogen, which fully corrects on mixing studies. Specific factor IX functional activity assay confirms the diagnosis. Treatment requires intravenous replacement with recombinant or plasma-derived purified factor IX concentrates. Unlike hemophilia A, desmopressin (DDAVP) is entirely ineffective in hemophilia B, as desmopressin mobilizes vWF and factor VIII from endothelial stores but has no influence on hepatic factor IX synthesis.

\medterm{Hexokinase}-- The enzyme catalyzing the initial, irreversible step of intracellular glucose utilization: the ATP-dependent phosphorylation of D-glucose to glucose-6-phosphate (G6P), effectively trapping glucose within the cytoplasm. Four distinct mammalian isozymes exist (Types I--IV). Hexokinases I--III are ubiquitously expressed across extrahepatic tissues (notably brain, erythrocytes, and muscle). They possess a very low Michaelis constant ($K_m < 0.1$ mM) for glucose, ensuring maximal catalytic activity and glucose uptake even during profound fasting or hypoglycemia, but display a low maximal velocity ($V_{\max}$) to prevent intracellular substrate overaccumulation. Hexokinases I--III are potently allosterically inhibited by their product, G6P (feedback inhibition), preventing unnecessary phosphorylation when downstream glycolysis or glycogenesis is saturated. In contrast, Hexokinase IV (glucokinase), expressed in hepatocytes and pancreatic $\beta$-cells, has a high $K_m$ and is not inhibited by G6P.

\medterm{HMG-CoA Reductase}-- The rate-limiting and committed enzyme of de novo cholesterol biosynthesis, located in the endoplasmic reticulum membrane with its catalytic domain facing the cytosol. It catalyzes the four-electron, two-step reduction of (3S)-hydroxy-3-methylglutaryl-CoA (HMG-CoA) to mevalonate, utilizing two molecules of NADPH and releasing coenzyme A. The enzyme is subjected to intricate multivalent regulation: (1) Transcriptional control: low intracellular sterol levels permit the SREBP-2/SCAP complex to migrate from the ER to the Golgi apparatus, where site-1 and site-2 proteases cleave SREBP-2 to enter the nucleus and upregulate HMG-CoA reductase transcription; (2) Sterol-accelerated degradation: elevated cholesterol and oxysterols bind Insig proteins, promoting ubiquitination and proteasomal degradation of the enzyme; (3) Covalent modification: AMP-activated protein kinase (AMPK) phosphorylates and inactivates the enzyme during energy deficit, while insulin promotes dephosphorylation and activation. Clinically, HMG-CoA reductase is the direct molecular target of statin medications (e.g., atorvastatin, rosuvastatin), which act as reversible competitive inhibitors, dramatically lowering intracellular cholesterol, upregulating hepatic LDL receptors, and reducing circulating atherogenic LDL particles.

\medterm{Hormone}-- Chemical messenger secreted by endocrine glands into blood, acting on distant target organs with specific receptors. Classes: peptides/proteins (insulin, GH), steroids (cortisol, estrogen), amines (thyroid hormone, catecholamines). Mechanisms: cell surface receptors (GPCRs, RTKs) or intracellular receptors (nuclear receptors). Feedback regulation: negative feedback predominates.

\medterm{Hormone-Sensitive Lipase}-- A neutral intracellular lipase (HSL) predominantly expressed in white and brown adipose tissue (and to a lesser degree in skeletal muscle, heart, and steroidogenic tissues) that catalyzes the rate-limiting hydrolytic removal of fatty acids from diacylglycerols (and to a lesser extent triacylglycerols and monoacylglycerols) during adipocyte lipolysis. In the basal, unphosphorylated state, HSL resides in the adipocyte cytosol, separated from lipid droplets coated by the protective structural protein perilipin-1. Under catabolic conditions, catecholamines (epinephrine, norepinephrine) bind $\beta$-adrenergic receptors, increasing cAMP and activating protein kinase A (PKA). PKA dual-phosphorylates both HSL (promoting its translocation to the lipid droplet) and perilipin-1 (releasing CGI-58 to activate adipose triglyceride lipase, ATGL). Phosphorylated HSL then hydrolyzes diacylglycerol to monoacylglycerol and free fatty acid. In contrast, insulin potently halts lipolysis by activating phosphodiesterase-3B (PDE-3B) and protein phosphatases, decreasing cAMP, dephosphorylating HSL, and promoting lipid storage.

\medterm{Hydrogen Bond}-- Weak electrostatic interaction between a hydrogen atom covalently bonded to an electronegative atom (N, O, F) and another electronegative atom. Critical for protein structure (alpha helix, beta sheet), DNA double helix (base pairing), enzyme-substrate binding, and water properties (high boiling point, solvent properties). Bond energy: ~4-5 kcal/mol.

\medterm{Hyperammonemia}-- A life-threatening clinical and biochemical emergency characterized by an abnormally elevated concentration of neurotoxic ammonia in arterial or venous blood ($>50 \ \mu\text{mol/L}$ in adults, $>100 \ \mu\text{mol/L}$ in neonates). Ammonia ($NH_3/NH_4^+$) is generated endogenously from the catabolism of dietary and cellular proteins (transamination and oxidative deamination of amino acids) and by urease-producing gut microbes. Under normal physiology, the liver detoxifies ammonia into non-toxic, water-soluble urea via the five enzymes of the urea cycle. Etiologies are divided into: (1) Primary Inborn Errors of Metabolism: congenital defects in urea cycle enzymes, most commonly X-linked Ornithine Transcarbamylase (OTC) deficiency (manifesting with hyperammonemia, low citrulline, and high urinary orotic acid), or autosomal recessive carbamoyl phosphate synthetase I (CPS-I) deficiency; and (2) Acquired Hyperammonemia: severe acute or chronic hepatic failure (cirrhosis, portosystemic shunting bypassing hepatocytes, Reye syndrome, or valproic acid toxicity). Excess circulating ammonia readily crosses the blood-brain barrier. In cerebral astrocytes, ammonia is detoxified by glutamine synthetase, condensing ammonia with $\alpha$-ketoglutarate-derived glutamate to form glutamine. Intracellular accumulation of osmotic glutamine causes astrocytic swelling, cerebral edema, disruption of the blood-brain barrier, and intracranial hypertension. Depletion of $\alpha$-ketoglutarate also cripples the mitochondrial citric acid cycle in neurons, arresting ATP production. Patients present with asterixis ('flapping' tremor), vomiting, cerebral edema, lethargy, seizures, coma, and death. Management involves dietary protein restriction, osmotic laxatives (lactulose, which acidifies colonic contents to trap $NH_4^+$), ammonia-scavenging nitrogen-binding agents (sodium phenylbutyrate, sodium benzoate), and emergency hemodialysis.

\medterm{Inflammation}-- A complex, highly orchestrated physiological defense response of vascularized tissue to harmful stimuli, including microbial infection, tissue necrosis, physical trauma, or toxic chemicals, aimed at neutralizing the offending agent and initiating tissue repair. The biochemical cascade is initiated by resident sentinel cells (macrophages, dendritic cells, mast cells) whose pattern recognition receptors (PRRs, such as Toll-like receptors) detect pathogen-associated molecular patterns (PAMPs, e.g., LPS) or damage-associated molecular patterns (DAMPs, e.g., ATP, HMGB1). Receptor engagement activates NF-$\kappa$B and the NLRP3 inflammasome, stimulating the release of primary inflammatory mediators: vasoactive amines (histamine, serotonin causing immediate vasodilation and increased venular permeability), lipid autacoids (prostaglandins PGE2/PGI2 mediating vasodilation and hyperalgesia; leukotrienes LTB4 recruiting neutrophils, LTC4/LTD4/LTE4 inducing bronchoconstriction and vascular permeability), and cytokines (TNF-$\alpha$, IL-1$\beta$, IL-6). These mediators induce endothelial adhesion molecule expression (E-selectin, ICAM-1, VCAM-1), orchestrating neutrophil and monocyte margination, rolling, tight adhesion, and diapedesis into the extravascular space. The classic cardinal signs (rubor, calor, tumor, dolor, and functio laesa) reflect these microvascular hemodynamic alterations. Successful acute inflammation resolves through active biochemical class switching to pro-resolving lipid mediators (lipoxins, resolvins, protectins); failure to clear the stimulus results in chronic inflammation, characterized by tissue destruction, macrophage/lymphocyte infiltration, neoangiogenesis, and progressive fibrosis.

\medterm{Insulin}-- Peptide hormone (51 amino acids) produced by beta cells of pancreatic islets. Secreted in response to hyperglycemia. Actions: promotes glucose uptake (GLUT4 translocation in muscle/adipose), glycogen synthesis (activates glycogen synthase), lipogenesis, protein synthesis; inhibits gluconeogenesis, glycogenolysis, lipolysis, proteolysis. Receptor: tyrosine kinase. Deficiency or resistance causes diabetes mellitus. Analogues: rapid-acting (lispro, aspart), long-acting (glargine, detemir).

\medterm{Insulin Signaling}-- A fundamental anabolic signal transduction cascade initiated when insulin binds to the extracellular $\alpha$-subunits of the cell-surface insulin receptor, a preformed heterotetrameric receptor tyrosine kinase ($\alpha_2\beta_2$). Ligand binding induces trans-autophosphorylation of specific tyrosine residues on the intracellular catalytic $\beta$-subunits, creating docking sites that recruit and phosphorylate insulin receptor substrate proteins (IRS-1 through IRS-4). Phosphorylated IRS proteins diverge into two primary downstream signaling arms: (1) The PI3K-Akt Pathway (metabolic arm): Phosphorylated IRS recruits phosphoinositide 3-kinase (PI3K), which converts $PIP_2$ to $PIP_3$ at the inner leaflet of the plasma membrane. $PIP_3$ recruits phosphoinositide-dependent kinase-1 (PDK1) and protein kinase B (Akt). Activated Akt phosphorylates AS160 (Akt substrate of 160 kDa), triggering the exocytic translocation of glucose transporter-4 (GLUT4) storage vesicles to the plasma membrane in skeletal muscle and adipose tissue, driving rapid glucose uptake. Akt also phosphorylates and inactivates glycogen synthase kinase-3 (GSK-3), disinhibiting glycogen synthase to promote glycogenesis, and activates mTORC1 to stimulate protein translation while suppressing FoxO1-dependent gluconeogenic gene expression; and (2) The MAPK/ERK Pathway (mitogenic arm): IRS recruits Grb2-SOS, activating the Ras-Raf-MEK-ERK kinase cascade to stimulate cellular proliferation and gene transcription. Pathologically, serine/threonine phosphorylation of IRS-1 by inflammatory kinases (JNK, IKK) or elevated free fatty acids inhibits insulin signaling, driving clinical insulin resistance and type 2 diabetes mellitus.

\medterm{Insulin-like Growth Factor}-- A family of structural and functional peptide hormones comprising IGF-1 (somatomedin C, a 70-amino-acid peptide) and IGF-2 (a 67-amino-acid peptide) that share high sequence homology with proinsulin. IGF-1 is synthesized predominantly by hepatocytes under the endocrine control of growth hormone (GH), and locally by peripheral mesenchymal tissues where it acts via autocrine/paracrine mechanisms. In circulation, $>98\%$ of IGF-1 is bound to a family of high-affinity IGF-binding proteins (primarily IGFBP-3) and an acid-labile subunit (ALS), forming a 150-kDa ternary complex that extends its biological half-life from minutes to over 15 hours. IGF-1 binds to the IGF-1 receptor (IGF-1R), a transmembrane receptor tyrosine kinase closely homologous to the insulin receptor. Receptor activation stimulates both the PI3K-Akt and MAPK/ERK signaling cascades, driving potent anabolic actions: stimulating cellular proliferation, inhibiting apoptosis, promoting amino acid uptake, and accelerating protein synthesis. In the skeletal system, IGF-1 acts on chondrocytes at epiphyseal growth plates to stimulate clonal expansion and matrix secretion, driving longitudinal bone growth during childhood and puberty. Because serum GH fluctuates wildly in pulsatile bursts, measuring serum IGF-1 provides a stable, integrated marker used clinically to diagnose and monitor acromegaly, gigantism, and growth hormone deficiency.

\medterm{Interleukin-6}-- A multifunctional, 21--28 kDa pleiotropic cytokine synthesized and secreted by monocytes, macrophages, vascular endothelial cells, fibroblasts, and contracting skeletal myocytes during infection, tissue trauma, or vigorous physical exertion. IL-6 binds to either a membrane-bound receptor (IL-6R / CD126, present on hepatocytes and leukocytes) or soluble IL-6R (sIL-6R, mediating 'trans-signaling' across all cell types), triggering homodimerization of the ubiquitous signal-transducing $\beta$-subunit glycoprotein 130 (gp130) and activating the intracellular JAK-STAT (primarily STAT3) and MAPK cascades. In systemic inflammation, IL-6 functions as the primary endocrine inducer of the hepatic acute-phase response: it transcriptionally upregulates hepatic production of C-reactive protein (CRP), fibrinogen, serum amyloid A, and hepcidin (which induces anemia of chronic disease), while downregulating negative acute-phase reactants such as albumin and transferrin. IL-6 also drives the differentiation of naive CD4+ T-cells into pro-inflammatory Th17 cells, promotes B-cell maturation into immunoglobulin-secreting plasma cells, and stimulates osteoclastogenesis via RANKL. Unrestrained hypersecretion of IL-6 drives the life-threatening 'cytokine release syndrome' (cytokine storm) observed in severe COVID-19 and CAR-T cell therapy, managed clinically with targeted IL-6 receptor antagonists (tocilizumab, sarilumab).

\medterm{Iron Deficiency Anemia}-- The most common nutritional deficiency and form of anemia worldwide, caused by inadequate systemic iron stores to support normal erythropoiesis. Etiologies include chronic occult gastrointestinal blood loss (peptic ulcers, colorectal carcinoma, angiodysplasia), menorrhagia in premenopausal women, malabsorption (celiac disease, post-gastrectomy), or increased physiological demand (pregnancy, childhood growth). Iron deficiency progresses through three discrete phases: (1) Storage iron depletion: serum ferritin drops ($<15--30$ ng/mL) while hemoglobin, serum iron, and erythrocyte morphology remain normal; (2) Iron-deficient erythropoiesis: serum iron falls, total iron-binding capacity (TIBC) rises ($>400$ mcg/dL), transferrin saturation plummets ($<15\%$), and erythrocyte zinc protoporphyrin accumulates; and (3) Overt iron deficiency anemia: hemoglobin synthesis is crippled, resulting in microcytic, hypochromic erythrocytes ($MCV < 80$ fL, $MCH < 27$ pg) with high red cell distribution width (RDW, reflecting anisocytosis). Peripheral blood smear shows pencil cells, target cells, and central pallor $>1/3$ the erythrocyte diameter. Clinical features include fatigue, pallor, pica (compulsive consumption of ice, dirt, or starch), pagophagia, koilonychia (spoon-shaped nails), atrophic glossitis, angular cheilitis, and Plummer-Vinson syndrome (triad of microcytic anemia, atrophic glossitis, and esophageal webs). Treatment involves oral ferrous sulfate (enhanced by ascorbic acid) or intravenous iron formulations, alongside mandatory endoscopic evaluation in adults to exclude gastrointestinal malignancy.

\medterm{Isocitrate Dehydrogenase}-- The rate-limiting, committed enzyme of the citric acid cycle (Krebs cycle), localized to the mitochondrial matrix (NAD+-dependent IDH3). It catalyzes the oxidative decarboxylation of D-isocitrate to $\alpha$-ketoglutarate, generating NADH and liberating the first molecule of $CO_2$ from the cycle. The reaction proceeds through an unstable enzyme-bound intermediate, oxalosuccinate, which undergoes spontaneous $\beta$-decarboxylation stabilized by a divalent metal cofactor ($Mg^{2+}$ or $Mn^{2+}$). IDH3 is allosterically activated by ADP (which lowers the $K_m$ for isocitrate) and $Ca^{2+}$, signaling increased muscle contraction and metabolic demand, while it is allosterically inhibited by ATP and NADH, reflecting high cellular energy charge. Cytosolic and mitochondrial NADP+-dependent isoforms (IDH1 and IDH2) participate in NADPH generation and cellular redox balance; somatic missense mutations in IDH1 (e.g., R132H) and IDH2 produce the oncometabolite (2R)-2-hydroxyglutarate (2-HG), driving epigenetic hypermethylation and oncogenesis in gliomas and acute myeloid leukemia.

\medterm{Ketone Bodies}-- Three water-soluble, four-carbon (and three-carbon) organic molecules---acetoacetate, $\beta$-hydroxybutyrate, and acetone---synthesized exclusively by hepatocytes in the mitochondrial matrix from fatty-acid-derived acetyl-CoA during states of low carbohydrate availability, prolonged fasting, starvation, or insulin deficiency. Ketogenesis is initiated when intense beta-oxidation floods the mitochondrial matrix with acetyl-CoA while oxaloacetate is actively drained by gluconeogenesis. The pathway involves four sequential steps: (1) Thiolase condenses two acetyl-CoA molecules to acetoacetyl-CoA; (2) Mitochondrial HMG-CoA synthase (the rate-limiting enzyme of ketogenesis) condenses a third acetyl-CoA to form 3-hydroxy-3-methylglutaryl-CoA (HMG-CoA); (3) HMG-CoA lyase cleaves HMG-CoA to release free acetoacetate and acetyl-CoA; and (4) Mitochondrial $\beta$-hydroxybutyrate dehydrogenase interconverts acetoacetate and $\beta$-hydroxybutyrate in an $NADH/NAD^+$-dependent reversible reaction (with $\beta$-hydroxybutyrate predominating in diabetic ketoacidosis due to high mitochondrial NADH). Spontaneous non-enzymatic decarboxylation of acetoacetate yields volatile acetone, which is excreted via the lungs, imparting a characteristic fruity or sweet odor to the breath. Ketone bodies are released into the systemic circulation and utilized as vital water-soluble fuels by extrahepatic tissues (brain, myocardium, skeletal muscle, renal cortex), where they are converted back to acetyl-CoA via thiophorase ($\beta$-ketoacyl-CoA transferase). Crucially, hepatocytes lack thiophorase, preventing the liver from futilely consuming the ketones it produces.

\medterm{Ketosis}-- A metabolic state characterized by elevated concentrations of circulating ketone bodies (acetoacetate, $\beta$-hydroxybutyrate, and acetone) in the bloodstream (ketonemia $>0.5$ mmol/L) and their appearance in urine (ketonuria). Ketone bodies are synthesized exclusively in hepatocytes from fatty-acid-derived acetyl-CoA when intracellular carbohydrate availability or insulin signaling is low. In hepatocytes, prolonged fasting, starvation, or intense exercise shifts the hormonal milieu toward a low insulin-to-glucagon ratio, massively stimulating adipocyte lipolysis via hormone-sensitive lipase. A huge flux of free fatty acids enters the liver and undergoes mitochondrial beta-oxidation, flooding the matrix with acetyl-CoA. Concurrently, gluconeogenesis actively consumes oxaloacetate (converting it to malate for export to the cytoplasm), which depletes the oxaloacetate required for citrate synthase to condense with acetyl-CoA. Diverted from the citric acid cycle, acetyl-CoA is funneled into mitochondrial ketogenesis via HMG-CoA synthase. Physiological ketosis (fasting ketosis, ketogenic diets) is a controlled, adaptive state (ketones 1--5 mmol/L) providing vital water-soluble fuel to the brain, heart, and skeletal muscle while sparing muscle proteolysis; blood pH remains within normal limits. In contrast, Pathological Ketoacidosis (most notably Diabetic Ketoacidosis [DKA], but also alcoholic ketoacidosis) is an uncontrolled, life-threatening emergency caused by absolute insulin deficiency and massive counter-regulatory hormone excess (glucagon, epinephrine, cortisol). Unrestrained lipolysis yields overwhelming ketone production ($>10--20$ mmol/L), exceeding systemic buffering capacity and causing a severe high anion gap metabolic acidosis, osmotic diuresis, dehydration, electrolyte depletion, and Kussmaul respiration.

\medterm{Lactate}-- The conjugate base of lactic acid and a normal product of glycolysis when pyruvate is reduced by lactate dehydrogenase. Lactate can be oxidized by tissues or transported to the liver for conversion back to glucose through the Cori cycle; persistent elevation may indicate impaired production-consumption balance.

\medterm{Lactate Dehydrogenase}-- A ubiquitous tetrameric cytoplasmic oxidoreductase that catalyzes the reversible interconversion of pyruvate and L-lactate with the simultaneous oxidation/reduction of NADH and $NAD^+$. Composed of varying combinations of two distinct 35-kDa subunits---heart (H) and muscle (M)---encoded by the LDHA and LDHB genes, forming five distinct tissue-specific isoenzymes: LDH-1 ($H_4$, myocardium, erythrocytes), LDH-2 ($H_3M_1$, reticuloendothelial system), LDH-3 ($H_2M_2$, lungs, spleen, platelets), LDH-4 ($HM_3$, kidneys, pancreas), and LDH-5 ($M_4$, skeletal muscle, hepatocytes). In anaerobic glycolysis (strenuous exercise, tissue hypoxia, mature erythrocytes), LDH reduces pyruvate to lactate, critically regenerating $NAD^+$ to permit continued flux through glyceraldehyde-3-phosphate dehydrogenase. Total serum LDH reference range is approximately 140--280 U/L. Serum LDH is a sensitive biomarker of cellular injury, hemolysis, myocardial infarction (characterised historically by an 'LDH flipped pattern', LDH-1 > LDH-2), pulmonary embolism, liver disease, and tumor burden in hematologic malignancies (e.g., lymphomas, leukemias) and germ cell tumors.

\medterm{Lactose Intolerance}-- A common gastrointestinal syndrome caused by deficiency of the brush-border enzyme lactase (lactase-phlorizin hydrolase), a $\beta$-galactosidase localized to the apical microvilli of mature enterocytes in the mid-jejunum. Lactase catalyzes the hydrolytic cleavage of the dietary disaccharide lactose into its constituent absorbable monosaccharides: D-glucose and D-galactose. In lactase deficiency, unhydrolyzed lactose cannot be absorbed across the intestinal mucosa and remains trapped within the gastrointestinal lumen. The osmotic force of unabsorbed lactose draws large volumes of water into the small intestine, producing osmotic diarrhea. Upon reaching the colon, resident anaerobic microflora ferment the undigested lactose into short-chain fatty acids (acetate, propionate, butyrate) and voluminous gases, including hydrogen ($H_2$), carbon dioxide ($CO_2$), and methane ($CH_4$). This gas production and mucosal distention produce characteristic clinical symptoms: abdominal cramping, bloating, flatulence, borborygmi, and explosive watery diarrhea within 30 minutes to two hours following dairy ingestion. The most common form is primary adult lactase deficiency (lactase non-persistence), an autosomal recessive trait affecting $>65\%$ of the world adult population, where lactase expression naturally declines after weaning; secondary lactase deficiency results from mucosal injury to brush-border enterocytes (celiac disease, viral gastroenteritis, Crohn disease). Diagnostic confirmation is achieved via the non-invasive hydrogen breath test (ingestion of lactose produces a rise in expired breath hydrogen $>20$ ppm) or stool acidity testing (pH $<5.5$). Management involves dietary lactose restriction, lactase enzyme replacement supplements, and calcium and vitamin D supplementation.

\medterm{LDL}-- Low-density lipoprotein, the primary circulating cholesterol-transport lipoprotein particle (density 1.019--1.063 g/mL, diameter ~22 nm), derived intravascularly from the catabolism of intermediate-density lipoprotein (IDL) and VLDL. LDL particles possess a hydrophobic core consisting almost entirely of cholesteryl esters and triglycerides, enveloped by an amphipathic phospholipid monolayer containing a single, unbroken copy of apolipoprotein B-100 (ApoB-100). ApoB-100 serves as the high-affinity ligand that binds to cell-surface LDL receptors (LDLR) on hepatocytes (clearing ~70\% of circulating LDL) and peripheral extrahepatic tissues. Upon binding, the LDLR-LDL complex undergoes clathrin-coated pit endocytosis; endosomal acidification dissociates the receptor (which recycles to the cell surface), while the particle is routed to the lysosome, where acid lipases hydrolyze cholesteryl esters into free cholesterol. Elevated intracellular cholesterol suppresses SREBP-2, downregulating both HMG-CoA reductase and LDL receptor transcription, and activates ACAT to esterify excess cholesterol for storage. When plasma LDL concentrations are elevated ($>100$ mg/dL, or $>70$ mg/dL in high-risk patients), LDL particles penetrate vascular endothelial junctions and become trapped in the arterial intima, where they undergo oxidative modification (oxLDL). Macrophage scavenger receptors (SR-A and CD36) take up oxLDL without negative feedback, transforming macrophages into lipid-laden foam cells that drive atherosclerotic plaque formation. Heterozygous Familial Hypercholesterolemia (LDLR mutation) causes severe hypercholesterolemia ($>300$ mg/dL), tendon xanthomas, and premature myocardial infarction before age 50.

\medterm{Leptin}-- A 16-kDa anorexigenic peptide hormone encoded by the LEP (ob) gene and secreted predominantly by white adipocytes in direct proportion to total body fat mass, functioning as the primary endocrine signal of long-term energy sufficiency to the central nervous system. Leptin crosses the blood-brain barrier via receptor-mediated transport and binds the long signaling isoform of the leptin receptor (LepRb / Ob-Rb) in the arcuate nucleus of the hypothalamus. Leptin receptor engagement activates the Janus kinase 2 / signal transducer and activator of transcription 3 (JAK2-STAT3) signaling pathway, resulting in two coordinated neuroendocrine actions: (1) Stimulation of anorexigenic pro-opiomelanocortin (POMC) and cocaine- and amphetamine-regulated transcript (CART) neurons, which release $\alpha$-melanocyte-stimulating hormone ($\alpha$-MSH) to activate melanocortin-4 receptors (MC4R), suppressing appetite and stimulating energy expenditure; and (2) Inhibition of orexigenic neuropeptide Y (NPY) and agouti-related peptide (AgRP) neurons. In starvation or rapid weight loss, falling leptin levels stimulate voracious hunger and suppress energy expenditure, thyroid hormone, and the reproductive axis (causing hypothalamic amenorrhea). Although congenital leptin deficiency or leptin receptor mutations cause severe early-onset morbid obesity (responsive to recombinant human leptin / metreleptin), the vast majority of human obesity is characterized by hyperleptinemia accompanied by central 'leptin resistance' due to defective transport across the blood-brain barrier and upregulation of intracellular feedback inhibitors (SOCS3 and PTP1B).

\medterm{Lesch-Nyhan Syndrome}-- A severe X-linked recessive disorder of purine metabolism caused by complete deficiency of the enzyme hypoxanthine-guanine phosphoribosyltransferase (HGPRT), encoded by the HPRT1 gene on the X chromosome. HGPRT functions as the core enzyme of the purine salvage pathway, catalyzing the condensation of 5-phosphoribosyl-1-pyrophosphate (PRPP) with free purine bases: hypoxanthine + PRPP $\to$ inosine monophosphate (IMP) and guanine + PRPP $\to$ guanosine monophosphate (GMP). Absence of HGPRT produces a dual metabolic catastrophe: (1) Free hypoxanthine and guanine cannot be salvaged and are diverted into purine catabolism, where xanthine oxidase oxidizes them into massive quantities of insoluble uric acid; and (2) Intracellular PRPP accumulates while IMP and GMP levels drop, massively disinhibiting the rate-limiting committed enzyme of de novo purine synthesis (PRPP amidotransferase), driving hyperactive purine production. Affected males present in infancy with 'orange sand' (sodium urate crystals) in their diapers, progressing to severe hyperuricemia, tophi, nephrolithiasis, and gouty arthritis. Neurologically, HGPRT deficiency disrupts dopamine neurodevelopment in the basal ganglia, manifesting as choreoathetosis, dystonia, severe spasticity, intellectual disability, and a pathognomonic compulsive self-mutilation behavior (severe biting of lips, tongue, and fingers, often requiring protective restraints or dental extraction). Treatment with xanthine oxidase inhibitors (allopurinol or febuxostat) successfully manages hyperuricemia and prevents renal failure and gout, but has no efficacy against the devastating neurological and behavioral manifestations.

\medterm{Leukemia}-- A heterogeneous group of malignant neoplastic proliferations of hematopoietic stem and progenitor cells in the bone marrow and peripheral blood, classified by cell lineage (myeloid vs. lymphoid) and clinical acuity (acute vs. chronic): (1) Acute Myeloblastic Leukemia (AML): clonal expansion of immature myeloid blasts ($>20\%$ in marrow), commonly associated with chromosomal translocations ($t(15;17)$ generating the PML-RARA fusion protein in acute promyelocytic leukemia, treated with all-trans retinoic acid [ATRA] and arsenic trioxide; Auer rods seen on smear); (2) Acute Lymphoblastic Leukemia (ALL): the most common childhood malignancy, characterized by lymphoblasts ($t(12;21)$ ETV6-RUNX1 with favorable prognosis, or $t(9;22)$ Philadelphia chromosome BCR-ABL1 with poor prognosis); (3) Chronic Myelogenous Leukemia (CML): a myeloproliferative disorder driven by the reciprocal translocation $t(9;22)(q34;q11)$, creating the BCR-ABL1 fusion gene encoding a constitutively active tyrosine kinase that drives massive granulocytosis, exquisitely responsive to BCR-ABL1 tyrosine kinase inhibitors (imatinib, dasatinib); and (4) Chronic Lymphocytic Leukemia (CLL): the most common adult leukemia in Western countries, characterized by clonal proliferation of mature-appearing CD5+/CD19+/CD23+ B-cells with smudge cells on peripheral blood smear. Blasts infiltrate bone marrow, crowding out normal hematopoiesis, presenting clinically with anemia (fatigue), thrombocytopenia (petechiae, bleeding), and neutropenia (recurrent opportunistic infections).

\medterm{Lipase}-- A glycoprotein enzyme (pancreatic triacylglycerol lipase) synthesized and secreted by pancreatic acinar cells into the duodenum, responsible for the primary digestion of dietary fats. Working in concert with pancreatic colipase (which anchors lipase to the bile-salt-coated lipid interface) and bile salts, pancreatic lipase selectively hydrolyzes dietary triacylglycerols at the $sn$-1 and $sn$-3 positions, yielding 2-monoacylglycerols and free fatty acids for micellar absorption by enterocytes. Serum lipase reference range is typically 10--140 U/L. Serum lipase is the preferred diagnostic biomarker for acute pancreatitis: it rises within 4--8 hours of pancreatic acinar injury, peaks at 24 hours, and remains elevated for 8--14 days (significantly longer than serum amylase). A serum lipase elevation $>3$ times the upper limit of normal provides high diagnostic sensitivity (85--100\%) and specificity (95--99\%) for acute pancreatitis, far outperforming amylase which is confounded by salivary and renal disorders.

\medterm{Lipid}-- Diverse group of hydrophobic molecules including triglycerides (energy storage), phospholipids (membranes), cholesterol (membranes, steroid precursor), fat-soluble vitamins (A, D, E, K). Triglycerides: glycerol + 3 fatty acids. Fatty acids: saturated (no double bonds) vs. unsaturated (one or more double bonds). Essential fatty acids: linoleic acid (omega-6), alpha-linolenic acid (omega-3).

\medterm{Lipogenesis}-- The anabolic biochemical pathway by which excess dietary carbohydrates and metabolic carbon intermediates are converted into fatty acids and subsequently esterified into triacylglycerols (triglycerides) for long-term energy storage. Occurring predominantly in the cytoplasm of hepatocytes and adipocytes, lipogenesis is powerfully stimulated in the postprandial state by high insulin and abundant ATP. When cellular energy charge is elevated, isocitrate dehydrogenase in the TCA cycle is inhibited by ATP and NADH, causing mitochondrial citrate to accumulate. Citrate is transported across the inner mitochondrial membrane into the cytoplasm via the tricarboxylate transport protein. In the cytosol, ATP-citrate lyase (induced by insulin) cleaves citrate into oxaloacetate and acetyl-CoA. Acetyl-CoA is then carboxylated to malonyl-CoA by the rate-limiting, committed enzyme acetyl-CoA carboxylase (ACC, activated allosterically by citrate and covalently by insulin-mediated dephosphorylation). Fatty acid synthase (FAS) subsequently utilizes acetyl-CoA, malonyl-CoA, and NADPH (supplied by the pentose phosphate pathway and malic enzyme) to synthesize palmitate. Newly formed fatty acids are activated to fatty acyl-CoAs and esterified to glycerol-3-phosphate (derived from glycolysis via DHAP or glycerol phosphorylation) to generate triacylglycerols, which the liver packages into VLDL for peripheral delivery.

\medterm{Lipolysis}-- The catabolic biochemical process by which stored triacylglycerols (triglycerides) within adipocyte cytoplasmic lipid droplets are sequentially hydrolyzed into free fatty acids (FFAs) and glycerol for release into the bloodstream. Lipolysis proceeds through a cascade of three distinct neutral lipases: (1) Adipose Triglyceride Lipase (ATGL): hydrolyzes triacylglycerol at the $sn-1$ or $sn-2$ position to generate diacylglycerol (DAG) and a free fatty acid, activated by CGI-58 upon perilipin phosphorylation; (2) Hormone-Sensitive Lipase (HSL): the primary hormonally regulated, rate-limiting enzyme that hydrolyzes diacylglycerol into monoacylglycerol (MAG) and a second FFA; and (3) Monoacylglycerol Lipase (MGL): constitutively hydrolyzes monoacylglycerol into glycerol and a third FFA. Hormonally, lipolysis is stimulated during fasting, exercise, stress, or hypoglycemia by catecholamines (epinephrine, norepinephrine), glucagon, ACTH, and cortisol via GPCR-$G_{\alpha s}$ activation, elevating cAMP and activating protein kinase A (PKA). PKA dual-phosphorylates HSL (promoting translocation to the droplet) and perilipin-1 (releasing CGI-58). In contrast, insulin potently inhibits lipolysis by activating phosphodiesterase-3B (PDE-3B) and protein phosphatase 1, degrading cAMP and dephosphorylating HSL. Released free fatty acids bind serum albumin and are delivered to heart, skeletal muscle, and liver for beta-oxidation, while glycerol is cleared by the liver for gluconeogenesis or glycolysis via glycerol kinase.

\medterm{Lipoprotein}-- Complex of lipids and proteins transporting hydrophobic lipids in aqueous blood. Classes: chylomicrons (dietary lipids), VLDL (endogenous triglycerides), LDL (cholesterol delivery), HDL (reverse cholesterol transport). ApoB-100 in LDL; apoA-1 in HDL. Elevated LDL cholesterol is a major risk factor for atherosclerosis. Statins lower LDL by upregulating LDL receptors.

\medterm{Lipoprotein Lipase}-- A water-soluble homodimeric enzyme bound to the luminal surface of capillary endothelial cells by negatively charged heparan sulfate proteoglycans, highly expressed in capillaries of adipose tissue, skeletal muscle, and cardiac muscle. LPL catalyzes the intravascular hydrolysis of core triglycerides from circulating triglyceride-rich lipoproteins (chylomicrons and VLDL) into free fatty acids and glycerol. Released fatty acids are taken up by adjacent adipocytes for re-esterification and storage, or by myocytes for $\beta$-oxidation. LPL catalytic activation requires apolipoprotein C-II (ApoC-II), an essential protein cofactor present on the surface of chylomicrons and VLDL, and is competitively inhibited by ApoC-III. In adipose tissue, LPL synthesis and luminal translocation are strongly stimulated by insulin in the fed state, whereas in heart and skeletal muscle, LPL is upregulated during fasting. Genetic deficiency of LPL or ApoC-II results in Familial Chylomicronemia Syndrome (Type I hyperlipoproteinemia), characterized by massive hypertriglyceridemia (>1,000 mg/dL), eruptive xanthomas, lipemia retinalis, hepatosplenomegaly, and recurrent acute pancreatitis without premature atherosclerosis.

\medterm{Liver Enzymes}-- A standard panel of serum diagnostic enzymes used clinically to assess hepatocellular integrity and biliary tract function. Alanine aminotransferase (ALT, reference 7--56 U/L) is a cytosolic enzyme predominantly expressed in hepatocytes, making it the most specific biomarker for hepatocellular injury. Aspartate aminotransferase (AST, reference 10--40 U/L) exists as cytosolic (20\%) and mitochondrial (80\%) isozymes distributed across liver, cardiac muscle, skeletal muscle, and kidneys. An AST-to-ALT ratio $>2:1$ is classically indicative of alcoholic liver disease, attributable to alcohol-induced mitochondrial injury and deficiency of pyridoxal phosphate (vitamin B6, required more by ALT than AST). Marked transaminase elevations ($>1,000$ U/L) signal acute viral hepatitis, ischemic hepatitis ('shock liver'), or acetaminophen toxicity. In contrast, alkaline phosphatase (ALP, 44--147 U/L) is localized to the canalicular membrane of hepatocytes and bone osteoblasts; its elevation indicates cholestasis or infiltrative disease. Gamma-glutamyl transferase (GGT, 9--48 U/L) is biliary-specific and induced by alcohol; an elevated GGT confirms that an elevated ALP originates from the hepatobiliary tree rather than bone.

\medterm{Lymphoma}-- A diverse group of clonal hematologic malignancies originating from mature lymphocytes or their committed precursors, presenting primarily as discrete nodal or extranodal solid tumor masses. Historically and clinically divided into two major categories: (1) Hodgkin Lymphoma (HL): characterized pathognomonic multinucleated neoplastic giant cells of B-cell origin termed Reed-Sternberg cells (classically 'owl-eyed' binucleated cells expressing CD15 and CD30, negative for CD20 and CD45) residing within a rich non-neoplastic inflammatory background of reactive T-cells, eosinophils, and histiocytes. Subtypes include nodular sclerosis (most common, young women, mediastinal mass), mixed cellularity (associated with EBV), lymphocyte-rich, and lymphocyte-depleted. HL typically spreads in an orderly, contiguous manner along adjacent lymph node chains, frequently presenting with painless cervical lymphadenopathy and systemic 'B symptoms' (unexplained fever, drenching night sweats, and weight loss $>10\%$); and (2) Non-Hodgkin Lymphoma (NHL): a heterogeneous spectrum ($>85\%$ B-cell origin, remainder T-cell or NK-cell) characterized by non-contiguous, diffuse nodal and extranodal involvement. Prominent B-cell NHL subtypes include Diffuse Large B-Cell Lymphoma (DLBCL, the most common aggressive lymphoma), Follicular Lymphoma ($t(14;18)$ overexpressing anti-apoptotic BCL-2), Burkitt Lymphoma ($t(8;14)$ overexpressing c-MYC, showing a 'starry sky' histopathology, associated with EBV), and Mantle Cell Lymphoma ($t(11;14)$ overexpressing cyclin D1).

\medterm{Magnesium}-- The second most abundant intracellular cation ($Mg^{2+}$, normal serum concentration: 1.7--2.2 mg/dL or 0.7--1.1 mmol/L), essential for the structural integrity and catalytic activity of over 300 enzyme systems. In biological systems, virtually all ATP must be chelated to magnesium ($Mg^{2+}$-ATP complex) to undergo hydrolysis or transfer phosphate groups; consequently, magnesium is an obligatory cofactor for all kinases (hexokinase, PFK-1, creatine kinase), DNA and RNA polymerases, ATPases, and adenylate cyclase. Magnesium also regulates neuromuscular excitability: it acts as a physiological calcium channel blocker at the neuromuscular junction (inhibiting presynaptic acetylcholine release) and stabilizes cardiac and neuronal cell membranes. Intestinal absorption occurs in the ileum and colon (via claudins and TRPM6), and renal handling involves reabsorption of 60--70\% in the cortical thick ascending limb of Henle (via paracellin-1/claudin-16). Hypomagnesemia ($<1.7$ mg/dL) results from chronic alcoholism, malnutrition, proton pump inhibitors, loop/thiazide diuretics, diarrhea, or aminoglycosides. Clinically, hypomagnesemia causes tremors, tetany, Chvostek and Trousseau signs, hyperreflexia, and fatal cardiac arrhythmias (prolonged QT interval, ventricular tachycardia, and Torsades de Pointes). Crucially, hypomagnesemia causes refractory hypokalemia (due to loss of intracellular $Mg^{2+}$ inhibition on renal outer medullary potassium [ROMK] channels, causing massive renal potassium wasting) and secondary hypocalcemia (due to impaired parathyroid hormone [PTH] release and end-organ PTH resistance). Intravenous magnesium sulfate is the definitive treatment for Torsades de Pointes, eclampsia/preeclampsia seizure prophylaxis, and severe acute asthma exacerbations. Hypermagnesemia ($>2.5$ mg/dL), seen in renal failure combined with magnesium-containing antacids or laxatives, causes loss of deep tendon reflexes, flaccid paralysis, respiratory depression, bradycardia, and cardiac arrest, treated immediately with intravenous calcium gluconate.

\medterm{Malate Dehydrogenase}-- An oxidoreductase enzyme that catalyzes the reversible oxidation of L-malate to oxaloacetate with the simultaneous reduction of $NAD^+$ to NADH. Two distinct compartmentalized isoenzymes exist in eukaryotic cells: mitochondrial malate dehydrogenase (MDH2), a critical component of the citric acid cycle (Krebs cycle), and cytosolic malate dehydrogenase (MDH1). In the TCA cycle, MDH2 drives the final step, regenerating oxaloacetate; although the standard free-energy change is endergonic ($\Delta G^{\circ\prime} = +29.7$ kJ/mol), the reaction proceeds readily in vivo because the citrate synthase reaction constantly pulls oxaloacetate concentrations to extremely low levels. Together, MDH1 and MDH2 form the catalytic backbone of the malate-aspartate shuttle, which transfers reducing equivalents of cytosolic NADH generated during glycolysis into the mitochondrial matrix (yielding ~2.5 ATP per NADH via Complex I) in cardiac muscle, liver, and kidneys, bypassing the impermeable inner mitochondrial membrane.

\medterm{Maple Syrup Urine Disease}-- A rare autosomal recessive inborn error of amino acid metabolism caused by deficiency in the branched-chain $\alpha$-ketoacid dehydrogenase (BCKDH) multienzyme complex. The BCKDH complex, localized in the inner mitochondrial membrane, resembles pyruvate dehydrogenase and $\alpha$-ketoglutarate dehydrogenase, requiring five cofactors: thiamine pyrophosphate (TPP, vitamin B1), lipoate, coenzyme A, FAD, and NAD+. The enzyme catalyzes the oxidative decarboxylation of $\alpha$-ketoacids derived from the transamination of the three branched-chain amino acids (BCAAs): leucine, isoleucine, and valine. BCKDH failure causes massive accumulation of leucine, isoleucine, valine, and their corresponding neurotoxic $\alpha$-ketoacids ($\alpha$-ketoisocaproate, $\alpha$-keto-$\beta$-methylvalerate, and $\alpha$-ketoisovalerate) in tissues, cerebrospinal fluid, and plasma. Elevated isoleucine imparts a sweet, caramelized burnt-sugar or maple-syrup odor to the urine and earwax. Neonates present within the first week of life with poor feeding, lethargy, alternating hypertonia and hypotonia (dystonic 'fencing' posture), cerebral edema, coma, and central hypoventilation due to leucine-induced neurotoxicity and brain swelling. Diagnosis is established by elevated BCAAs and presence of alloisoleucine on plasma amino acid chromatography. Management consists of emergency hemodialysis for acute crises, lifelong dietary restriction of BCAAs using specialized amino acid formulas, and thiamine supplementation in thiamine-responsive variants.

\medterm{Messenger RNA (mRNA)}-- Single-stranded RNA molecule transcribed from DNA template, carrying genetic code from nucleus to ribosomes for protein synthesis. Contains 5' cap (methylguanosine), 5' UTR, coding sequence (codons), 3' UTR, and poly-A tail. Eukaryotic mRNA undergoes processing: splicing (intron removal, exon ligation), capping, polyadenylation. mRNA stability determines protein output.

\medterm{Metabolic Acidosis}-- A primary acid-base disturbance characterized by a pathological reduction in arterial plasma bicarbonate ($[HCO_3^-] < 22$ mEq/L) and a fall in systemic blood pH ($\text{pH} < 7.35$), caused by the abnormal accumulation of nonvolatile acids, excessive gastrointestinal or renal loss of bicarbonate, or impaired renal tubular acid excretion. Compensatory respiratory hyperventilation ensues rapidly (within minutes), stimulating medullary chemoreceptors to blow off volatile $CO_2$; Winter's formula calculates the expected compensatory $P_{CO_2}$: $P_{CO_2} = (1.5 \times [HCO_3^-]) + 8 \pm 2$ mmHg (a measured $P_{CO_2}$ above this indicates concomitant respiratory acidosis; below indicates concomitant respiratory alkalosis). Clinical evaluation hinges upon calculating the serum Anion Gap ($[Na^+] - [Cl^- + HCO_3^-]$, normal 8--12 mEq/L): (1) High Anion Gap Metabolic Acidosis (HAGMA, $>12$ mEq/L): unmeasured nonvolatile anions accumulate, driven by lactic acidosis (tissue hypoperfusion, sepsis, shock), ketoacidosis (diabetic, alcoholic, starvation), end-stage renal failure (sulfate, phosphate retention), or toxic ingestions (methanol, ethylene glycol, salicylates); and (2) Normal Anion Gap / Hyperchloremic Metabolic Acidosis (NAGMA, 8--12 mEq/L): bicarbonate is lost and replaced by chloride to preserve electrical neutrality, caused by severe diarrhea (loss of intestinal $HCO_3^-$), renal tubular acidosis (Types 1, 2, and 4), or carbonic anhydrase inhibitors (acetazolamide). Severe metabolic acidosis reduces myocardial contractility, provokes peripheral vasodilation, causes hyperkalemia (due to $H^+/K^+$ cellular shifting), and produces deep, rapid Kussmaul respirations.

\medterm{Metabolic Alkalosis}-- A primary acid-base disturbance characterized by an elevation in arterial plasma bicarbonate ($[HCO_3^-] > 28$ mEq/L) and a rise in blood pH ($\text{pH} > 7.45$), accompanied by secondary compensatory respiratory hypoventilation (increasing $P_{CO_2}$, with expected compensation: $\Delta P_{CO_2} = 0.7 \times \Delta [HCO_3^-] \pm 5$ mmHg, though hypoventilation is clinically capped by hypoxemic respiratory drive). Metabolic alkalosis develops in two phases: generation (net loss of hydrogen ions or addition of bicarbonate) and maintenance (impaired renal capacity to excrete excess bicarbonate due to volume depletion, hypokalemia, or chloride depletion). Clinically and pathophysiologically categorized by urinary chloride concentration: (1) Chloride-Responsive (Saline-Responsive, urinary $Cl^- < 20$ mEq/L): caused by loss of gastric acid containing HCl (vomiting, nasogastric suction) or renal chloride and volume wasting from loop or thiazide diuretics; volume contraction and chloride depletion stimulate renal aldosterone secretion and proximal tubular bicarbonate reclamation, maintaining alkalosis. This form corrects readily with intravenous isotonic ($0.9\%$) saline resuscitation and potassium chloride replacement; and (2) Chloride-Resistant (Saline-Unresponsive, urinary $Cl^- > 20$ mEq/L): typically accompanied by hypertension and hypervolemia, driven by mineralocorticoid excess (primary hyperaldosteronism / Conn syndrome, Cushing syndrome, licorice ingestion, or genetic channelopathies like Liddle syndrome), where unrestrained cortical collecting duct epithelial sodium channel (ENaC) activity drives continuous $H^+$ and $K^+$ excretion via intercalated and principal cells. Metabolic alkalosis shifts ionized calcium onto albumin, provoking acute tetany, paresthesias, neuromuscular irritability, and cardiac arrhythmias.

\medterm{Metabolic Syndrome}-- A constellation of interconnected metabolic, biochemical, and physiological derangements that synergistically multiply the risk of developing atherosclerotic cardiovascular disease (a 2-fold increase) and type 2 diabetes mellitus (a 5-fold increase). According to the harmonized National Cholesterol Education Program (NCEP ATP III) criteria, diagnosis requires the presence of any three or more of the following five components: (1) Abdominal (central) obesity: waist circumference $\ge 102$ cm ($\ge 40$ in) in men or $\ge 88$ cm ($\ge 35$ in) in women; (2) Hypertriglyceridemia: fasting plasma triglycerides $\ge 150$ mg/dL (1.7 mmol/L), or specific drug treatment; (3) Low HDL cholesterol: serum HDL-C $<40$ mg/dL (1.0 mmol/L) in men or $<50$ mg/dL (1.3 mmol/L) in women, or specific drug treatment; (4) Systemic hypertension: systolic blood pressure $\ge 130$ mmHg and/or diastolic blood pressure $\ge 85$ mmHg, or antihypertensive medication use; and (5) Impaired fasting glucose: fasting plasma glucose $\ge 100$ mg/dL (5.6 mmol/L), or drug treatment for elevated glucose. The core pathophysiological driver is visceral adiposity coupled with systemic insulin resistance: hypertrophied visceral adipocytes exhibit high lipolytic rates, releasing an excess flux of non-esterified free fatty acids (NEFAs) directly into the portal circulation. This hepatic lipid overload drives hepatic steatosis, accelerates VLDL synthesis (triggering hypertriglyceridemia and low HDL via CETP exchange), promotes gluconeogenesis, impairs skeletal muscle glucose uptake, and fosters a pro-inflammatory and pro-thrombotic state characterized by elevated CRP, TNF-$\alpha$, IL-6, and plasminogen activator inhibitor-1 (PAI-1).

\medterm{Metabolism}-- Sum of all chemical reactions in the body: catabolism (breakdown, releases energy) and anabolism (synthesis, consumes energy). Central pathways: glycolysis, TCA cycle, oxidative phosphorylation, pentose phosphate pathway. Metabolic regulation ensures energy homeostasis. Disorders: inborn errors of metabolism (PKU, glycogen storage diseases), diabetes, obesity.

\medterm{Michaelis-Menten Kinetics}-- A model of enzyme kinetics describing the rate of enzymatic reactions as a function of substrate concentration. It is defined by the equation $V = \frac{V_{\max}[S]}{K_m + [S]}$, where $V$ is reaction velocity, $V_{\max}$ is maximum velocity, and $K_m$ (Michaelis constant) is the substrate concentration at which the reaction velocity is half of $V_{\max}$. $K_m$ is inversely related to enzyme-substrate affinity. Competitive inhibitors increase $K_m$ without altering $V_{\max}$, whereas noncompetitive inhibitors decrease $V_{\max}$ without altering $K_m$.

\medterm{MicroRNA}-- Small (approximately 21--25 nucleotides), conserved, non-coding single-stranded RNA molecules that function as master post-transcriptional regulators of gene expression in eukaryotic cells. MicroRNAs (miRNAs) are transcribed from the genome by RNA polymerase II as primary transcripts (pri-miRNAs) possessing local hairpin stem-loop structures. In the nucleus, the microprocessor complex (comprising the RNase III enzyme Drosha and DGCR8) cleaves the pri-miRNA into a ~70-nucleotide precursor hairpin (pre-miRNA), which exportin-5 transports through nuclear pores into the cytoplasm. In the cytoplasm, the RNase III endonuclease Dicer cleaves the pre-miRNA hairpin loop, yielding a mature double-stranded miRNA duplex. One strand (the guide strand) is loaded into an Argonaute protein to assemble the RNA-Induced Silencing Complex (RISC). The RISC complex binds complementary target messenger RNAs predominantly within their 3'-untranslated regions (3'-UTRs) via a conserved 6- to 8-nucleotide 'seed sequence'. Perfect base-pairing complementarity triggers Argonaute-mediated endonucleolytic cleavage and mRNA degradation, whereas imperfect complementarity (the predominant mechanism in mammals) induces translational repression and accelerated mRNA deadenylation. MiRNAs coordinate embryonic development, cell differentiation, and metabolic balance; dysregulated miRNAs act as oncogenes ('oncomirs') or tumor suppressors in human malignancies.

\medterm{Milk-Alkali Syndrome}-- Also designated calcium-alkali syndrome, a clinical triad consisting of hypercalcemia, metabolic alkalosis, and acute or chronic renal dysfunction resulting from the ingestion of large quantities of calcium (typically $>2--4$ g/day) and absorbable alkaline salts (most commonly calcium carbonate antacids or over-the-counter calcium supplements). The absorbed calcium suppresses parathyroid hormone (PTH) secretion; however, the massive concurrent alkali load overwhelms the renal capacity to excrete bicarbonate, particularly because hypercalcemia impairs the renal concentrating ability (inducing nephrogenic diabetes insipidus via calcium-sensing receptor activation in the medullary collecting duct) and causes renal vasoconstriction, reducing glomerular filtration rate. Volume contraction from polyuria and vomiting stimulates proximal tubular sodium and bicarbonate reabsorption, perpetuating the metabolic alkalosis and worsening hypercalcemia in a self-reinforcing vicious cycle. Patients present with nausea, vomiting, lethargy, confusion, polyuria, polydipsia, band keratopathy, and metastatic soft-tissue calcification. Treatment involves immediate cessation of calcium and alkali ingestion, vigorous volume resuscitation with intravenous isotonic saline, and loop diuretics once volume status is restored.

\medterm{mTOR}-- Mechanistic (or mammalian) target of rapamycin, an evolutionarily conserved 289-kDa serine/threonine protein kinase belonging to the phosphatidylinositol 3-kinase-related kinase (PIKK) family that acts as the central coordinator of cell growth, proliferation, protein synthesis, and survival. mTOR operates as the catalytic core of two structurally and functionally distinct multiprotein complexes: (1) mTOR Complex 1 (mTORC1, containing Raptor): Sensitive to rapamycin, mTORC1 integrates four major upstream regulatory cues: growth factors (via the PI3K-Akt-TSC1/2-Rheb signaling cascade), intracellular amino acid availability (sensed via Rag GTPases at the lysosomal surface), oxygen tension, and cellular energy status (suppressed by AMPK during ATP depletion). When active, mTORC1 promotes anabolic biomass accumulation by phosphorylating p70S6 kinase (p70S6K, which phosphorylates ribosomal protein S6) and eukaryotic translation initiation factor 4E-binding protein-1 (4E-BP1, liberating eIF4E to initiate cap-dependent translation), while suppressing catabolic processes such as autophagy (by phosphorylating ULK1); (2) mTOR Complex 2 (mTORC2, containing Rictor): Insensitive to acute rapamycin exposure, mTORC2 regulates spatial cell organization, actin cytoskeleton dynamics, and phosphorylates Akt at Ser473 for maximal Akt kinase activation. Hyperactivation of mTOR signaling occurs widely across human cancers, while mTOR inhibitors (rapamycin/sirolimus, everolimus) are utilized clinically as immunosuppressants in organ transplantation and as targeted antineoplastic agents.

\medterm{Multiple Myeloma}-- A neoplastic clonal plasma cell dyscrasia characterized by the uncontrolled, malignant proliferation of monoclonal plasma cells in the bone marrow, leading to excessive production of non-functional monoclonal immunoglobulins or free immunoglobulin light chains (M-protein / paraprotein). Most commonly, the paraprotein is IgG (approximately 55\%) or IgA (20\%), accompanied by monoclonal kappa ($\kappa$) or lambda ($\lambda$) light chains that are filtered across the glomerulus and excreted in urine as Bence-Jones proteins. Overproduction of RANKL and suppression of osteoprotegerin by malignant plasma cells activates osteoclasts while inhibiting osteoblasts, generating extensive osteolytic bone destruction. The classic clinical features are summarized by the 'CRAB' criteria: (1) Hypercalcemia (resulting from severe osteoclastic bone resorption); (2) Renal insufficiency (primarily myeloma cast nephropathy, where excess filtered light chains precipitate with Tamm-Horsfall mucoprotein in distal tubules, causing tubular obstruction and toxic damage); (3) Anemia (normocytic normochromic, due to bone marrow infiltration by neoplastic plasma cells and relative erythropoietin deficiency); and (4) Bone lesions (punched-out lytic lesions on skeletal radiography, bone pain, and pathological fractures). Diagnostic evaluation includes serum and urine protein electrophoresis (SPEP/UPEP demonstrating a narrow, dense monoclonal 'M-spike'), serum free light chain assay, bone marrow aspirate/biopsy ($>10\%$ clonal plasma cells), and skeletal imaging.

\medterm{Mutation}-- A permanent, heritable or acquired alteration in the primary nucleotide sequence of an organism's genome, originating from spontaneous replication errors by DNA polymerases, endogenous chemical damage (depurination, deamination, oxidative stress), or environmental mutagens (ionizing radiation, UV light, chemical alkylating agents). Genomic mutations are classified by structural type: (1) Point Mutations (single base substitutions): (a) Silent mutations (synonymous change encoding the same amino acid due to genetic code degeneracy); (b) Missense mutations (amino acid substitution, e.g., Glu6Val in sickle cell anemia); and (c) Nonsense mutations (conversion of an amino acid codon to a premature stop codon [UAA, UAG, UGA], leading to truncated, non-functional proteins or nonsense-mediated mRNA decay); (2) Frameshift Mutations: insertions or deletions of a number of nucleotides not divisible by three, which shifts the translational reading frame, altering all downstream codons and typically producing a premature stop codon (e.g., severe Duchenne muscular dystrophy, cystic fibrosis $\Delta$F508); (3) Splice-Site Mutations: mutations at conserved 5' donor (GU) or 3' acceptor (AG) intron-exon boundaries, leading to aberrant intron retention or exon skipping (common in $\beta$-thalassemias); and (4) Trinucleotide Repeat Expansions: unstable microsatellite expansions that exhibit anticipation across generations (e.g., CAG repeats in Huntington disease; CGG in Fragile X syndrome). Somatic mutations accumulate in non-germline cells, driving oncogenesis, while germline mutations are transmitted to offspring.

\medterm{NAD+ (Nicotinamide Adenine Dinucleotide)}-- Coenzyme involved in redox reactions, existing in oxidized (NAD+) and reduced (NADH) forms. Accepts electrons and hydrogen (NAD+ + 2e- + H+ -> NADH). Key roles: glycolysis (glyceraldehyde-3-phosphate dehydrogenase), TCA cycle (isocitrate, alpha-ketoglutarate, malate dehydrogenases), beta-oxidation. NADH donates electrons to electron transport chain. NAD+ also substrate for sirtuins (deacetylases) and PARPs (DNA repair).

\medterm{Niacin (Vitamin B3)}-- A water-soluble vitamin encompassing nicotinic acid and nicotinamide, which are converted intracellularly into the essential pyridine dinucleotide coenzymes: nicotinamide adenine dinucleotide ($NAD^+$) and nicotinamide adenine dinucleotide phosphate ($NADP^+$). In addition to dietary intake (meat, fish, enriched grains), niacin can be synthesized de novo in hepatocytes from the essential amino acid L-tryptophan via the kynurenine pathway (requiring 60 mg of dietary tryptophan to yield 1 mg of niacin), an enzymatic conversion that requires pyridoxal phosphate (vitamin B6) and riboflavin (vitamin B2) as cofactors. $NAD^+$ serves as the primary electron acceptor in catabolic oxidative pathways (glycolysis, pyruvate dehydrogenase, TCA cycle, fatty acid beta-oxidation), generating NADH to drive mitochondrial oxidative phosphorylation. In contrast, NADPH functions as the primary electron donor in reductive biosynthetic pathways (fatty acid and cholesterol synthesis, steroidogenesis) and protects cells against oxidative damage by maintaining reduced glutathione. Severe niacin deficiency causes Pellagra, classically characterized by the '4 D's': Dermatitis (a symmetrical, hyperkeratotic, hyperpigmented rash on sun-exposed skin, classically forming Casal's necklace around the neck), Diarrhea (due to mucosal atrophy of the gastrointestinal tract), Dementia (memory loss, hallucinations, encephalopathy), and Death if untreated. Pellagra occurs in severe malnutrition, alcoholism, Hartnup disease (defective renal and intestinal neutral amino acid transport, losing tryptophan), and carcinoid syndrome (where serotonin synthesis consumes $>70\%$ of tryptophan pools). Pharmacologically, high-dose nicotinic acid lowers plasma triglycerides and LDL while raising HDL, frequently producing prostaglandin-mediated cutaneous flushing that can be blunted by aspirin.

\medterm{Norepinephrine}-- Noradrenaline, an essential catecholamine that functions both as the primary chemical neurotransmitter released by postganglionic sympathetic neurons and as a circulating hormone secreted by chromaffin cells of the adrenal medulla (~20\% of output). Synthesized from dopamine within synaptic storage vesicles by the copper- and ascorbic acid-dependent enzyme dopamine $\beta$-hydroxylase (DBH). Norepinephrine is a potent agonist of $\alpha_1$- and $\alpha_2$-adrenergic receptors, possesses robust activity at myocardial $\beta_1$-receptors, but exhibits minimal affinity for $\beta_2$-receptors (distinguishing its physiological actions from epinephrine). Its predominant hemodynamic effect is mediated via vascular $\alpha_1$-receptors coupled to $G_{\alpha q}$, triggering phospholipase C activation, $IP_3$ generation, and intracellular calcium release, which provokes intense, widespread systemic arterial and venous vasoconstriction. This markedly increases total peripheral resistance (systemic vascular resistance, SVR) and elevates both systolic and diastolic blood pressure; the severe rise in afterload triggers a baroreceptor-mediated vagal reflex that often results in reflex bradycardia despite direct myocardial $\beta_1$ stimulation. Clinically, intravenous norepinephrine is the gold-standard, first-line vasopressor recommended in international guidelines for the hemodynamic resuscitation of septic, cardiogenic, and distributive shock.

\medterm{Oxaloacetate}-- A critical four-carbon dicarboxylic acid intermediate situated at the nexus of multiple major metabolic pathways, including the citric acid cycle (Krebs cycle), gluconeogenesis, the malate-aspartate shuttle, and amino acid transamination. In the mitochondrial matrix, oxaloacetate condenses with the acetyl moiety of acetyl-CoA in the first step of the citric acid cycle, catalyzed by citrate synthase to form citrate. It is regenerated in the cycle's final step by the oxidation of L-malate via malate dehydrogenase (MDH2). In gluconeogenesis, oxaloacetate is synthesized from pyruvate by the anaplerotic, biotin-dependent enzyme pyruvate carboxylase; because the inner mitochondrial membrane lacks a dedicated transporter for oxaloacetate, it is converted to malate (via mitochondrial MDH2) or transaminated to aspartate (via AST), transported into the cytosol, and reconverted to oxaloacetate. Cytosolic phosphoenolpyruvate carboxykinase (PEPCK) then decarboxylates and phosphorylates oxaloacetate using GTP to generate phosphoenolpyruvate (PEP). In amino acid metabolism, oxaloacetate undergoes reversible transamination with glutamate via AST (PLP cofactor) to generate aspartate, which donates nitrogen atoms to the urea cycle and purine/pyrimidine synthesis. Depletion of oxaloacetate during starvation or diabetic ketoacidosis (due to maximal gluconeogenic drain) impairs citrate synthase, shunting acetyl-CoA into ketogenesis.

\medterm{Oxidative Phosphorylation}-- Process generating ATP using energy released by electron transfer through the electron transport chain. Protons pumped across inner mitochondrial membrane create electrochemical gradient (proton motive force). ATP synthase (Complex V) uses proton flow to phosphorylate ADP to ATP. P/O ratio: ~2.5 ATP per NADH, ~1.5 ATP per FADH2. Total: ~30-32 ATP per glucose molecule. Inhibited by rotenone (Complex I), antimycin A (Complex III), cyanide/azide (Complex IV), oligomycin (ATP synthase). Uncouplers (DNP, thermogenin) dissipate gradient as heat.

\medterm{Oxidative Stress}-- A physiological or pathological state of cellular disruption characterized by an imbalance between the excessive generation of reactive oxygen species (ROS: superoxide anion [$O_2^{\bullet-}$], hydrogen peroxide [$H_2O_2$], hydroxyl radical [$^{\bullet}OH$]) or reactive nitrogen species (RNS: peroxynitrite [$ONOO^-$]) and the biological capacity of cellular antioxidant systems to detoxify these reactive intermediates or repair the resulting damage. Under basal conditions, ROS are produced as unavoidable byproducts of mitochondrial oxidative phosphorylation (Complexes I and III) and enzymatic reactions (NADPH oxidase, xanthine oxidase, cytochrome P450 enzymes). When ROS production surpasses endogenous enzymatic defenses (superoxide dismutase, catalase, glutathione peroxidase) and non-enzymatic scavengers (glutathione, vitamins C and E), unrestrained free radical attack damages all major classes of biological macromolecules: (1) Lipid Peroxidation: hydroxyl radicals attack methylene carbons in polyunsaturated fatty acids of cell and organelle membranes, generating lipid peroxyl radicals that propagate autocatalytic membrane destruction, forming cytotoxic aldehydes (malondialdehyde [MDA], 4-hydroxynonenal [4-HNE]) and disrupting membrane fluidity and barrier function; (2) Protein Oxidation: ROS oxidize amino acid side chains (forming carbonyls, sulfoxides, and disulfides), leading to enzyme inactivation, structural misfolding, and cross-linking; and (3) Nucleic Acid Damage: hydroxyl radicals oxidize DNA bases, predominantly generating 8-hydroxy-2'-deoxyguanosine (8-OHdG), which causes GC-to-TA transversion mutations during replication, as well as single- and double-strand DNA breaks. Oxidative stress is a central pathogenic mechanism in atherosclerosis, neurodegenerative disorders (Alzheimer, Parkinson), ischemia-reperfusion injury, diabetic microvascular complications, and aging.

\medterm{Pantothenic Acid (Vitamin B5)}-- A water-soluble B-complex vitamin composed of pantoic acid linked to $\beta$-alanine by an amide bond. Pantothenic acid is phosphorylated and converted by cells into two critical bioenergetic components: (1) Coenzyme A (CoA-SH), an indispensable acyl-group carrier containing a reactive terminal thiol (-SH) group. CoA-SH activates carboxylic acids by forming high-energy thioester bonds, yielding essential metabolic intermediates such as acetyl-CoA, succinyl-CoA, malonyl-CoA, and fatty acyl-CoAs that drive the citric acid cycle, fatty acid beta-oxidation, cholesterol synthesis, and the synthesis of ketone bodies and acetylcholine; and (2) The 4'-phosphopantetheine prosthetic arm of the acyl carrier protein (ACP) domain of fatty acid synthase, which tethers growing fatty acid chains during cytoplasmic de novo lipogenesis. Because pantothenic acid is ubiquitous in virtually all natural food sources (from the Greek 'pantothen', meaning 'from everywhere'), isolated nutritional deficiency is exceptionally rare. Experimental severe deficiency produces 'burning feet syndrome' (erythromelalgia-like paresthesias and hyperesthesia of the lower extremities), muscle cramps, fatigue, gastrointestinal distress, and impaired adrenal steroidogenesis.

\medterm{Penicillin}-- Beta-lactam antibiotic inhibiting bacterial cell wall synthesis. Targets penicillin-binding proteins (PBPs), transpeptidases that cross-link peptidoglycan. Bactericidal against growing bacteria. Spectrum: gram-positive (Streptococcus, Staphylococcus), some gram-negative (Neisseria, Treponema). Resistance: beta-lactamase production, altered PBPs. Derivatives: penicillin G (natural), amoxicillin (aminopenicillin), methicillin (resistant to staphylococcal beta-lactamase), piperacillin (extended spectrum).

\medterm{Pentose Phosphate Pathway}-- Also known as the hexose monophosphate (HMP) shunt or phosphogluconate pathway, an alternative cytoplasmic glucose catabolic pathway that diverges from glycolysis at glucose-6-phosphate (G6P) without directly consuming or generating ATP. The pathway fulfills two distinct and vital physiological functions: generating NADPH (reducing equivalents) and producing ribose-5-phosphate (for nucleotide and nucleic acid biosynthesis). The pathway operates in two functional phases: (1) Irreversible Oxidative Phase: Glucose-6-phosphate is oxidized by the rate-limiting, committed enzyme glucose-6-phosphate dehydrogenase (G6PD) to 6-phosphoglucono-$\delta$-lactone, which is hydrolyzed to 6-phosphogluconate and oxidatively decarboxylated by 6-phosphogluconate dehydrogenase to yield ribulose-5-phosphate, generating 2 molecules of NADPH and 1 molecule of $CO_2$ per G6P. G6PD is allosterically regulated by the $NADP^+/\text{NADPH}$ ratio, being strongly inhibited by NADPH; (2) Reversible Non-Oxidative Phase: Catalyzed by transketolase (an enzyme requiring thiamine pyrophosphate [TPP, vitamin B1] as a cofactor) and transaldolase, which interconvert 5-carbon pentose phosphates (ribose-5-phosphate and xylulose-5-phosphate) with glycolytic intermediates (fructose-6-phosphate and glyceraldehyde-3-phosphate). This reversibility allows cells to produce ribose-5-phosphate without generating NADPH when nucleotide synthesis is needed, or return excess pentoses to glycolysis. NADPH generated by the oxidative phase is essential for: reductive biosynthesis of fatty acids and cholesterol (liver, lactating mammary glands, adipose), steroid hormone synthesis (adrenal cortex, gonads), maintenance of reduced glutathione to prevent erythrocyte oxidative hemolysis, and the generation of microbicidal reactive oxygen species by phagocyte NADPH oxidase during respiratory burst.

\medterm{Peptide Bond}-- A stable covalent amide linkage formed between the $\alpha$-carboxyl group (-COOH) of one amino acid and the $\alpha$-amino group (-NH$_2$) of an adjacent amino acid, releasing a molecule of water through a condensation (dehydration) reaction. In biological systems, peptide bond synthesis occurs within the peptidyl transferase center of the large ribosomal subunit, an RNA-catalyzed (ribozyme) reaction where the nucleophilic $\alpha$-amino group of an aminoacyl-tRNA in the ribosomal A-site attacks the ester carbonyl carbon of the peptidyl-tRNA in the P-site. The resulting peptide bond (-CO-NH-) possesses approximately 40\% partial double-bond character due to resonance delocalization of the nitrogen atom's lone electron pair into the carbonyl $\pi$-system. This partial double-bond character restricts rotation about the C-N bond, causing the six atoms of the peptide group ($C_\alpha$, $C$, $O$, $N$, $H$, and adjacent $C_\alpha$) to lie within a rigid, coplanar conformation. Free conformational rotation within a polypeptide backbone is restricted strictly to the single bonds flanking the $\alpha$-carbon: the phi ($\phi$) angle between $N$ and $C_\alpha$, and the psi ($\psi$) angle between $C_\alpha$ and the carbonyl carbon, whose sterically permitted values are mapped by the Ramachandran plot to determine secondary protein structures ($\alpha$-helices, $\beta$-pleated sheets). Due to steric hindrance between side chains, the trans configuration of the peptide bond is strongly favored over cis (except in proline residues, where ~10\% adopt cis conformation). In vitro, peptide bonds are exceptionally stable at physiological pH (uncatalyzed half-life $>500$ years), requiring specialized proteolytic enzymes (peptidases, proteases) for cellular hydrolysis.

\medterm{pH}-- A logarithmic scale quantifying the concentration of free, unbuffered hydrogen ions in an aqueous solution: $\text{pH} = -\log_{10}[H^+]$, or $[H^+] = 10^{-\text{pH}}$. In human arterial blood, extracellular pH is maintained within an exceptionally narrow physiological range of 7.35 to 7.45 (corresponding to $[H^+]$ of 35 to 45 nmol/L). Acidemia is defined as an arterial $\text{pH} < 7.35$, whereas alkalemia is defined as $\text{pH} > 7.45$; severe deviations outside 6.80 to 7.80 are generally fatal due to catastrophic protein denaturation, enzyme inactivation, and disruption of cardiac and neuromuscular action potentials. Physiological pH is regulated dynamically across three sequential lines of defense: (1) Chemical buffer systems (instantaneous): primarily the extracellular bicarbonate-carbonic acid system ($H^+ + HCO_3^- \rightleftharpoons H_2CO_3 \rightleftharpoons CO_2 + H_2O$, described by the Henderson-Hasselbalch equation $\text{pH} = 6.1 + \log_{10}([HCO_3^-] / (0.03 \times P_{CO_2}))$), supported by intracellular proteins (hemoglobin, histidine imidazole groups) and bone phosphate; (2) Respiratory compensation (minutes to hours): central and peripheral chemoreceptors alter alveolar ventilation to retain or blow off volatile $CO_2$; and (3) Renal compensation (hours to days): proximal and distal renal tubules excrete fixed nonvolatile acids (as titratable acid and ammonium, $NH_4^+$) and reclaim or synthesize new bicarbonate.

\medterm{Phenylketonuria}-- The archetypal autosomal recessive inborn error of amino acid metabolism (prevalence $\approx 1:10,000$ live births), caused by pathogenic mutations in the PAH gene encoding phenylalanine hydroxylase (PAH), or rarely ($<2\%$) by defects in the biosynthesis or regeneration of its essential cofactor, tetrahydrobiopterin ($BH_4$, via dihydropteridine reductase). Under physiological conditions, PAH converts the essential amino acid L-phenylalanine into L-tyrosine in hepatocytes, requiring molecular oxygen and $BH_4$. Enzymatic deficiency blocks this conversion, making tyrosine a nutritionally essential amino acid and causing massive accumulation of phenylalanine in plasma, tissues, and the central nervous system. Excess phenylalanine is diverted into alternative minor catabolic pathways, producing phenylketones (phenylpyruvate, phenyllactate, and phenylacetate), which are excreted in large quantities in urine, imparting a characteristic musty or 'mousy' body odor. Untreated PKU results in severe irreversible intellectual disability, microcephaly, seizures, eczema, behavioral disturbances, and prominent hypopigmentation (fair skin, blonde hair, and blue eyes due to phenylalanine inhibition of tyrosinase, the rate-limiting enzyme in melanin synthesis). Universal newborn screening via tandem mass spectrometry allows immediate institution of a strict, lifelong phenylalanine-restricted diet supplemented with specialized amino acid medical foods and sapropterin ($BH_4$ analogue) for responsive variants. In pregnant women with PKU, strict maternal phenylalanine control is mandatory to prevent Maternal PKU Syndrome (microcephaly, congenital heart defects, and fetal intellectual disability).

\medterm{Phosphate}-- An essential trivalent polyatomic inorganic anion ($PO_4^{3-}$, normal serum range: 2.5--4.5 mg/dL or 0.8--1.45 mmol/L) that exists in extracellular fluid predominantly as monohydrogen phosphate ($HPO_4^{2-}$) and dihydrogen phosphate ($H_2PO_4^-$) in a 4:1 physiological ratio at pH 7.4. Phosphate plays an indispensable structural, biochemical, and regulatory role in human physiology: (1) Bioenergetics: forms the high-energy phosphoanhydride bonds of adenosine triphosphate (ATP), creatine phosphate, and 1,3-bisphosphoglycerate, storing and transferring free energy; (2) Macromolecular Architecture: forms the phosphodiester backbone of DNA and RNA, and constitutes the polar headgroup of membrane phospholipids (phosphatidylcholine, phosphatidylethanolamine); (3) Signal Transduction: reversible protein phosphorylation and dephosphorylation catalyzed by kinases and phosphatases serves as the primary on/off switch for enzyme activity; (4) Skeletal Mineralization: complexed with calcium to form hydroxyapatite crystals ($Ca_{10}(PO_4)_6(OH)_2$) in bone and dental enamel; (5) Oxygen Delivery: forms 2,3-bisphosphoglycerate (2,3-BPG) in erythrocytes, modulating hemoglobin oxygen affinity; and (6) Acid-Base Buffering: acts as the primary 'titratable acid' buffer in renal tubular urine ($H^+ + HPO_4^{2-} \to H_2PO_4^-$), facilitating net hydrogen ion excretion. Phosphate homeostasis is regulated by parathyroid hormone (PTH, which promotes renal phosphate wasting by downregulating NaPi-IIa cotransporters), fibroblast growth factor 23 (FGF23, phosphaturic), and calcitriol (increases intestinal absorption).

\medterm{Phosphofructokinase-1}-- The primary committed, rate-limiting, and most extensively regulated enzyme of glycolysis, catalyzing the irreversible, ATP-dependent phosphorylation of D-fructose-6-phosphate (F6P) to D-fructose-1,6-bisphosphate (F-1,6-BP). Located in the cytosol as a homotetramer, PFK-1 is subjected to complex allosteric regulation reflecting cellular energy charge and hormonal signals: (1) Allosteric inhibitors: high physiological concentrations of ATP (which binds to a distinct low-affinity regulatory site, lowering affinity for F6P) and citrate (the product of citrate synthase, signaling an abundant supply of TCA cycle intermediates); (2) Allosteric activators: high concentrations of AMP (which relieves ATP-mediated inhibition) and, most potently, fructose-2,6-bisphosphate (F-2,6-BP). F-2,6-BP is produced in the liver by the bifunctional enzyme phosphofructokinase-2/fructose-2,6-bisphosphatase (PFK-2/FBPase-2). Insulin promotes dephosphorylation of the bifunctional enzyme, activating PFK-2 and elevating F-2,6-BP to stimulate glycolysis; glucagon, via cAMP and protein kinase A (PKA), phosphorylates the enzyme, activating FBPase-2 and depleting F-2,6-BP, thereby shutting down PFK-1 and diverting carbon into gluconeogenesis.

\medterm{Phospholipid}-- Lipid containing phosphate group, major component of cell membranes. Structure: glycerol backbone, two fatty acid tails (hydrophobic), phosphate group with polar head group (hydrophilic). Forms bilayers spontaneously in aqueous environment. Types: phosphatidylcholine (lecithin), phosphatidylethanolamine, phosphatidylserine, phosphatidylinositol. Membrane fluidity influenced by fatty acid saturation and cholesterol content.

\medterm{Polycythemia Vera}-- A clonal, BCR-ABL1-negative chronic myeloproliferative neoplasm characterized by the autonomous, cytokine-independent overproduction of morphologically normal red blood cells, frequently accompanied by leukocytosis and thrombocytosis. Greater than 95\% of cases are driven by a somatic gain-of-function missense mutation in exon 14 of the JAK2 gene on chromosome 9, causing a valine-to-phenylalanine substitution at codon 617 (JAK2 V617F) in the autoinhibitory pseudokinase (JH2) domain. This mutation renders the Janus kinase 2 tyrosine kinase constitutively active, driving persistent downstream phosphorylation of STAT5, PI3K-Akt, and MAPK signaling pathways even in the complete absence of erythropoietin (EPO). Laboratory testing reveals an elevated red cell mass (hematocrit $>49\%$ in men, $>48\%$ in women; hemoglobin $>16.5$ g/dL in men, $>16.0$ g/dL in women) accompanied by characteristically suppressed, subnormal serum EPO levels (distinguishing it from secondary polycythemia caused by chronic hypoxia, smoking, or EPO-secreting renal cell carcinomas). The resulting hyperviscosity reduces cerebral blood flow, manifesting with headaches, dizziness, visual disturbances, facial plethoric flushing, and aquagenic pruritus (severe itching following warm showers, mediated by basophil histamine and prostaglandin release). Patients face severe risk of arterial and venous thrombosis, including atypical site thromboses such as Budd-Chiari syndrome (hepatic vein thrombosis) and portal vein thrombosis. First-line therapy consists of therapeutic phlebotomy (maintaining hematocrit $<45\%$), low-dose aspirin, and cytoreductive therapy (hydroxyurea or interferon-alfa) for high-risk patients.

\medterm{Polymerase Chain Reaction (PCR)}-- A revolutionary, Nobel Prize-winning molecular biology technique that enables the rapid, exponential in vitro enzymatic amplification of a specific, targeted segment of double-stranded DNA across millions to billions of copies from minimal initial starting material. The reaction requires five core chemical components: the template DNA containing the target region, a pair of synthetic, single-stranded oligonucleotide primers that flank the target sequence on opposite antiparallel strands, a heat-stable DNA polymerase (most famously Taq polymerase, isolated from the thermophilic bacterium Thermus aquaticus), free deoxynucleotide triphosphates (dNTPs: dATP, dCTP, dGTP, dTTP), and a divalent magnesium ($Mg^{2+}$) buffer cofactor. PCR is carried out in an automated thermal cycler through 25 to 35 repeating cycles, each comprising three discrete temperature-controlled steps: (1) Thermal Denaturation (94--96$^\circ$C): heat breaks the hydrogen bonds between complementary base pairs, yielding two separated single strands of template DNA; (2) Primer Annealing (50--65$^\circ$C): the temperature is lowered to permit oligonucleotide primers to hybridize specifically to their complementary target sequences on the single-stranded templates; and (3) Primer Extension / Elongation (72$^\circ$C): Taq polymerase synthesizes new complementary DNA strands in the 5'-to-3' direction by incorporating dNTPs onto the free 3'-OH ends of the primers. Because the newly synthesized strands serve as templates in subsequent cycles, target DNA accumulates exponentially ($2^n$, where $n$ is cycle count). Advanced variants include Reverse Transcription PCR (RT-PCR, converting viral or cellular RNA to cDNA for SARS-CoV-2 and HIV detection) and Quantitative Real-Time PCR (qPCR, measuring amplicon accumulation via fluorescent probes to quantify viral load and gene expression).

\medterm{Prion}-- A proteinaceous infectious particle completely devoid of nucleic acids that causes fatal, transmissible neurodegenerative conditions known as transmissible spongiform encephalopathies (TSEs). The normal physiological prion protein ($PrP^C$) is a 33--35 kDa cell-surface glycoprotein tethered by a GPI anchor to neuronal membranes, characterized predominantly by an $\alpha$-helical secondary structure ($42\% \ \alpha$-helix, $3\% \ \beta$-sheet) that is soluble in detergents and susceptible to digestion by proteinase K. In prion disease, $PrP^C$ undergoes a profound conformational refolding into the pathogenic, scrapie isoform ($PrP^{Sc}$), which is enriched in $\beta$-sheet architecture ($30\% \ \alpha$-helix, $43\% \ \beta$-sheet). $PrP^{Sc}$ acts as a conformational template, binding normal $PrP^C$ and catalyzing its conversion into additional $PrP^{Sc}$ molecules. The resulting $PrP^{Sc}$ molecules are insoluble in non-ionic detergents and highly resistant to standard proteases, heat, formalin, and ionizing radiation. They aggregate into neurotoxic amyloid fibrils and plaques, provoking vacuolar neuronal degeneration (spongiform change), astrocytic gliosis, and neuronal apoptosis without an inflammatory response. Human prion diseases include Creutzfeldt-Jakob disease (CJD, presenting with rapidly progressive dementia, myoclonus, and periodic sharp wave complexes on EEG), variant CJD (vCJD, linked to bovine spongiform encephalopathy), fatal familial insomnia, and kuru.

\medterm{Protein}-- Polymer of amino acids linked by peptide bonds. Structure: primary (sequence), secondary (alpha helix, beta sheet), tertiary (3D folding), quaternary (multiple subunits). Functions: enzymes, structural, transport, signaling, immune defense, movement. Folding determined by amino acid sequence (Anfinsen's dogma). Chaperones assist folding. Misfolding causes diseases (Alzheimer's, prion diseases). Denaturation by heat, pH, chemicals.

\medterm{Protein C}-- A 62-kDa vitamin K-dependent serine protease zymogen synthesized in the liver that serves as the centerpiece of the body's natural anticoagulant system. Protein C activation occurs on endothelial cell surfaces when thrombin binds to the transmembrane receptor thrombomodulin; this alters thrombin's substrate specificity, transforming it from a procoagulant into a potent activator of protein C (accelerated 20-fold by the endothelial protein C receptor, EPCR). Activated protein C (APC), in concert with its vitamin K-dependent cofactor protein S, proteolytically cleaves and irreversibly inactivates activated factor V (FVa) and activated factor VIII (FVIIIa), dismantling the prothrombinase and tenase complexes and halting further thrombin generation. APC also promotes fibrinolysis by inactivating plasminogen activator inhibitor-1 (PAI-1). Hereditary protein C deficiency (autosomal dominant) predisposes to recurrent venous thromboembolism. Crucially, upon initiating warfarin therapy, protein C levels drop precipitously before procoagulant factors due to its short biological half-life (6--8 hours), creating a temporary hypercoagulable window that can precipitate microvascular thrombosis and warfarin-induced skin necrosis unless bridged with heparin.

\medterm{Protein S}-- A 75-kDa vitamin K-dependent plasma glycoprotein synthesized primarily by hepatocytes, vascular endothelial cells, and megakaryocytes that functions as an indispensable non-enzymatic cofactor for activated protein C (APC). Protein S binds APC on phospholipid surfaces, dramatically enhancing the catalytic cleavage and inactivation of procoagulant factor Va and factor VIIIa. In human plasma, protein S circulates in two equilibrium states: approximately 60\% is bound with high affinity to C4b-binding protein (C4b-BP), a regulatory protein of the complement cascade, rendering it functionally inactive; only the remaining 40\% circulates as uncomplexed 'free' protein S, which is the biologically active anticoagulant form. Because C4b-BP is an acute-phase reactant that rises during systemic inflammation, inflammatory states shift the equilibrium, decreasing free protein S and fostering a prothrombotic state. Hereditary protein S deficiency (autosomal dominant) causes venous thromboembolism (deep vein thrombosis and pulmonary embolism) and confers risk for warfarin-induced skin necrosis and neonatal purpura fulminans in homozygous states.

\medterm{Protein Synthesis}-- Translation of mRNA code into polypeptide chain. Initiation: ribosome assembly at start codon (AUG). Elongation: aminoacyl-tRNAs bring amino acids; peptide bonds form in peptidyl transferase center; ribosome translocates. Termination: stop codon (UAA, UAG, UGA) recognized by release factors. Post-translational modifications: phosphorylation, glycosylation, acetylation, ubiquitination. Energy: GTP per amino acid added.

\medterm{Prothrombin Time}-- A clinical screening coagulation test that evaluates the extrinsic and common pathways of the coagulation cascade (specifically factors VII, X, V, II [prothrombin], and I [fibrinogen]). Performed by adding recombinant tissue factor, phospholipids, and calcium to citrated platelet-poor plasma and measuring the seconds elapsed until fibrin clot formation (reference range: 11--13.5 seconds). To standardize results across laboratories utilizing tissue factor reagents of varying sensitivity, the prothrombin time is converted to the International Normalized Ratio: $\text{INR} = (\text{PT}_{\text{patient}} / \text{PT}_{\text{control}})^{\text{ISI}}$, where ISI is the International Sensitivity Index. The PT/INR is the primary test used to monitor oral vitamin K antagonist therapy (warfarin), with a standard target INR of 2.0 to 3.0 (or 2.5 to 3.5 for mechanical mitral heart valves). Because factor VII has the shortest biological half-life (4--6 hours) of all coagulation factors, the PT/INR is the most sensitive and earliest laboratory indicator of acute hepatic synthetic failure (e.g., fulminant hepatitis) and acute vitamin K deficiency.

\medterm{Purines}-- Nitrogenous bases containing a double-ring structure (a six-membered ring fused to a five-membered imidazole ring), which serve as essential components of nucleotides in DNA and RNA. The primary purines are adenine (A) and guanine (G). They are synthesized de novo starting from ribose-5-phosphate or recycled via salvage pathways (catalyzed by HGPRT and APRT). Defective purine metabolism leads to gout (excess uric acid accumulation) and Lesch-Nyhan syndrome (HGPRT deficiency causing severe neurological impairment and self-mutilation).

\medterm{Pyridoxine (Vitamin B6)}-- A collective term for three natural pyridine vitamer precursors (pyridoxine, pyridoxal, and pyridoxamine) that are phosphorylated by pyridoxal kinase into the active coenzyme pyridoxal 5'-phosphate (PLP). PLP acts as a versatile coenzyme for $>140$ distinct enzymatic reactions, predominantly in amino acid metabolism: (1) Transamination: PLP forms a Schiff base intermediate with amino acid $\alpha$-amino groups, required by ALT and AST; (2) Decarboxylation: synthesis of neurotransmitters including GABA (from glutamate via GAD), dopamine (from DOPA), serotonin (from 5-HTP), and histamine (from histidine); (3) Heme Synthesis: the initial, rate-limiting step of porphyrin synthesis catalyzed by $\delta$-aminolevulinic acid synthase (ALAS), localized to the mitochondrial matrix; (4) Transsulfuration: conversion of homocysteine to cystathionine by cystathionine $\beta$-synthase; and (5) Glycogenolysis: cofactor for glycogen phosphorylase in muscle and liver. Nutritional deficiency is seen in chronic alcoholism and is induced by medications that form inactive hydrazone complexes with PLP (isoniazid, penicillamine). Deficiency manifests as peripheral neuropathy, convulsions/seizures (due to impaired GABA synthesis), microcytic hypochromic sideroblastic anemia (due to impaired heme synthesis, accumulating iron in ringed sideroblasts), stomatitis, and glossitis. Conversely, chronic massive megadoses of pyridoxine ($>200--500$ mg/day) cause severe, irreversible sensory axonal neuropathy affecting posterior columns.

\medterm{Pyrimidines}-- Nitrogenous bases containing a single six-membered heterocyclic ring, serving as structural building blocks for nucleic acids. The primary pyrimidines are cytosine (C), thymine (T, present in DNA), and uracil (U, present in RNA). Unlike purines, the pyrimidine ring is synthesized first (using carbamoyl phosphate and aspartate) and subsequently attached to ribose-5-phosphate. Disorders of pyrimidine metabolism include hereditary orotic aciduria, which is characterized by megaloblastic anemia, developmental delay, and excessive orotic acid excretion.

\medterm{Pyruvate}-- The pivotal three-carbon $\alpha$-ketoacid end-product of glycolysis, situated at the central crossroads of carbohydrate, lipid, and amino acid metabolism. In the cytoplasm, 1 molecule of glucose yields 2 molecules of pyruvate via 10 enzymatic reactions, generating a net of 2 ATP and 2 NADH. Depending on cellular oxygen availability, mitochondrial function, and energetic demand, pyruvate has four primary metabolic fates: (1) Oxidative Decarboxylation to Acetyl-CoA: Catalyzed by the mitochondrial pyruvate dehydrogenase (PDH) complex under aerobic conditions, generating NADH and $CO_2$, feeding carbon into the TCA cycle or lipogenesis; (2) Reduction to L-Lactate: Catalyzed by cytoplasmic lactate dehydrogenase (LDH) under anaerobic conditions (or in erythrocytes), regenerating $NAD^+$ to permit continued glycolytic flux; (3) Carboxylation to Oxaloacetate: Catalyzed by mitochondrial pyruvate carboxylase (requiring biotin and ATP), providing an essential anaplerotic replenishment of TCA cycle intermediates and initiating gluconeogenesis; and (4) Transamination to L-Alanine: Catalyzed by alanine aminotransferase (ALT, requiring pyridoxal phosphate [PLP, vitamin B6]), participating in the Cahill (glucose-alanine) cycle to safely transport muscle amino groups to the liver.

\medterm{Pyruvate Dehydrogenase}-- A massive multienzyme complex located in the mitochondrial matrix that catalyzes the irreversible oxidative decarboxylation of pyruvate to acetyl-CoA, releasing $CO_2$ and generating NADH. This critical reaction links cytosolic glycolysis to the mitochondrial citric acid cycle and lipogenesis. The complex consists of three core catalytic enzymes: pyruvate dehydrogenase ($E_1$), dihydrolipoyl transacetylase ($E_2$), and dihydrolipoyl dehydrogenase ($E_3$), which require five distinct cofactors: thiamine pyrophosphate (TPP, vitamin B1), lipoic acid, coenzyme A (pantothenic acid, B5), flavin adenine dinucleotide (FAD, B2), and nicotinamide adenine dinucleotide (NAD+, B3). Regulation occurs via allosteric feedback and covalent modification: pyruvate dehydrogenase kinase (PDK) phosphorylates and inactivates $E_1$ in response to high ATP, NADH, and acetyl-CoA; pyruvate dehydrogenase phosphatase (PDP) dephosphorylates and activates $E_1$ in response to $Ca^{2+}$ (signaling contraction) and insulin. Pyruvate dehydrogenase deficiency (most commonly X-linked $E_1$ subunit mutations) impairs aerobic carbohydrate oxidation, causing microcephaly, developmental delay, and severe lactic acidosis; treatment includes a ketogenic diet with high fat and minimal carbohydrates to generate acetyl-CoA via beta-oxidation.

\medterm{Reactive Oxygen Species}-- Chemically reactive chemical species derived from molecular oxygen ($O_2$) that contain oxygen atoms with unpaired valence electrons or peroxide bonds, generated as normal metabolic byproducts of mitochondrial aerobic respiration and cellular enzymatic reactions. The primary endogenous source is electron leakage from Complex I and Complex III of the mitochondrial respiratory chain, where single electrons reduce $O_2$ to the superoxide anion radical ($O_2^{\bullet-}$). Superoxide is dismutated by mitochondrial manganese superoxide dismutase (Mn-SOD/SOD2) or cytosolic copper-zinc superoxide dismutase (Cu/Zn-SOD/SOD1) into hydrogen peroxide ($H_2O_2$). In the presence of reduced transition metals ($Fe^{2+}$ or $Cu^+$), hydrogen peroxide undergoes the non-enzymatic Fenton reaction to generate the hydroxyl radical ($^{\bullet}OH$), the most destructive and non-selective free radical in biological systems ($Fe^{2+} + H_2O_2 \to Fe^{3+} + ^{\bullet}OH + OH^-$). In phagocytes (neutrophils and macrophages), NADPH oxidase intentionally generates large volumes of superoxide during the respiratory burst, which myeloperoxidase (MPO) converts with chloride ions into hypochlorous acid ($HOCl$, bleach) to destroy engulfed microbes. When ROS generation exceeds cellular antioxidant defenses (glutathione peroxidase, catalase, vitamins C and E), oxidative stress ensues, causing lipid peroxidation of polyunsaturated fatty acids (forming malondialdehyde), protein carbonylation, and oxidative DNA base damage (notably 8-oxoguanine).

\medterm{Riboflavin (Vitamin B2)}-- A water-soluble vitamin composed of a dimethylisoalloxazine ring linked to a ribitol side chain. Riboflavin is converted sequentially in the cytoplasm by flavokinase to flavin mononucleotide (FMN) and by FAD synthetase to flavin adenine dinucleotide (FAD), using ATP. FAD and FMN serve as tightly bound or covalently linked prosthetic groups for numerous oxidoreductases in cellular bioenergetics: (1) Succinate dehydrogenase (Complex II of the electron transport chain and TCA cycle, reducing $FAD \to FADH_2$); (2) Acyl-CoA dehydrogenase (the first step of fatty acid beta-oxidation); (3) Pyruvate dehydrogenase and $\alpha$-ketoglutarate dehydrogenase (as the prosthetic group of the dihydrolipoyl dehydrogenase $E_3$ subunit); and (4) Glutathione reductase (which transfers electrons from NADPH to FAD to reduce oxidized glutathione, maintaining erythrocytic antioxidant defense). Riboflavin deficiency (ariboflavinosis) is typically seen in chronic malnutrition, alcoholism, or anorexia nervosa, manifesting clinically with the classic '2 C's': Cheilosis/angular stomatitis (fissuring and inflammation of the angles of the mouth), Corneal vascularization, glossitis (magenta tongue), and seborrheic dermatitis. Erythrocyte glutathione reductase activity coefficient (EGRAC) before and after FAD stimulation serves as the definitive diagnostic assay.

\medterm{Ribosomal RNA (rRNA)}-- RNA component of ribosomes, catalyzing peptide bond formation (peptidyl transferase activity). Eukaryotic ribosome: 60S (28S, 5.8S, 5S rRNAs + proteins) + 40S (18S rRNA + proteins) = 80S. Prokaryotic: 50S (23S, 5S rRNAs) + 30S (16S rRNA) = 70S. rRNA genes are highly conserved, used in phylogenetics.

\medterm{Ribozyme}-- An RNA molecule possessing intrinsic catalytic activity capable of accelerating specific biochemical reactions in the absence of protein enzymes, demonstrating that RNA can serve simultaneously as both genetic information carrier and catalytic macromolecule (supporting the 'RNA World' hypothesis of early evolution). Ribozymes fold into precise three-dimensional tertiary architectures that stabilize transition states through acid-base catalysis, conformational strain, and coordination of divalent metal ion cofactors (such as $Mg^{2+}$). The most prominent biological example in all living cells is the peptidyl transferase center of the large ribosomal subunit (28S rRNA in eukaryotes, 23S rRNA in prokaryotes), which catalyzes peptide bond formation during protein translation purely through RNA-mediated catalysis without direct polypeptide participation. Other essential ribozymes include: (1) The catalytic RNA component of ribonuclease P (RNase P), a ribonucleoprotein that cleaves the 5' leader sequence of precursor tRNAs; (2) Self-splicing group I and group II introns; and (3) Small self-cleaving ribozymes (hammerhead, hairpin, hepatitis delta virus ribozymes) that perform site-specific cleavage during rolling-circle viral genome replication.

\medterm{Sepsis}-- A life-threatening medical emergency defined biochemically and clinically as organ dysfunction caused by a dysregulated, aberrant host response to infection (Sequential Organ Failure Assessment [SOFA] score increase $\ge 2$). Sepsis is initiated when widespread pathogen invasion triggers an overwhelming, uncontrolled release of pro-inflammatory cytokines ('cytokine storm': TNF-$\alpha$, IL-1$\beta$, IL-6, and HMGB1) and complement activation products (C3a, C5a). This systemic inflammatory cascade inflicts widespread microvascular endothelial injury, leading to systemic degradation of the protective endothelial glycocalyx, microvascular hyperpermeability, and massive capillary fluid extravasation into interstitial tissues. Concurrently, systemic expression of tissue factor on monocytes and denuded endothelium triggers extensive disseminated intravascular coagulation (DIC), while loss of endogenous antithrombin and activated protein C impairs anticoagulation, depositing microvascular fibrin thrombi that obstruct capillary perfusion. Cellular hypoxia, mitochondrial dysfunction ('cytopathic hypoxia'), and uncoupled oxidative phosphorylation force cells into anaerobic glycolysis, generating profound hyperlactatemia ($>2$ mmol/L). In Septic Shock, persistent hypotension requires vasopressors to maintain a mean arterial pressure (MAP) $\ge 65$ mmHg despite adequate fluid resuscitation, accompanied by a serum lactate $>2$ mmol/L, reflecting severe cellular metabolic and circulatory collapse with hospital mortality exceeding 40\%.

\medterm{Sialuria}-- An exceedingly rare autosomal dominant inborn error of carbohydrate metabolism caused by heterozygous missense mutations in the allosteric regulatory domain of the GNE gene (encoding UDP-N-acetylglucosamine 2-epimerase / N-acetylmannosamine kinase), the bifunctional rate-limiting enzyme in the biosynthesis of CMP-sialic acid (CMP-Neu5Ac). Under physiological conditions, CMP-sialic acid acts as an allosteric feedback inhibitor that binds the epimerase domain of GNE to shut off further sialic acid production when cellular requirements are met. The pathogenic mutations disrupt this allosteric feedback pocket, rendering the enzyme completely unresponsive to feedback inhibition. Consequently, cells synthesize massive excess quantities of free, uncomplexed N-acetylneuraminic acid (sialic acid), which saturates intracellular utilization and spills into the circulation and urine (French-type sialuria, characterized by massive urinary excretion of grams of free sialic acid per day). Clinical features include variable developmental delay, mild facial coarsening, hepatosplenomegaly, and recurrent upper respiratory infections, often showing clinical improvement with age.

\medterm{Sickle Cell Trait}-- The heterozygous carrier state for sickle cell hemoglobin ($HbAS$), resulting from inheritance of one normal adult $\beta$-globin allele ($HbA$) and one sickle $\beta$-globin allele ($HbS$, characterized by a point mutation in codon 6 of the $HBB$ gene on chromosome 11, substituting valine for glutamic acid: Glu6Val). The hydrophobic valine residue creates a sticky patch on the deoxyhemoglobin molecule. Because red blood cells in individuals with sickle cell trait typically contain $>50--60\%$ normal $HbA$ and only $35--40\%$ $HbS$ (due to higher translational efficiency and affinity of $\alpha$-chains for $\beta^A$ over $\beta^S$), the intracellular concentration of deoxy-HbS remains below the critical threshold required for hemoglobin polymer nucleation under normal physiological conditions. Consequently, individuals with sickle cell trait are clinically asymptomatic, possess normal life expectancies, and exhibit no baseline anemia, hemolysis, or peripheral blood smear abnormalities. Sickle cell trait confers significant evolutionary protection against severe falciparum malaria. However, under extreme non-physiological conditions of severe hypoxia, profound lactic acidosis, dehydration, or extreme physical exertion at high altitude, erythrocytes can sickle, rarely predisposing to renal medullary sickling (causing microscopic or macroscopic painless hematuria, loss of urinary concentrating ability / isosthenuria, and renal papillary necrosis), splenic infarction at high altitudes, and an increased risk of renal medullary carcinoma.

\medterm{Sitosterolemia}-- A rare autosomal recessive lipid storage disorder (also termed phytosterolemia) caused by loss-of-function mutations in either the ABCG5 or ABCG8 gene, which encode the heterodimeric ATP-binding cassette transporter sterolin (ABCG5/ABCG8) in the apical membrane of intestinal enterocytes and the canalicular membrane of hepatocytes. Under normal physiology, sterolin acts as an active efflux pump that excretes dietary plant sterols (sitosterol, campesterol, stigmasterol) out of enterocytes back into the gut lumen (limiting plant sterol absorption to $<5\%$) and excretes hepatic sterols into bile. Defective ABCG5/ABCG8 leads to hyperabsorption of plant sterols and failure of biliary sterol excretion, causing a 30- to 100-fold elevation in plasma plant sterol concentrations. Clinical features include extensive tendon and tuberous xanthomas, premature severe coronary artery atherosclerosis, recurrent acute hemolysis, and macrothrombocytopenia with stomatocytosis (plant sterols incorporate into erythrocyte and platelet membranes, disrupting membrane fluidity and mechanical stability). Routine clinical cholesterol assays cannot distinguish plant sterols from cholesterol; diagnosis requires gas chromatography-mass spectrometry. Treatment includes a diet strictly restricted in plant sterols (avoiding vegetable oils, nuts, avocados, margarine) and the cholesterol absorption inhibitor ezetimibe (which inhibits NPC1L1, drastically lowering plant sterol levels).

\medterm{Succinate Dehydrogenase}-- A unique heterotetrameric iron-sulfur flavoprotein that participates simultaneously in two foundational bioenergetic pathways: the citric acid cycle (Krebs cycle) and the mitochondrial electron transport chain (as Complex II). Unlike all other TCA cycle enzymes which are soluble in the matrix, succinate dehydrogenase is anchored directly into the inner mitochondrial membrane. It consists of two hydrophilic catalytic subunits projecting into the matrix (SDHA, containing covalently bound FAD, and SDHB, containing three iron-sulfur clusters $[2Fe-2S], [4Fe-4S], [3Fe-4S]$) and two hydrophobic transmembrane anchoring subunits (SDHC and SDHD, containing a heme $b$ group and ubiquinone binding sites). The enzyme oxidizes succinate to fumarate while reducing bound $FAD \to FADH_2$; the electrons are then transferred through the iron-sulfur centers to reduce ubiquinone (coenzyme Q) to ubiquinol ($QH_2$). Because electron transfer through Complex II does not traverse the inner membrane, it does NOT pump protons and does not directly contribute to the proton gradient. Inactivating germline mutations in SDH subunits lead to accumulation of succinate (an oncometabolite inhibiting $\alpha$-ketoglutarate-dependent dioxygenases, stabilizing HIF-1$\alpha$) and predispose to familial paragangliomas and pheochromocytomas.

\medterm{Telomere}-- Specialized nucleoprotein structures situated at the physical termini of linear eukaryotic chromosomes that preserve chromosomal stability, prevent end-to-end chromosomal fusion, and protect coding genomic sequences from degradation during cell division. Human telomeric DNA consists of thousands of tandem repeats of the non-coding hexanucleotide sequence 5'-TTAGGG-3', spanning approximately 10 to 15 kilobases at birth. The extreme 3' terminus terminates in a single-stranded overhang of 100--300 nucleotides that invades the proximal double-stranded repeat tract, forming a loop structure called the T-loop (telomeric loop). The T-loop is stabilized by a dedicated six-protein complex termed shelterin (TRF1, TRF2, RAP1, TIN2, TPP1, POT1), which caps the chromosome end and hides it from being recognized as an accidental double-strand break by the ATM/ATR DNA damage response machinery. Because standard DNA polymerases require an RNA primer and can only synthesize DNA in the 5'-to-3' direction, the removal of the terminal RNA primer leaves a 3' gap that cannot be replicated (the 'end-replication problem'). Consequently, telomeres shorten by 50--100 base pairs with each cycle of mitotic division, eventually reaching the Hayflick limit, which activates p53/p21 to trigger irreversible replicative cellular senescence or apoptosis. Germ cells, stem cells, and $>85\%$ of human cancers overcome this limit by upregulating telomerase, a ribonucleoprotein reverse transcriptase (TERT) that utilizes an internal RNA template (TERC) to lengthen telomeres.

\medterm{Thalassemia}-- A group of inherited autosomal recessive hemoglobinopathies characterized by absent or reduced synthesis of one or more of the globin polypeptide chains ($\alpha$ or $\beta$) that constitute adult hemoglobin $A$ ($\alpha_2\beta_2$). The resulting imbalance between $\alpha$- and $\beta$-globin chains causes precipitation of the excess, unpaired globin chains within erythroid precursors, driving ineffective intramedullary erythropoiesis and severe peripheral extravascular hemolysis: (1) $\beta$-Thalassemia: caused primarily by point mutations (splicing or promoter defects) in the HBB gene on chromosome 11, resulting in reduced ($\beta^+$) or absent ($\beta^0$) $\beta$-chain synthesis. $\beta$-Thalassemia minor (heterozygous) presents with mild microcytic anemia ($MCV < 75$ fL) with elevated $HbA_2$ ($>3.5\%$) on hemoglobin electrophoresis; $\beta$-Thalassemia major (Cooley anemia, homozygous $\beta^0/\beta^0$) causes profound transfusion-dependent microcytic anemia manifesting in infancy (as fetal $HbF$ $[\alpha_2\gamma_2]$ wanes), marked hepatosplenomegaly, skeletal deformities ('crew-cut' skull, chipmunk facies due to massive marrow expansion), and fatal secondary hemochromatosis from chronic transfusions; (2) $\alpha$-Thalassemia: caused primarily by gene deletions across the four $\alpha$-globin alleles on chromosome 16: 1-gene deletion (silent carrier, normal), 2-gene deletion ($\alpha$-thalassemia trait, mild microcytosis), 3-gene deletion (Hemoglobin H disease, excess $\beta$-chains form $\beta_4$ tetramers causing chronic hemolytic anemia), and 4-gene deletion (Hemoglobin Barts, excess $\gamma$-chains form $\gamma_4$ tetramers with extreme oxygen affinity, causing fatal hydrops fetalis in utero).

\medterm{Thiamine (Vitamin B1)}-- A water-soluble B-complex vitamin that is phosphorylated by thiamine pyrophosphokinase to its active coenzyme form, thiamine pyrophosphate (TPP). TPP is an obligatory cofactor for four essential multienzyme complexes involved in carbohydrate and branched-chain amino acid catabolism: (1) Pyruvate dehydrogenase (linking glycolysis to the citric acid cycle); (2) $\alpha$-Ketoglutarate dehydrogenase (in the citric acid cycle); (3) Transketolase (in the non-oxidative branch of the pentose phosphate pathway); and (4) Branched-chain $\alpha$-ketoacid dehydrogenase (in leucine, isoleucine, and valine degradation). Thiamine deficiency impairs aerobic glucose oxidation and ATP generation, with high-energy tissues (brain and myocardium) being most vulnerable. Severe deficiency manifests as: Wet Beriberi, characterized by high-output congestive heart failure, peripheral vasodilation, dyspnea, and brawny edema; Dry Beriberi, presenting with symmetrical peripheral polyneuropathy, sensory loss, muscle wasting, and diminished deep tendon reflexes; and Wernicke Encephalopathy, an acute neuropsychiatric emergency seen frequently in chronic alcohol use disorder, characterized by the classic triad of confusion/encephalopathy, ophthalmoplegia/nystagmus, and gait ataxia. If untreated, it progresses to irreversible Korsakoff Psychosis, marked by anterograde and retrograde amnesia and confabulation due to hemorrhagic necrosis of the mammillary bodies and medial thalamus. Erythrocyte transketolase activity before and after in vitro TPP addition confirms deficiency. In malnourished patients, intravenous thiamine MUST always be administered before glucose infusions to prevent precipitating acute Wernicke encephalopathy.

\medterm{Thyroid Hormone}-- Iodinated amino acid-derived hormones synthesized and secreted by the follicular cells of the thyroid gland, comprising prohormone thyroxine ($T_4$, 3,5,3',5'-tetraiodothyronine, accounting for ~80\% of secretion) and biologically active triiodothyronine ($T_3$, 3,5,3'-triiodothyronine, accounting for ~20\% of secretion, with the remainder generated by peripheral tissue outer-ring deiodination of $T_4$ via 5'-deiodinase). Synthesis requires active iodide trapping via the sodium-iodide symporter (NIS), oxidation of iodide and organification into tyrosyl residues of thyroglobulin by thyroid peroxidase (TPO), followed by coupling of monoiodotyrosine (MIT) and diiodotyrosine (DIT). Lipophilic free $T_3$ enters target cells and binds nuclear thyroid hormone receptors (TR-$\alpha$ and TR-$\beta$), forming heterodimers with retinoid X receptors (RXR) that bind thyroid response elements (TREs) to regulate transcription. Thyroid hormone is the master determinant of basal metabolic rate (BMR): it upregulates $Na^+/K^+$-ATPase expression across tissues, increases mitochondrial uncoupling protein-1 (UCP-1 / thermogenin) to stimulate calorigenesis and oxygen consumption, stimulates glycogenolysis, gluconeogenesis, and lipolysis, and upregulates cardiac $\beta_1$-adrenergic receptor density (increasing heart rate and contractility). Thyroid hormone is absolutely required for fetal and neonatal neurogenesis and epiphyseal bone maturation; untreated congenital hypothyroidism causes irreversible intellectual disability and stunted skeletal growth (cretinism). In adults, hypothyroidism causes weight gain, fatigue, cold intolerance, bradycardia, and myxedema; hyperthyroidism causes weight loss, heat intolerance, tremors, tachycardia, and atrial fibrillation.

\medterm{Transcription}-- Synthesis of RNA from DNA template by RNA polymerase. Eukaryotic: RNA Pol II transcribes mRNA, rRNA, snRNA; RNA Pol I transcribes rRNA; RNA Pol III transcribes tRNA, 5S rRNA. Steps: initiation (promoter recognition, transcription factors), elongation (5' to 3'), termination. Promoter elements: TATA box, CAAT box, GC box. Post-transcriptional modifications: 5' cap, 3' poly-A tail, splicing.

\medterm{Translation}-- Decoding of mRNA sequence into polypeptide chain by ribosomes. Occurs in cytoplasm. tRNAs carry specific amino acids and recognize codons via anticodons. Steps: initiation, elongation, termination. Energy: GTP hydrolysis. Antibiotics targeting translation: tetracyclines (30S), chloramphenicol (50S), erythromycin (50S), streptomycin (30S).

\medterm{Trimethylaminuria}-- An autosomal recessive metabolic disorder (commonly known as 'fish odor syndrome') caused by loss-of-function mutations in the FMO3 gene on chromosome 1, which encodes flavin-containing monooxygenase 3. FMO3 is a hepatic microsomal enzyme responsible for the NADPH- and oxygen-dependent N-oxygenation of dietary trimethylamine (TMA) into odorless trimethylamine N-oxide (TMAO). Trimethylamine is a volatile, tertiary aliphatic amine generated in large quantities by colonic bacterial fermentation of dietary precursors rich in choline, lecithin, and L-carnitine (abundant in egg yolks, saltwater fish, liver, legumes, and red meat). In FMO3 deficiency, unmetabolized trimethylamine accumulates in the bloodstream and is excreted unchanged in sweat, urine, saliva, and expired breath, imparting a pungent, sickening odor resembling rotting fish. Although physically benign without neurological, hepatic, or renal toxicity, trimethylaminuria causes profound psychosocial disability, depression, anxiety, social isolation, and suicide risk. Diagnosis is confirmed by demonstrating an elevated ratio of free TMA to total TMA in urine, confirmed by FMO3 genetic sequencing. Management focuses on dietary avoidance of choline- and carnitine-rich foods, short courses of oral non-absorbable antibiotics (metronidazole, neomycin) to suppress intestinal TMA-producing bacteria, riboflavin (vitamin B2) supplementation (an FMO3 prosthetic cofactor that maximizes residual enzyme activity), and acidic soaps (pH 5.5--6.5) to convert volatile volatile TMA into non-volatile water-soluble salts.

\medterm{Troponin}-- A heterotrimeric protein complex situated on the actin thin filament of striated muscle, composed of three subunits: troponin C (binds $Ca^{2+}$), troponin I (inhibits actomyosin ATPase), and troponin T (anchors the complex to tropomyosin). Cardiac-specific isoforms of troponin I (cTnI) and troponin T (cTnT) possess unique amino acid sequences distinct from skeletal muscle isoforms, establishing them as the definitive, gold-standard clinical biomarkers for diagnosing acute myocardial infarction (AMI). Following myocardial ischemic necrosis and sarcolemmal rupture, cytosolic and structurally bound cardiac troponins leak into the coronary and systemic circulation. High-sensitivity cardiac troponin (hs-cTn) assays detect elevations within 2 to 3 hours of symptom onset, peaking at 12 to 24 hours, and remaining detectable for 7 to 14 days, providing both acute diagnostic sensitivity and a wide detection window for subacute infarctions. The Fourth Universal Definition of Myocardial Infarction requires a rise and/or fall of cardiac troponin with at least one value above the 99th percentile upper reference limit, accompanied by clinical evidence of acute myocardial ischemia.

\medterm{Tumor Necrosis Factor}-- Tumor necrosis factor-alpha (TNF-$\alpha$), a potent homotrimeric pro-inflammatory cytokine ($~17$ kDa per monomer) produced primarily by activated macrophages, monocytes, T-lymphocytes, and natural killer cells in response to pathogen-associated molecular patterns (such as bacterial lipopolysaccharide, LPS). Initially expressed as a 26-kDa transmembrane precursor, it is proteolytically cleaved by TNF-alpha converting enzyme (TACE / ADAM17) to release soluble circulating homotrimeric TNF. TNF binds two cell-surface receptors: TNFR1 (ubiquitously expressed, containing an intracellular death domain that recruits TRADD to activate NF-$\kappa$B for inflammation and pro-survival, or FADD and caspase-8 to trigger apoptosis) and TNFR2 (predominantly immune cells, mediating tissue repair). Systemically, TNF-$\alpha$ is a master driver of acute and chronic inflammation: it activates vascular endothelial cells to express E-selectin, ICAM-1, and VCAM-1, promoting leukocyte extravasation; acts on the hypothalamus as an endogenous pyrogen to induce fever (via local PGE2 synthesis); and induces hepatic acute-phase reactant synthesis. In metabolic tissues, chronic low-grade elevation of TNF-$\alpha$ in obesity triggers serine phosphorylation of IRS-1 via IKK and JNK, directly inducing peripheral insulin resistance. In chronic illness, sustained TNF secretion suppresses appetite and stimulates lipoprotein lipase inhibition, driving profound muscle wasting and fat loss (cachexia / cachectin). Anti-TNF biologics (infliximab, adalimumab, etanercept) are cornerstone therapeutics for Crohn disease, rheumatoid arthritis, and psoriasis.

\medterm{Ubiquinone (Coenzyme Q10)}-- A lipid-soluble benzoquinone derivative featuring a side chain of ten repeating isoprenoid units (CoQ10), localized within the hydrophobic core of the inner mitochondrial membrane. Ubiquinone functions as an indispensable, mobile, two-electron and two-proton electron carrier connecting Complexes I and II to Complex III of the respiratory chain. It undergoes reversible reduction through three distinct oxidation states: fully oxidized ubiquinone (Q), a transient free-radical semiquinone intermediate ($QH^\bullet$), and fully reduced ubiquinol ($QH_2$). Ubiquinol collects electrons from NADH via NADH-ubiquinone oxidoreductase (Complex I), from $FADH_2$ via succinate dehydrogenase (Complex II), and from fatty acid beta-oxidation via electron-transferring flavoprotein dehydrogenase (ETFDH), shuttling them to Complex III via the Q-cycle to drive proton pumping into the intermembrane space. Beyond bioenergetics, ubiquinol acts as a potent, lipid-phase antioxidant protecting mitochondrial membranes from lipid peroxidation. CoQ10 is synthesized endogenously via the mevalonate/cholesterol biosynthetic pathway; consequently, HMG-CoA reductase inhibitors (statins) inhibit CoQ10 synthesis, which may contribute to statin-induced myopathy.

\medterm{Urea Cycle}-- A metabolic pathway occurring in both the mitochondria and cytoplasm of hepatocytes that converts toxic ammonia ($NH_3$), derived from amino acid catabolism, into nontoxic, water-soluble urea for excretion by the kidneys. The cycle consists of five enzymatic steps, starting with carbamoyl phosphate synthetase I (CPS I), which is the rate-limiting enzyme activated by N-acetylglutamate. Inherited deficiencies in any urea cycle enzyme (most commonly ornithine transcarbamylase deficiency, which is X-linked) lead to hyperammonemia, causing lethargy, seizures, cerebral edema, and cognitive impairment.

\medterm{Valinemia}-- A rare autosomal recessive metabolic disorder of branched-chain amino acid catabolism caused by selective deficiency of the enzyme valine aminotransferase (valine transaminase), which specifically catalyzes the reversible transamination of L-valine with $\alpha$-ketoglutarate to yield $\alpha$-ketoisovalerate and glutamate. Because the deficiency is restricted solely to the transaminase for valine, the transamination and decarboxylation of the other two branched-chain amino acids (leucine and isoleucine) proceed normally. Consequently, patients exhibit isolated, massive elevations of valine in plasma, cerebrospinal fluid, and urine, without elevations of leucine, isoleucine, or branched-chain $\alpha$-ketoacids (distinguishing it from maple syrup urine disease). Clinical manifestations in reported infants have included poor feeding, frequent vomiting, failure to thrive, generalized hypotonia, and psychomotor developmental delay, although some cases exhibit a benign or asymptomatic course. Diagnosis is confirmed by quantitative plasma amino acid chromatography showing isolated hypervalinemia. Treatment consists of dietary valine restriction using specialized amino acid formulas, titrating valine intake to maintain plasma concentrations within the normal physiological range while supporting normal somatic growth.

\medterm{Variegate Porphyria}-- An autosomal dominant acute hepatic porphyria caused by heterozygous mutations in the PPOX gene on chromosome 1, resulting in an approximate 50\% reduction in the activity of protoporphyrinogen oxidase. Protoporphyrinogen oxidase is the penultimate enzyme of the heme biosynthetic pathway, localized to the inner mitochondrial membrane, where it catalyzes the six-electron oxidation of protoporphyrinogen IX to protoporphyrin IX. Variegate porphyria is unique among the porphyrias in that it produces a hybrid clinical phenotype, combining both acute neurovisceral attacks and cutaneous photosensitivity: (1) Acute Neurovisceral Crises: Precipitated by cytochrome P450-inducing drugs (barbiturates, sulfonamides, antiepileptics), fasting, infection, or hormonal fluctuations, which deplete free regulatory heme and induce hepatic $\delta$-aminolevulinic acid synthase-1 (ALAS1). This drives massive accumulation of neurotoxic porphyrin precursors ($\delta$-aminolevulinic acid [ALA] and porphobilinogen [PBG]), causing severe agonizing abdominal pain, peripheral motor neuropathy (flaccid paralysis, respiratory failure), autonomic instability (tachycardia, hypertension), and neuropsychiatric disturbances (psychosis, seizures); (2) Cutaneous Lesions: Accumulated porphyrinogens diffuse into the skin, where solar UV light excites them to generate singlet oxygen, causing severe mechanical fragility, bullae, erosions, and hyperpigmentation on sun-exposed areas. A unique fluorescence emission peak at 626 nm on plasma fluorometry at neutral pH is pathognomonic, distinguishing it from acute intermittent porphyria and porphyria cutanea tarda. Acute attacks are treated with intravenous hemin (which feedback-inhibits ALAS1) and glucose infusions.

\medterm{Vitamin A}-- A group of fat-soluble retinoids (retinol, retinal, retinoic acid, and retinyl esters) and provitamin A carotenoids ($\beta$-carotene) essential for vision, cellular differentiation, immune competence, and embryonic morphogenesis. In the retina, 11-cis-retinal combines with the apoprotein opsin to form rhodopsin in rod photoreceptor cells; photon absorption isomerizes 11-cis-retinal to all-trans-retinal, activating transducin and generating the visual nerve impulse. In non-retinal tissues, all-trans-retinoic acid functions as a nuclear hormone, binding retinoic acid receptors (RAR and RXR) to regulate the transcription of genes controlling epithelial cellular differentiation and mucus secretion. Vitamin A deficiency is a leading cause of preventable childhood blindness worldwide, manifesting sequentially as: nyctalopia (night blindness), xerophthalmia (pathological corneal dryness), Bitot spots (foamy keratin plaques on the bulbar conjunctiva), keratomalacia (corneal softening and perforation), and squamous metaplasia of respiratory and urinary tract epithelia, causing recurrent infections. Conversely, acute hypervitaminosis A causes nausea, vomiting, increased intracranial pressure (pseudotumor cerebri), and papilledema; chronic toxicity causes hepatosplenomegaly, bone pain, hypercalcemia, and alopecia. Retinoids (such as isotretinoin) are extremely teratogenic, causing microtia, craniofacial dysmorphism, and conotruncal heart defects if ingested during early pregnancy.

\medterm{Vitamin B12 (Cobalamin)}-- A complex, cobalt-containing water-soluble vitamin with a corrin ring structure, synthesized exclusively by microorganisms and obtained in humans strictly from animal dietary sources (meat, dairy, fish, eggs). Gastric acid and pepsin liberate cobalamin from dietary proteins; it binds salivary haptocorrin (R-factor), which is digested by pancreatic proteases in the duodenum, allowing cobalamin to bind intrinsic factor (IF, secreted by gastric parietal cells). The cobalamin-IF complex resists digestion and undergoes receptor-mediated endocytosis via cubam receptors in the terminal ileum. Large hepatic stores (2--5 mg) protect against deficiency for 3 to 5 years. Cobalamin serves as an essential coenzyme for two vital enzymes: (1) Methionine Synthase (cytoplasmic, methylcobalamin-dependent): transfers a methyl group from $N^5$-methyl-THF to homocysteine, generating methionine and regenerating free THF; failure traps folate as $N^5$-methyl-THF ('methyl trap'), preventing nucleotide synthesis and causing megaloblastic anemia with hypersegmented neutrophils; and (2) Methylmalonyl-CoA Mutase (mitochondrial, adenosylcobalamin-dependent): isomerizes methylmalonyl-CoA to succinyl-CoA; failure causes accumulation of toxic methylmalonic acid (MMA). Cobalamin deficiency causes severe macrocytic megaloblastic anemia and Subacute Combined Degeneration (SCD) of the spinal cord (demyelination of the dorsal columns [loss of vibration/proprioception], lateral corticospinal tracts [spasticity, paraplegia], and spinocerebellar tracts). Pernicious anemia (autoimmune destruction of gastric parietal cells or IF) and terminal ileum resection are classic causes. Diagnosis is confirmed by elevated serum methylmalonic acid and homocysteine.

\medterm{Vitamin C (Ascorbic Acid)}-- A water-soluble hexose-derived lactone that functions as an essential electron donor, biological reducing agent, and antioxidant. Because humans lack L-gulonolactone oxidase (the terminal enzyme in the ascorbate biosynthetic pathway), vitamin C must be obtained from dietary fruits and vegetables. Ascorbic acid functions as a specific cofactor for copper- and iron-containing mixed-function monooxygenases and dioxygenases, maintaining active-site metal ions in their reduced states ($Fe^{2+}$ and $Cu^+$): (1) Collagen Biosynthesis: acts as an obligatory cofactor for prolyl 4-hydroxylase and lysyl hydroxylase, which hydroxylate proline and lysine residues in procollagen chains inside the rough endoplasmic reticulum, allowing stable hydrogen-bonded triple-helix formation and covalent cross-linking; (2) Catecholamine Biosynthesis: cofactor for dopamine $\beta$-hydroxylase, which converts dopamine to norepinephrine; (3) Carnitine Synthesis: required for two dioxygenase steps in carnitine biosynthesis; and (4) Intestinal Iron Absorption: reduces dietary non-heme ferric iron ($Fe^{3+}$) to more soluble ferrous iron ($Fe^{2+}$) in the stomach. Vitamin C deficiency causes Scurvy, characterized by defective collagen synthesis, capillary fragility, and impaired wound healing. Clinical signs include painful corkscrew hairs with perifollicular petechiae and purpura, spongy bleeding gums, subperiosteal hematomas, hemarthroses, and impaired bone formation in children.

\medterm{Vitamin D}-- A fat-soluble secosteroid hormone precursor obtained through diet (ergocalciferol/D2 from fungi; cholecalciferol/D3 from animal sources) or synthesized endogenously in the epidermis when ultraviolet B (UVB) light converts 7-dehydrocholesterol to pre-vitamin D3, which thermally isomerizes into cholecalciferol. Cholecalciferol is transported to the liver and hydroxylated at carbon-25 by cytochrome P450 25-hydroxylase (CYP2R1) to form 25-hydroxyvitamin D [25(OH)D, calcifediol], the primary circulating storage form used clinically to assess vitamin D status (normal: 30--100 ng/mL). In the proximal renal tubule, 25(OH)D is hydroxylated at carbon-1 by $1\alpha$-hydroxylase (CYP27B1)---an enzyme strongly stimulated by parathyroid hormone (PTH) and hypophosphatemia, and inhibited by FGF23 and hypercalcemia---to form the biologically active hormone 1,25-dihydroxyvitamin D [1,25(OH)$_2$D, calcitriol]. Calcitriol binds the nuclear vitamin D receptor (VDR), upregulating calcium-transporting proteins (TRPV6, calbindin-D28k, PMCA1b) in enterocytes to massively enhance intestinal calcium and phosphate absorption, and stimulates osteoblasts to express RANKL, supporting bone remodeling. Deficiency causes Rickets in children (characterized by failure of epiphyseal plate mineralization, craniotabes, rachitic rosary, and bowing deformities of long bones) and Osteomalacia in adults (impaired mineralization of osteoid, diffuse bone pain, and pseudofractures).

\medterm{Vitamin E}-- A group of eight fat-soluble lipid compounds comprising four tocopherols ($\alpha$, $\beta$, $\gamma$, $\delta$) and four tocotrienols, with $\alpha$-tocopherol being the most biologically active form in humans due to preferential hepatic retention by $\alpha$-tocopherol transfer protein ($\alpha$-TTP). Vitamin E functions as the premier lipid-soluble, chain-breaking antioxidant in human biology, intercalating directly into cell membranes, myelin sheaths, and circulating lipoproteins. It specifically intercepts lipid peroxyl radicals ($LOO^\bullet$) generated during free radical attack on polyunsaturated fatty acids (PUFAs), donating a phenolic hydrogen atom to terminate the autocatalytic chain reaction of lipid peroxidation. The resulting oxidized tocopheroxyl radical is relatively unreactive and is reduced back to active $\alpha$-tocopherol by ascorbic acid (vitamin C) and glutathione. Because vitamin E is abundant in vegetable oils, nuts, and seeds, primary deficiency is exceptionally rare; secondary deficiency occurs in severe fat malabsorption syndromes (cystic fibrosis, abetalipoproteinemia, chronic cholestasis, celiac disease). Clinically, vitamin E deficiency presents with symptoms that closely mimic Friedreich ataxia and vitamin B12 deficiency: posterior column degeneration (loss of proprioception and vibratory sensation), spinocerebellar tract ataxia, areflexia, muscle weakness, and hemolytic anemia (due to increased erythrocyte membrane fragility from unrestrained oxidative peroxidation).

\medterm{Vitamin K}-- A group of fat-soluble naphthoquinone compounds comprising phylloquinone (vitamin K1, synthesized by green plants) and menaquinones (vitamin K2, synthesized by intestinal bacterial microflora). Vitamin K functions as an obligatory coenzyme for the endoplasmic reticulum enzyme $\gamma$-glutamyl carboxylase. This enzyme converts specific peptide-bound glutamate (Glu) residues into $\gamma$-carboxyglutamate (Gla) residues on several crucial plasma proteins: coagulation factors II (prothrombin), VII, IX, and X, natural anticoagulant proteins C and S, and the bone matrix protein osteocalcin. The introduced dicarboxylic Gla domains possess high affinity for divalent calcium ($Ca^{2+}$), enabling these clotting factors to undergo conformational shifts and bind negatively charged phospholipid membranes exposed on activated platelets via calcium bridges, assembly without which the tenase and prothrombinase complexes cannot function. During carboxylation, active vitamin K hydroquinone is oxidized to vitamin K epoxide; the enzyme vitamin K epoxide reductase (VKORC1) reduces it back to hydroquinone. Warfarin and coumarin anticoagulants selectively inhibit VKORC1, preventing vitamin K recycling. Neonates have sterile intestines, poor placental transfer, and low breast-milk vitamin K, putting them at extreme risk for Vitamin K Deficiency Bleeding (VKDB / hemorrhagic disease of the newborn, presenting with intracranial hemorrhage), universally prevented by intramuscular phytonadione at birth. In adults, deficiency develops from prolonged broad-spectrum antibiotics, severe malabsorption, or biliary obstruction, prolonging the prothrombin time (PT/INR) first (due to factor VII's short 4--6 hour half-life) followed by aPTT.

\medterm{Vitamins}-- Organic micronutrients required in small amounts for metabolic function, cannot be synthesized in sufficient quantities. Fat-soluble (A, D, E, K) absorbed with lipids, stored in liver/fat. Water-soluble (B complex, C) not stored, excess excreted. Deficiency causes specific syndromes: A (night blindness), B1 (beriberi), B3 (pellagra), B12 (megaloblastic anemia), C (scurvy), D (rickets), K (bleeding).

\medterm{VLDL}-- Very low-density lipoprotein, an endogenous triglyceride-rich lipoprotein (density 0.95--1.006 g/mL, diameter 30--80 nm) synthesized and secreted exclusively by hepatocytes to transport endogenously formed triglycerides from the liver to peripheral tissues (adipose tissue and skeletal muscle). Composed of an apolar core dominated by triglycerides (55--65\% by mass) and cholesteryl esters, surrounded by a monolayer containing free cholesterol, phospholipids, and a single molecule of full-length apolipoprotein B-100 (ApoB-100). Hepatic VLDL assembly requires the microsomal triglyceride transfer protein (MTP); inactivating mutations in MTP cause abetalipoproteinemia. Upon entering the circulation, VLDL acquires ApoC-II and ApoE from HDL particles. At capillary endothelia, lipoprotein lipase (LPL), activated by ApoC-II, hydrolyzes core triglycerides into free fatty acids for cellular energy or re-esterification. Progressive triglyceride depletion shrinks the VLDL particle, converting it into intermediate-density lipoprotein (IDL / VLDL remnant). Approximately half of IDL particles are cleared by hepatic LDL receptors via ApoE-mediated uptake, while the remainder undergo further triglyceride hydrolysis by hepatic lipase (HL) and shed ApoE to become cholesterol-enriched LDL. VLDL secretion is heavily stimulated by obesity, insulin resistance, alcohol consumption, and high-carbohydrate diets.

\medterm{Von Willebrand Disease}-- The most common inherited bleeding disorder in humans (prevalence $\approx 1\%$), caused by quantitative deficiencies or qualitative defects in von Willebrand factor (vWF), a massive multimeric plasma glycoprotein synthesized by vascular endothelial cells (stored in Weibel-Palade bodies) and megakaryocytes ($\alpha$-granules). vWF performs two essential hemostatic functions: (1) It acts as a molecular bridge between exposed subendothelial collagen (via its A3 domain) and platelet glycoprotein Ib-IX-V receptors (via its A1 domain) under high shear stress in injured arterioles, mediating initial platelet adhesion; and (2) It binds and stabilizes circulating coagulation factor VIII, extending its half-life from 2 hours to 12 hours. The disorder is classified into three major types: Type 1 (autosomal dominant, 70--80\% of cases, mild-to-moderate quantitative reduction of structurally normal vWF); Type 2 (autosomal dominant or recessive, qualitative dysfunction of vWF multimers); and Type 3 (autosomal recessive, near-total absence of vWF and extremely low factor VIII, mimicking severe hemophilia). Clinical manifestations include mucocutaneous bleeding (epistaxis, menorrhagia, gingival hemorrhage, easy bruising). Diagnostic evaluation reveals normal or low platelet count, prolonged bleeding time, normal PT, variable prolongation of aPTT, reduced vWF antigen, and decreased ristocetin cofactor activity. Treatment includes desmopressin (DDAVP), which triggers vWF release from endothelial stores, and vWF/FVIII replacement concentrates.

\medterm{Wilson Disease}-- An autosomal recessive disorder of copper metabolism caused by loss-of-function mutations in the ATP7B gene on chromosome 13, which encodes a copper-transporting P-type ATPase localized to the trans-Golgi network of hepatocytes. ATP7B performs two vital functions: (1) It incorporates copper into apoceruloplasmin to synthesize functional ceruloplasmin; and (2) It facilitates the excretion of excess copper into bile, the body's sole major excretory route for copper. Loss of ATP7B causes systemic copper overload: copper initially accumulates in hepatocytes, generating toxic hydroxyl free radicals via the Fenton reaction, leading to lipid peroxidation, mitochondrial damage, and hepatocellular necrosis. Spilled into the systemic circulation, unbound free copper deposits in extrahepatic tissues, especially the basal ganglia (putamen and globus pallidus), cerebral cortex, kidneys, and the Descemet membrane of the cornea. Clinical manifestations include hepatic disease (acute hepatitis, chronic cirrhosis, fulminant hepatic failure) and neuropsychiatric manifestations (parkinsonism, resting and intention tremors, dysarthria, dystonia, choreoathetosis, depression, and psychosis). Ophthalmic slit-lamp examination reveals pathognomonic Kayser-Fleischer rings (golden-brown copper rings at the corneal limbus). Diagnostic laboratory hallmarks include low serum ceruloplasmin ($<20$ mg/dL), markedly elevated 24-hour urinary copper excretion ($>100$ mcg/day), and elevated hepatic parenchymal copper concentration on liver biopsy. Treatment includes copper chelating agents (D-penicillamine, trientine) and oral zinc salts (which induce intestinal enterocyte metallothionein to block dietary copper absorption).

\medterm{Zinc}-- An essential catalytic, structural, and regulatory trace metal micronutrient ($Zn^{2+}$) required for the physiological function of over 300 enzymes and more than 2,000 transcription factors in the human body. Unlike iron, zinc does not undergo redox cycling, functioning instead as a stable Lewis acid: (1) Catalytic cofactor: required by carbonic anhydrase (acid-base balance and $CO_2$ transport), alcohol dehydrogenase, alkaline phosphatase, carboxypeptidase A, RNA and DNA polymerases, and cytosolic copper-zinc superoxide dismutase (Cu/Zn-SOD); (2) Structural role: stabilizes 'zinc finger' motifs, loop structures coordinated by cysteine and histidine residues that fit into the major groove of double-stranded DNA to regulate sequence-specific gene transcription; (3) Immune competence: supports thymulin production, T-helper cell differentiation, and natural killer cell activity; and (4) Wound healing: required for matrix metalloproteinase (MMP) activity and collagen synthesis. Ingested zinc is absorbed in the duodenum and jejunum via ZIP4 transporters. Acrodermatitis enteropathica is a rare autosomal recessive disorder caused by SLC39A4 (ZIP4) mutations, presenting upon weaning from breast milk with the classic triad of periorificial and acral vesiculobullous dermatitis, intractable diarrhea, and diffuse alopecia. Acquired zinc deficiency causes impaired wound healing, dysgeusia (loss of taste), anosmia, immune dysfunction, hypogonadism, and delayed sexual maturation.
